\chapter{Spectral de Rham Reflection and Infinitesimal Descent}
\label{chap:spectral-de-Rham}

\section{Orientation: from infinitesimal invariance to de Rham descent}
\label{sec:de-Rham-purpose}

The preceding chapters supplied two complementary pieces of structure. The
differentiated spectral doctrine constructed structured square-zero extensions,
finite infinitesimal substitutions, and their endpoint semantics; formal
reconstruction then showed how suitably complete nonlinear geometry can be
recovered from infinitesimal data. The present chapter asks a different
question. We identify contexts that differ only by finite infinitesimal
variation and ask what descent theory survives after that identification.

The principal result is a de Rham descent calculus attached to that
identification. We construct the universal left-exact reflection that inverts
finite infinitesimal variation, separate its universal target from the effective
quotient actually presented by points of the original context, and show that
stable scalar descent on the resulting effective relation produces intrinsic
forms, a homotopy-coherent differential, products, closedness data, and
primitives. Two later comparisons then measure different boundaries of this
calculus: first-order restriction recovers the cotangent complex, whereas
reduced-classical reflection adds the convergence needed to pass beyond finite
infinitesimal stages.

The first step is therefore a localization problem. Affine square-zero
substitutions generate a pullback-stable congruence, and the universal
left-exact localization at that congruence is the spectral de Rham reflection.
Its local objects are precisely the stacks that cannot distinguish a finite
infinitesimal thickening from its centre. The following display records the
eventual reflector; its local-equivalence class and relative form are constructed
before the notation is used systematically:
\[
  L^{\mathrm{dR}}:
  \mathcal H_{\Sp}
  \rightleftarrows
  \mathcal H_{\Sp}^{\mathrm{dR}}:
  \iota^{\mathrm{dR}}.
\]

A second issue appears as soon as we try to extract a descent object from the
relative unit. The effective image of a localization unit can be strictly
smaller than the whole reflected object. Its kernel-pair nerve therefore
presents that effective image. We separate the universal and effective parts:
the relative unit factors through an effective quotient \(Q_{X/S}\), and the
resulting kernel-pair groupoid presents \(Q_{X/S}\). This distinction will remain active
throughout the chapter. Stable scalar functions on the groupoid recover
functions on \(Q_{X/S}\); functions on the full reflected object are compared
to them by restriction.

The dependency structure therefore branches into a universal object and an
effective presentation; the effective branch supplies the descent calculus,
while the universal branch remains available for comparison. The next diagram is only a route map: its
symbols label objects introduced later, and no notation in it is used
systematically before its construction.
\[
\begin{tikzcd}[column sep=large,row sep=large]
  & \text{finite infinitesimal substitutions}
      \arrow[d,"\text{left-exact localization}"]
  &
  \\
  & \text{relative de Rham reflection}
      \arrow[dl,"\text{effective image}"']
      \arrow[dr,"\text{full reflected object}"]
  &
  \\
  Q_{X/S}
      \arrow[d,"\text{\v{C}ech relation}"']
  &&
  X_{\mathrm{dR}/S}
      \arrow[d,"\text{stable scalar functions}"]
  \\
  \mathcal R^{\mathrm{dR}}_{X/S,\bullet}
      \arrow[d,"\text{stable scalar functions}"']
  &&
  \mathrm{DR}^{\mathrm{full}}_{X/S}
      \arrow[dl,"\text{restriction}"]
  \\
  C^\bullet_{\mathrm{dR}}(X/S)
      \arrow[r,"\operatorname*{Tot}"']
  &
  \mathrm{DR}_{X/S}
      \arrow[d,"\text{totalization filtration}"]
  &
  \\
  & \bigl(\Omega^\bullet_{\mathrm{dR}}(X/S),d_{\mathrm{dR}}\bigr)
  &
\end{tikzcd}
\]

The chapter follows this dependency order. We first construct the absolute
and indexed de Rham reflectors and identify finite-infinitesimal invariance.
We then factor the relative unit and prove that its effective relation is the
\v{C}ech groupoid that supports descent. Only then do we build the stable
scalar interface needed to take functions on that relation. Effective descent
produces the cochain algebra and its totalization; stable Dold--Kan then exposes
forms, the homotopy-coherent differential, products, closedness, and primitives.
The final sections establish descent and comprehension and treat the cotangent
and reduced-classical comparisons as two separate exits from the construction.

The conventions on universes, sheaves, effective epimorphisms, higher
coherence, and notation from Chapter~\Ref{chap:spectral-algebraic-geometry}
remain in force. In particular, once the functorial source of a coherence has
been identified, routine identity, composition, mate, and higher-pasting
coherences are governed by the standing convention and will not be restated at
every later occurrence.
\section{From infinitesimal substitution to indexed de Rham reflection}
\label{sec:de-Rham-reflective-doctrine}

The differentiated doctrine already supplies the arrows that should become
invisible: underlying square-zero infinitesimal substitutions. To turn that
test geometry into a modality, we must prove that the class generated by those
arrows is stable under arbitrary pullback. That pullback stability is the
substantive input: once it is established, accessible localization produces a
left-exact reflector, and left exactness then makes the construction compatible
with substitution.

This section therefore has three tasks. We first pass from affine square-zero
generators to a pullback-stable strongly saturated congruence. We then
construct the universal reflector and identify its local objects by
finite-infinitesimal invariance, including the endpoint comparison of the
preceding chapter. Finally we relativize the construction over a base and
prove coherent base change, thereby obtaining an indexed reflection rather
than merely an ambient localization.
\subsection{Affine square-zero generators and the pullback-stable congruence}
\label{sec:de-Rham-congruence}

The desired reflector will be left exact only if its local-equivalence class is
stable under pullback. We therefore begin one level below the reflector
itself. Smallness gives an accessible localization at the affine square-zero
arrows, while affine density lets us test pullback stability on affine base
changes. The point of the next lemma is to show that no additional
pullback-generators are required.

By Convention~\Ref{conv:size-spectral-ag}, the category of
\(\mathbb U\)-small affine spectral schemes is \(\mathbb V\)-small. Let
\(\Sigma_{\mathrm{sqz}}\) be the resulting \(\mathbb V\)-small collection of
morphisms of representables induced by the underlying square-zero infinitesimal substitutions
of \Ref{def:spectral-square-zero-substitution} between \(\mathbb U\)-small affine
spectral schemes. Thus a generator has the form
\[
  h_A\longrightarrow h_{A'}
\]
for a spectral square-zero extension \(A'\to A\).

For a map \(h_V\to h_{A'}\), pullback is represented by the affine base-change
square
\[
\begin{tikzcd}[column sep=large,row sep=large]
  h_V\times_{h_{A'}}h_A
    \arrow[r]
    \arrow[d]
  &
  h_A
    \arrow[d]
  \\
  h_V
    \arrow[r]
  &
  h_{A'}.
\end{tikzcd}
\]

\begin{lemma}[Affine density in slices]
\label{lem:de-Rham-affine-density}
Let \(S\in\mathcal H_{\Sp}\) and \(Y\to S\) be an object of the slice.
Then the canonical morphism
\[
  \operatorname*{colim}_{(T,u)\in\operatorname{AffPt}(Y)}h_T
  \longrightarrow Y
\]
is an equivalence in \((\mathcal H_{\Sp})_{/S}\).
\end{lemma}

\begin{proof}
\leavevmode
\begin{enumerate}[label=\textup{(\arabic*)}]
    \item \textbf{Establish the next claim.}
    Regard \(Y\) through the fully faithful inclusion
    \(\mathcal H_{\Sp}\hookrightarrow\PSh(\Aff_{\Sp})\).
    \item \textbf{Apply Yoneda density.}
    Yoneda density gives, in the presheaf slice over the
    underlying presheaf of \(S\),
    \[
    Y
    \simeq
    \operatorname*{colim}_{(T,u)\in\operatorname{AffPt}(Y)}h_T.
    \]
    \item \textbf{Record smallness of the indexing category.}
    The indexing category is \(\mathbb V\)-small by
    Convention~\Ref{conv:size-spectral-ag}.
    \item \textbf{Use colimit preservation by sheafification.}
    Spectral Zariski sheafification
    preserves colimits, every representable is already a spectral Zariski sheaf,
    and \(Y\) is already local.
    \item \textbf{Use the stated construction.}
    Applying sheafification to the density colimit
    therefore gives the displayed equivalence in
    \((\mathcal H_{\Sp})_{/S}\).
    \item \textbf{Identify the affine-point indexing category.}
    The affine-point category is the canonical
    indexing object for this density presentation.
    Compare
    the Yoneda-density and sheaf-localization formalism of
    \cite{LurieHigherToposTheory}.
\end{enumerate}
\end{proof}

\begin{definition}[De Rham local-equivalence class]
Let \(W^{\mathrm{dR}}\) be the smallest strongly saturated class of morphisms in
\(\mathcal H_{\Sp}\) which contains \(\Sigma_{\mathrm{sqz}}\) and is
stable under arbitrary pullback.
\end{definition}

\begin{lemma}[Generated square-zero congruence]
\label{lem:generated-congruence}
The strongly saturated class \(\operatorname{Sat}(\Sigma_{\mathrm{sqz}})\) is
already stable under arbitrary pullback. Hence
\[
  W^{\mathrm{dR}}
  =
  \operatorname{Sat}(\Sigma_{\mathrm{sqz}}).
\]
If \(f:A\to B\) lies in \(W^{\mathrm{dR}}\), every iterated diagonal of \(f\)
also lies in \(W^{\mathrm{dR}}\).
\end{lemma}

\begin{proof}
\leavevmode
\begin{enumerate}[label=\textup{(\arabic*)}]
    \item \textbf{Base-change every square-zero generator.}
    Put
    \[
    W_0:=\operatorname{Sat}(\Sigma_{\mathrm{sqz}}),
    \]
    and let \(L_\Sigma\) be the accessible localization generated by the small
    family \(\Sigma_{\mathrm{sqz}}\), whose existence follows from
    \cite[Proposition~5.5.4.15]{LurieHigherToposTheory}. Its
    local-equivalence class is \(W_0\).

    Let \(h_U\to h_T\) be a generator and let \(Y\to h_T\) be arbitrary.
    Lemma~\ref{lem:de-Rham-affine-density} gives
    \[
    Y\simeq
    \operatorname*{colim}_{(V,u)\in\operatorname{AffPt}(Y)}h_V
    \]
    in the slice over \(h_T\). Colimits in an infinity-topos are universal, so
    \[
    Y\times_{h_T}h_U
    \simeq
    \operatorname*{colim}_{(V,u)}
    \bigl(h_V\times_{h_T}h_U\bigr),
    \]
    and the pullback arrow
    \[
    Y\times_{h_T}h_U\longrightarrow Y
    \]
    is the corresponding colimit in the arrow category. Every stage
    \[
    h_V\times_{h_T}h_U\longrightarrow h_V
    \]
    is represented by the affine base change \(V\times_TU\to V\), hence is again
    a square-zero infinitesimal substitution by
    \Ref{prop:square-zero-substitution-permanence}. Since \(L_\Sigma\) preserves
    colimits and inverts every stage map, it inverts their colimit. Thus every
    pullback of every generator belongs to \(W_0\).
    \item \textbf{Fix the indicated data.}
    Let
    \[
    W_0^{\mathrm{univ}}
    :=
    \left\{
    f:X\to Y
    \ \middle|\
    \text{every pullback of \(f\) belongs to \(W_0\)}
    \right\}.
    \]
    \item \textbf{Establish the indicated equivalence.}
    This class is stable under arbitrary pullback by definition and contains all
    equivalences.
    \item \textbf{Establish the next claim.}
    It is closed under retracts and small colimits in the arrow
    category: after any base change, retract diagrams remain retract diagrams and,
    by universality of colimits in an infinity-topos, arrow-category colimits
    become the corresponding colimits of the base-changed arrows.
    \item \textbf{Establish the next claim.}
    Strong
    saturation of \(W_0\) then applies.
    \item \textbf{Establish the next claim.}
    We use the arrow-category
    characterization of strong saturation from
    \cite[Section~5.5.4]{LurieHigherToposTheory}.
    \item \textbf{Establish the next claim.}
    It remains to verify three-for-two.
    \item \textbf{Fix the indicated data.}
    Let
    \[
    X\xrightarrow{f}Y\xrightarrow{g}Z
    \]
    be composable.
    \item \textbf{Treat the stated case.}
    If \(f,g\in W_0^{\mathrm{univ}}\), then every pullback of
    \(gf\) is a composite of pullbacks of \(f\) and \(g\), so
    \(gf\in W_0^{\mathrm{univ}}\).
    \item \textbf{Treat the stated case.}
    If \(f,gf\in W_0^{\mathrm{univ}}\), let \(Z'\to Z\) be arbitrary and put
    \[
    Y':=Y\times_ZZ',
    \qquad
    X':=X\times_ZZ'.
    \]
    \item \textbf{Use pullback stability.}
    Then \(X'\to Z'\), being a pullback of \(gf\), lies in \(W_0\), while
    \(X'\to Y'\), being a pullback of \(f\), lies in \(W_0\).
    \item \textbf{Apply two-out-of-three.}
    Two-out-of-three
    in \(W_0\) gives \(Y'\to Z'\in W_0\).
    Hence
    \(g\in W_0^{\mathrm{univ}}\).
    \item \textbf{Establish the next claim.}
    For the remaining cancellation case assume
    \(g,gf\in W_0^{\mathrm{univ}}\), and let \(b:Y'\to Y\) be arbitrary.
    \item \textbf{Introduce the notation.}
    Put
    \[
    \overline Y:=Y\times_ZY',
    \]
    where \(Y'\to Z\) is \(gb\).
    \item \textbf{Use pullback stability.}
    The projection
    \[
    p:\overline Y\longrightarrow Y'
    \]
    is a pullback of \(g\), hence belongs to
    \(W_0^{\mathrm{univ}}\).
    \item \textbf{Construct the indicated map.}
    The map \(b\) determines a canonical section
    \[
    s:Y'\longrightarrow\overline Y,
    \qquad
    ps\simeq\id_{Y'}.
    \]
    \item \textbf{Use pullback stability.}
    Every pullback of \(s\) remains a section of the corresponding pullback of
    \(p\).
    \item \textbf{Use the stated hypothesis.}
    Since every pullback of \(p\) lies in \(W_0\), two-out-of-three with
    the identity shows that every pullback of \(s\) lies in \(W_0\).
    Thus
    \(s\in W_0^{\mathrm{univ}}\).
    \item \textbf{Establish the next claim.}
    Now put
    \[
    X_Z':=X\times_ZY'.
    \]
    \item \textbf{Use pullback stability.}
    The map \(X_Z'\to Y'\) is a pullback of \(gf\), hence lies in \(W_0\).
    \item \textbf{Use pullback stability.}
    Moreover the canonical map
    \[
    X\times_YY'\longrightarrow X_Z'
    \]
    is the pullback of \(s\) along the morphism
    \(X_Z'\to\overline Y\) induced by \(f\).
    \item \textbf{Use the stated hypothesis.}
    Since
    \(s\in W_0^{\mathrm{univ}}\), this canonical map lies in \(W_0\).
    \item \textbf{Apply two-out-of-three.}
    Their
    composite is the base change of \(f\) along \(b\), so two-out-of-three gives
    \[
    X\times_YY'\longrightarrow Y'\in W_0.
    \]
    \item \textbf{Establish the next claim.}
    As \(b\) was arbitrary, \(f\in W_0^{\mathrm{univ}}\).
    \item \textbf{Conclude the stated claim.}
    Hence
    \(W_0^{\mathrm{univ}}\) is strongly saturated.
    \item \textbf{Establish the next claim.}
    Step \textup{(1)} shows
    \(\Sigma_{\mathrm{sqz}}\subseteq W_0^{\mathrm{univ}}\).
    Therefore
    \[
    W_0=\operatorname{Sat}(\Sigma_{\mathrm{sqz}})
    \subseteq W_0^{\mathrm{univ}}.
    \]
    \item \textbf{Prove the converse implication.}
    Conversely, the identity base change of any
    \(f\in W_0^{\mathrm{univ}}\) shows \(f\in W_0\).
    \item \textbf{Conclude the stated claim.}
    Thus
    \[
    W_0=W_0^{\mathrm{univ}},
    \]
    proving that \(W_0\) is stable under arbitrary pullback.
    \item \textbf{Use the stated construction.}
    The defining
    minimality of \(W^{\mathrm{dR}}\) now gives
    \[
    W^{\mathrm{dR}}
    =
    W_0.
    \]
    \item \textbf{Treat the stated case.}
    If \(f:A\to B\) lies in \(W^{\mathrm{dR}}\), the projection
    \[
    A\times_BA\longrightarrow A
    \]
    is a pullback of \(f\), hence belongs to \(W^{\mathrm{dR}}\).
    \item \textbf{Establish the next claim.}
    Its diagonal
    section
    \[
    A\longrightarrow A\times_BA
    \]
    composes with the projection to the identity.
    \item \textbf{Apply two-out-of-three.}
    Two-out-of-three therefore puts
    the diagonal in \(W^{\mathrm{dR}}\).
    \item \textbf{Establish the next claim.}
    Applying the same argument recursively
    proves the assertion for every iterated diagonal.
\end{enumerate}
\end{proof}

\subsection{The universal spectral de Rham reflector}
\label{sec:intrinsic-de-Rham-reflector}

The generated congruence is now known to be pullback-stable. This is exactly
the condition needed to promote accessible localization to a left-exact
localization of the spectral Zariski topos. After constructing the reflector,
we extend its affine generating property to global square-zero and finite
infinitesimal substitutions and record the resulting mapping property.

\begin{theorem}[Intrinsic spectral de Rham localization]
\label{thm:intrinsic-de-Rham-localization}
There is an accessible left-exact localization
\[
  L^{\mathrm{dR}}:
  \mathcal H_{\Sp}
  \rightleftarrows
  \mathcal H_{\Sp}^{\mathrm{dR}}
  :\iota^{\mathrm{dR}}
\]
displayed by the adjunction
\[
\begin{tikzcd}[column sep=huge,row sep=large]
  \mathcal H_{\Sp}
    \arrow[r,shift left=1.2ex,"L^{\mathrm{dR}}"]
  &
  \mathcal H_{\Sp}^{\mathrm{dR}}
    \arrow[l,shift left=1.2ex,"\iota^{\mathrm{dR}}"']
\end{tikzcd}
\]
with localization unit \(\eta_X^{\mathrm{dR}}:X\to X_{\mathrm{dR}}\), whose local-equivalence
class is \(W^{\mathrm{dR}}\). It is initial among accessible left-exact localizations
of \(\mathcal H_{\Sp}\) which invert
every affine square-zero infinitesimal substitution.
\end{theorem}

\begin{proof}
\leavevmode
\begin{enumerate}[label=\textup{(\arabic*)}]
    \item \textbf{Establish the locality claim.}
    The spectral Zariski topos is presentable, so localization at the
    \(\mathbb V\)-small family \(\Sigma_{\mathrm{sqz}}\) exists by the presentable
    localization theorem;
    see \cite[Proposition~5.5.4.15]{LurieHigherToposTheory}.
    \item \textbf{Establish the indicated equivalence.}
    Its local-equivalence
    class is the strongly saturated closure \(W^{\mathrm{dR}}\), which is
    pullback-stable by Lemma~\ref{lem:generated-congruence}.
    \item \textbf{Establish the locality claim.}
    The left-exactness
    criterion for a localization of an infinity-topos therefore makes the localization
    left exact; see \cite[Proposition~6.2.1.1]{LurieHigherToposTheory}.
    \item \textbf{Establish the locality claim.}
    Any accessible left-exact localization which inverts every affine square-zero
    infinitesimal substitution inverts every member of
    \(\Sigma_{\mathrm{sqz}}\).
    \item \textbf{Establish the indicated equivalence.}
    Its local-equivalence class is strongly saturated and
    therefore contains
    \[
    \operatorname{Sat}(\Sigma_{\mathrm{sqz}})
    =
    W^{\mathrm{dR}}.
    \]
    \item \textbf{Establish the next claim.}
    The universal factorization follows.
\end{enumerate}
\end{proof}

\begin{definition}[Spectral de Rham reflection]
Define the ambient idempotent left-exact endofunctor
\[
  \operatorname{dR}
  :=
  \iota^{\mathrm{dR}}L^{\mathrm{dR}}:
  \mathcal H_{\Sp}\longrightarrow\mathcal H_{\Sp}.
\]
For \(X\in\mathcal H_{\Sp}\), put
\[
  X_{\mathrm{dR}}:=\operatorname{dR}(X),
\]
and write
\[
  \eta_X^{\mathrm{dR}}:X\longrightarrow X_{\mathrm{dR}}
\]
for the localization unit.
\end{definition}

\begin{lemma}[Target descent for de Rham equivalences]
\label{lem:absolute-target-descent-de-Rham}
Let
\[
  f:A\longrightarrow B
\]
be a morphism in \(\mathcal H_{\Sp}\), let
\(e:B'\to B\) be an effective epimorphism, and put
\(A':=A\times_BB'\). If \(A'\to B'\) is a de Rham equivalence, then
\(A\to B\) is a de Rham equivalence.
\end{lemma}

\begin{proof}
\leavevmode
\begin{enumerate}[label=\textup{(\arabic*)}]
    \item \textbf{Fix the indicated data.}
    Let \(B'_\bullet\to B\) be the \v{C}ech nerve of \(e\).
    \item \textbf{Use the stated hypothesis.}
    Since the
    local-equivalence class of the left-exact localization \(L^{\mathrm{dR}}\) is
    pullback-stable, every map
    \[
    A\times_BB'_n\longrightarrow B'_n
    \]
    is a de Rham equivalence.
    \item \textbf{Treat the stated case.}
    If \(F\) is de Rham-local, effective descent in the
    ambient infinity-topos gives
    \[
    \Map(B,F)
    \simeq
    \operatorname*{Tot}_n\Map(B'_n,F)
    \]
    and
    \[
    \Map(A,F)
    \simeq
    \operatorname*{Tot}_n
    \Map(A\times_BB'_n,F).
    \]
    \item \textbf{Pass to totalization.}
    The degreewise comparison is an equivalence, hence so is its totalization.
    \item \textbf{Apply orthogonality.}
    Orthogonality to all de Rham-local objects therefore identifies \(A\to B\) as a
    de Rham equivalence.
\end{enumerate}
\end{proof}

\begin{proposition}[Global finite-infinitesimal arrows are de Rham equivalences]
\label{prop:global-finite-infinitesimal-de-Rham}
Every square-zero infinitesimal substitution of spectral schemes is a
de Rham equivalence. Consequently every finite spectral infinitesimal
substitution of
\Ref{def:finite-spectral-infinitesimal-substitution}
is a de Rham equivalence.
\end{proposition}

\begin{proof}
\leavevmode
\begin{enumerate}[label=\textup{(\arabic*)}]
    \item \textbf{Fix the indicated data.}
    Let \(j:U\to T\) be a global square-zero infinitesimal substitution.
    \item \textbf{Construct the indicated map.}
    Choose an affine spectral Zariski covering morphism
    \[
    e:\coprod_\alpha T_\alpha\longrightarrow T.
    \]
    \item \textbf{Use the stated hypothesis.}
    Since \(j\) is affine, each pullback
    \[
    U_\alpha:=U\times_TT_\alpha\longrightarrow T_\alpha
    \]
    is an affine square-zero infinitesimal substitution, hence belongs to
    \(\Sigma_{\mathrm{sqz}}\) up to equivalence and is inverted by
    \(L^{\mathrm{dR}}\).
    \item \textbf{Establish the indicated equivalence.}
    Their coproduct is therefore a de Rham equivalence.
    \item \textbf{Apply the cited lemma.}
    Lemma~\ref{lem:absolute-target-descent-de-Rham} descends this conclusion along
    \(e\), proving that \(j\) is a de Rham equivalence.
    \item \textbf{Establish the next claim.}
    The finite-composite statement follows from strong saturation and
    \Ref{def:finite-spectral-infinitesimal-substitution}.
\end{enumerate}
\end{proof}

\begin{corollary}[Global universal property of de Rham localization]
\label{cor:global-universal-property-de-Rham}
The localization
\[
  L^{\mathrm{dR}}:
  \mathcal H_{\Sp}
  \rightleftarrows
  \mathcal H_{\Sp}^{\mathrm{dR}}
  :\iota^{\mathrm{dR}}
\]
is initial among accessible left-exact localizations of
\(\mathcal H_{\Sp}\) which invert every global square-zero infinitesimal
substitution of spectral schemes.
\end{corollary}

\begin{proof}
\leavevmode
\begin{enumerate}[label=\textup{(\arabic*)}]
    \item \textbf{Verify inversion of global square-zero arrows.}
    Proposition~\ref{prop:global-finite-infinitesimal-de-Rham} shows that
    \(L^{\mathrm{dR}}\) itself inverts every global square-zero
    infinitesimal substitution of spectral schemes.
    \item \textbf{Prove the converse implication.}
    Conversely, any accessible left-exact localization which inverts every such global
    square-zero infinitesimal substitution in particular inverts every affine square-zero
    generator in \(\Sigma_{\mathrm{sqz}}\).
    \item \textbf{Use the stated construction.}
    The initiality statement of
    Theorem~\ref{thm:intrinsic-de-Rham-localization} therefore supplies the required
    universal factorization.
\end{enumerate}
\end{proof}

\begin{theorem}[de Rham locality criterion]
\label{thm:de-Rham-locality}
For \(F\in\mathcal H_{\Sp}\), the following are equivalent.
\begin{enumerate}
\item \(F\) is \(L^{\mathrm{dR}}\)-local.
\item For every affine square-zero infinitesimal substitution
\(\Spec A\to\Spec A'\), restriction induces
\[
  F(\Spec A')\xrightarrow{\simeq}F(\Spec A).
\]
\item Restriction is an equivalence along every finite spectral infinitesimal
substitution of \Ref{def:finite-spectral-infinitesimal-substitution}.
\end{enumerate}
\end{theorem}

\begin{proof}
\leavevmode
\begin{enumerate}[label=\textup{(\arabic*)}]
    \item \textbf{Apply orthogonality.}
    Conditions \textup{(1)} and \textup{(2)} are the orthogonality characterization
    of local objects for the localization generated by
    \(\Sigma_{\mathrm{sqz}}\).
    \item \textbf{Treat the stated case.}
    If \(F\) is local and \(j:U\to T\) is any global
    finite spectral infinitesimal substitution, then
    Proposition~\ref{prop:global-finite-infinitesimal-de-Rham} makes \(j\) a de Rham
    equivalence, so
    \[
    \Map(T,F)\xrightarrow{\simeq}\Map(U,F).
    \]
    \item \textbf{Establish the next claim.}
    This proves \textup{(1)}\(\Rightarrow\)\textup{(3)}.
    \item \textbf{Prove the converse implication.}
    Conversely, \textup{(3)} contains all affine square-zero infinitesimal
    substitutions and hence implies \textup{(2)}.
\end{enumerate}
\end{proof}

\begin{proposition}[Universal mapping property]
\label{prop:de-Rham-universal-mapping}
If \(Z\) is de Rham-local, composition with \(\eta_X^{\mathrm{dR}}\) induces a natural
equivalence
\[
  \Map_{\mathcal H_{\Sp}}(X_{\mathrm{dR}},Z)
  \xrightarrow{\simeq}
  \Map_{\mathcal H_{\Sp}}(X,Z).
\]
\end{proposition}

\begin{proof}
\leavevmode
\begin{enumerate}[label=\textup{(\arabic*)}]
    \item \textbf{Construct the indicated map.}
    This is the mapping-space adjunction for
    \(L^{\mathrm{dR}}\dashv\iota^{\mathrm{dR}}\).
\end{enumerate}
\end{proof}

\subsection{Recognition by finite-infinitesimal invariance}
\label{sec:finite-crystals}

The locality criterion can now be compared directly with the infinitesimal
endpoint semantics of the differentiated doctrine. On an infinitesimal arrow
the canonical endpoint comparison is nothing other than restriction from the
thickening to its centre. Thus the previous chapter already contains a
pointwise recognition test for the new reflector.

Let
\[
  i_!\dashv i^*\dashv i_*\dashv i^!
\]
be the endpoint adjoint string of
Chapter~\Ref{chap:differentiation-spectral-doctrine}. The canonical comparison
\(i_!\to i_*\) gives
\[
  \gamma_F:i_!F\longrightarrow i_*F.
\]
On a finite infinitesimal arrow \(j:U\to T\), the endpoint formulas give
\[
  (i_!F)(j)\simeq F(T),
  \qquad
  (i_*F)(j)\simeq F(U),
\]
and \(\gamma_F(j)\) is the restriction map \(F(T)\to F(U)\).

Thus the endpoint comparison is read pointwise through
\[
\begin{tikzcd}[column sep=huge,row sep=large]
  (i_!F)(j)
    \arrow[r,"\gamma_F(j)"]
    \arrow[d,"\simeq"']
  &
  (i_*F)(j)
    \arrow[d,"\simeq"]
  \\
  F(T)
    \arrow[r]
  &
  F(U).
\end{tikzcd}
\]


\begin{remark}[Comparison with crystals]
\label{rem:crystal-recognition-only}
Finite-infinitesimal invariance is the spectral analogue of the invariance condition
underlying the de Rham-prestack formulation of crystals; compare
\cite{GaitsgoryRozenblyumCrystals}.
\end{remark}

\begin{theorem}[Finite-infinitesimal recognition of de Rham locality]
\label{thm:crystalline-recognition}
For \(F\in\mathcal H_{\Sp}\), the following are equivalent:
\begin{enumerate}
\item \(F\) is de Rham-local;
\item \(F\) is invariant under every finite spectral infinitesimal substitution;
\item
\[
  \gamma_F:i_!F\longrightarrow i_*F
\]
is an equivalence in the infinitesimal Zariski topos \(\mathcal H_{\inf}\).
\end{enumerate}
Hence \(\mathcal H_{\Sp}^{\mathrm{dR}}\) is the reflective full subcategory
characterized by finite-infinitesimal invariance.
\end{theorem}

\begin{proof}
\leavevmode
\begin{enumerate}[label=\textup{(\arabic*)}]
    \item \textbf{Identify the stated objects.}
    The pointwise endpoint formulas identify condition \textup{(3)} with invariance
    along every affine finite infinitesimal substitution, since the infinitesimal Zariski topos of
    Chapter~\Ref{chap:differentiation-spectral-doctrine} is built from the affine
    finite-infinitesimal arrow category.
    \item \textbf{Establish the locality claim.}
    Such invariance contains the affine
    square-zero generators and therefore implies de Rham locality by
    Theorem~\ref{thm:de-Rham-locality}.
    \item \textbf{Prove the converse implication.}
    Conversely, a de Rham-local object is
    invariant under every global finite infinitesimal substitution by the same theorem,
    hence in particular on the affine endpoint category.
    \item \textbf{Establish the indicated equivalence.}
    This proves the equivalence
    of all three conditions.
\end{enumerate}
\end{proof}

\begin{corollary}[De Rham locality and terminal formal \texorpdfstring{\'etaleness}{etaleness}]
An object \(F\) is de Rham-local if and only if its terminal morphism is formally
\(\acute{e}\)tale with respect to the finite spectral infinitesimal substitutions of
Chapter~\Ref{chap:differentiation-spectral-doctrine}.
\end{corollary}


The absolute reflector now has both its universal property and its intrinsic
recognition criterion. To use it as part of the spectral doctrine, however,
the construction must live in every slice and commute with substitution. The
relative object to be justified in the next subsection is
\[
  X_{\mathrm{dR}/S}
  :=
  S\times_{S_{\mathrm{dR}}}X_{\mathrm{dR}},
  \qquad
  \eta_{X/S}^{\mathrm{dR}}:X\longrightarrow X_{\mathrm{dR}/S}.
\]
The next step gives the indexed form of the same localization.
\subsection{Relative reflection and substitution in context}
\label{sec:relative-de-Rham}

\begin{definition}[Relative de Rham object]
For \(X\to S\) in \(\mathcal H_{\Sp}\), define
\[
  X_{\mathrm{dR}/S}
  :=
  S\times_{S_{\mathrm{dR}}}X_{\mathrm{dR}}.
\]
The induced map \(\eta_{X/S}^{\mathrm{dR}}:X\to X_{\mathrm{dR}/S}\) is the relative de Rham
unit.
\end{definition}

\begin{construction}[Slice adjunction]
\label{con:de-Rham-slice-adjunction}
For fixed \(S\), applying absolute de Rham reflection to structure maps gives
\[
  L^{\mathrm{dR}}_S:
  (\mathcal H_{\Sp})_{/S}
  \longrightarrow
  (\mathcal H_{\Sp}^{\mathrm{dR}})_{/S_{\mathrm{dR}}},
  \qquad
  (X\to S)\longmapsto(X_{\mathrm{dR}}\to S_{\mathrm{dR}}).
\]
Pullback along \(\eta_S^{\mathrm{dR}}:S\to S_{\mathrm{dR}}\) gives
\[
  \iota^{\mathrm{dR}}_S:
  (\mathcal H_{\Sp}^{\mathrm{dR}})_{/S_{\mathrm{dR}}}
  \longrightarrow
  (\mathcal H_{\Sp})_{/S},
  \qquad
  E\longmapsto S\times_{S_{\mathrm{dR}}}E.
\]
For \(X\to S\) and local \(E\to S_{\mathrm{dR}}\), the pullback universal
property and the absolute localization adjunction give a canonical equivalence
\[
  \Map_S(X,\iota^{\mathrm{dR}}_SE)
  \simeq
  \Map_{S_{\mathrm{dR}}}(X_{\mathrm{dR}},E).
\]
Thus
\[
  L^{\mathrm{dR}}_S\dashv\iota^{\mathrm{dR}}_S.
\]

The relative reflector sits over the absolute one in the diagram
\[
\begin{tikzcd}[column sep=huge,row sep=large]
  (\mathcal H_{\Sp})_{/S}
    \arrow[r,"L^{\mathrm{dR}}_S"]
    \arrow[d]
  &
  (\mathcal H_{\Sp}^{\mathrm{dR}})_{/S_{\mathrm{dR}}}
    \arrow[d]
  \\
  \mathcal H_{\Sp}
    \arrow[r,"L^{\mathrm{dR}}"']
  &
  \mathcal H_{\Sp}^{\mathrm{dR}}.
\end{tikzcd}
\]

\end{construction}

\begin{lemma}[Local-map criterion]
\label{lem:local-map-orthogonality}
Let \(L:\mathcal Y\to\mathcal Y_L\) be a left-exact localization of an
infinity-topos with local-equivalence class \(W\). For a morphism \(u:U\to X\),
the following are equivalent:
\begin{enumerate}
\item \(u\) is right orthogonal to every morphism of \(W\);
\item the localization-unit square
\[
\begin{tikzcd}
U \ar[r,"\eta_U"] \ar[d,"u"'] & LU \ar[d,"Lu"]\\
X \ar[r,"\eta_X"'] & LX
\end{tikzcd}
\]
is Cartesian.
\end{enumerate}
\end{lemma}

\begin{proof}
\leavevmode
\begin{enumerate}[label=\textup{(\arabic*)}]
    \item \textbf{Treat the stated case.}
    If the square is Cartesian, then \(u\) is a pullback of the map
    \(Lu:LU\to LX\) between local objects.
    \item \textbf{Apply orthogonality.}
    A map between local objects is right
    orthogonal to every local equivalence, and right orthogonality is preserved by
    pullback, so \(u\) is right orthogonal to \(W\).
    \item \textbf{Prove the converse implication.}
    Conversely, put
    \[
    P:=X\times_{LX}LU,
    \]
    let \(\alpha:U\to P\) be the comparison, and let \(p:P\to X\) be the
    projection.
    \item \textbf{Apply left exactness.}
    Left exactness gives
    \[
    LP
    \simeq
    LX\times_{LLX}LLU
    \simeq
    LU,
    \]
    so \(L\alpha\) is an equivalence and hence \(\alpha\in W\).
    \item \textbf{Use pullback stability.}
    The map \(p\)
    is a pullback of \(Lu\), so \(p\) is right orthogonal to \(W\).
    \item \textbf{Construct the indicated map.}
    Applying
    \(\alpha\perp u\) to the square with upper map \(\id_U\) and lower map \(p\)
    produces
    \[
    s:P\longrightarrow U
    \]
    with
    \[
    s\alpha\simeq\id_U,
    \qquad
    us\simeq p.
    \]
    \item \textbf{Establish the next claim.}
    Now apply \(\alpha\perp p\).
    \item \textbf{Use the stated construction.}
    The maps \(\id_P\) and \(\alpha s\) are two
    fillers of the same lifting problem, so contractibility of the filler space gives
    \[
    \alpha s\simeq\id_P.
    \]
    \item \textbf{Conclude the stated claim.}
    Thus \(\alpha\) is an equivalence, which is precisely Cartesianity of the unit
    square.
\end{enumerate}
\end{proof}

\begin{theorem}[Slice de Rham localization]
\label{thm:slice-de-Rham}
The functor \(\iota^{\mathrm{dR}}_S\) of Construction~\ref{con:de-Rham-slice-adjunction} is fully
faithful. Its composite with the left adjoint defines an accessible left-exact
idempotent endofunctor
\[
  \operatorname{dR}_S
  :=\iota^{\mathrm{dR}}_SL^{\mathrm{dR}}_S:
  (\mathcal H_{\Sp})_{/S}
  \longrightarrow
  (\mathcal H_{\Sp})_{/S},
\]
with
\[
  \operatorname{dR}_S(X)
  \simeq
  X_{\mathrm{dR}/S}.
\]
Its essential image, denoted
\[
  \bigl((\mathcal H_{\Sp})_{/S}\bigr)^{\mathrm{dR}},
\]
is canonically equivalent to
\[
  (\mathcal H_{\Sp}^{\mathrm{dR}})_{/S_{\mathrm{dR}}}.
\]
An object \(F\to S\) belongs to this essential image if and only if, for every
affine square-zero infinitesimal substitution \(U\to T\) over \(S\),
\[
  \Map_S(T,F)\xrightarrow{\simeq}\Map_S(U,F).
\]
Equivalently, the same holds for every finite spectral infinitesimal substitution over \(S\).
\end{theorem}

\begin{proof}
\leavevmode
\begin{enumerate}[label=\textup{(\arabic*)}]
    \item \textbf{Use the stated construction.}
    For local \(E\to S_{\mathrm{dR}}\), left exactness gives
    \[
    \operatorname{dR}\bigl(S\times_{S_{\mathrm{dR}}}E\bigr)
    \simeq
    \operatorname{dR}(S)
    \times_{\operatorname{dR}(S_{\mathrm{dR}})}
    \operatorname{dR}(E)
    \simeq
    S_{\mathrm{dR}}\times_{S_{\mathrm{dR}}}E
    \simeq E.
    \]
    \item \textbf{Conclude the stated claim.}
    Thus the counit
    \(L^{\mathrm{dR}}_S\iota^{\mathrm{dR}}_S\to\id\) is an equivalence, and
    \(\iota^{\mathrm{dR}}_S\) is fully faithful.
    \item \textbf{Establish the next claim.}
    The formula for
    \(\operatorname{dR}_S\) follows from its definition.
    \item \textbf{Use pullback stability.}
    Accessibility follows from accessibility of \(L^{\mathrm{dR}}\) and pullback,
    and left exactness follows from left exactness of \(L^{\mathrm{dR}}\) and finite
    limits in the slice.
    \item \textbf{Identify the stated objects.}
    It remains to identify the local objects intrinsically.
    \item \textbf{Fix the indicated data.}
    Let
    \(F\to S\) be an object of the slice.
    \item \textbf{Invoke the cited result.}
    By
    Lemma~\ref{lem:local-map-orthogonality}, \(F\) lies in the essential image of
    \(\iota^{\mathrm{dR}}_S\) if and only if its structure map \(F\to S\) is right orthogonal to every
    absolute de Rham equivalence.
    \item \textbf{Use the stated hypothesis.}
    Since
    \[
    W^{\mathrm{dR}}=\operatorname{Sat}(\Sigma_{\mathrm{sqz}}),
    \]
    and, for a fixed right map, the class of left maps orthogonal to it is strongly
    saturated, this is equivalent to right orthogonality to every affine square-zero
    generator equipped with a structure map of its target to \(S\).
    \item \textbf{Construct the indicated map.}
    Written in the
    slice, that condition is exactly
    \[
    \Map_S(T,F)\xrightarrow{\simeq}\Map_S(U,F)
    \]
    for every affine square-zero infinitesimal substitution \(U\to T\) over \(S\).
    \item \textbf{Establish the next claim.}
    After the
    standing universe enlargement these tests form a small set, so they also exhibit
    \(L^{\mathrm{dR}}_S\) as the corresponding set-generated localization of the
    presentable slice.
    \item \textbf{Establish the next claim.}
    Finally, every global finite infinitesimal substitution over \(S\) has underlying
    arrow in \(W^{\mathrm{dR}}\) by
    Proposition~\ref{prop:global-finite-infinitesimal-de-Rham}.
    \item \textbf{Conclude the stated claim.}
    Hence every relative
    local object is invariant under all such arrows.
    \item \textbf{Prove the converse implication.}
    Conversely the finite class
    contains the affine square-zero generators, so finite-infinitesimal invariance
    implies the displayed affine criterion.
\end{enumerate}
\end{proof}

\begin{proposition}[Relative universal mapping property]
\label{prop:relative-de-Rham-universal-mapping}
If \(Z\to S\) is de Rham-local over \(S\), composition with the relative unit induces
\[
  \Map_S(X_{\mathrm{dR}/S},Z)
  \xrightarrow{\simeq}
  \Map_S(X,Z).
\]
The equivalence is natural in both variables.
\end{proposition}

\begin{proof}
\leavevmode
\begin{enumerate}[label=\textup{(\arabic*)}]
    \item \textbf{Construct the indicated map.}
    This is the mapping-space adjunction for the reflective subcategory of
    Theorem~\ref{thm:slice-de-Rham}.
\end{enumerate}
\end{proof}

\begin{theorem}[Arbitrary base change]
\label{thm:relative-de-Rham-base-change}
For every Cartesian square
\[
\begin{tikzcd}
X' \ar[r] \ar[d] & X \ar[d]\\
S' \ar[r] & S,
\end{tikzcd}
\]
there is a canonical equivalence
\[
  X'_{\mathrm{dR}/S'}
  \xrightarrow{\simeq}
  X_{\mathrm{dR}/S}\times_SS'.
\]
The comparison is the one supplied by finite-limit functoriality.
\end{theorem}

\begin{proof}
\leavevmode
\begin{enumerate}[label=\textup{(\arabic*)}]
    \item \textbf{Apply left exactness.}
    Left exactness gives
    \[
    \operatorname{dR}(X')
    \simeq
    \operatorname{dR}(X)
    \times_{\operatorname{dR}(S)}
    \operatorname{dR}(S').
    \]
    \item \textbf{Use the stated construction.}
    Substituting this into the defining pullback of
    \(X'_{\mathrm{dR}/S'}\) and cancelling the repeated pullback over
    \(\operatorname{dR}(S')\) gives the result.
    \item \textbf{Establish the next claim.}
    The coherence is the coherence of
    finite-limit functoriality.
\end{enumerate}
\end{proof}

\begin{theorem}[Indexed spectral de Rham reflection]
\label{thm:indexed-spectral-de-Rham-reflection}
The fibrewise reflectors
\[
  \operatorname{dR}_S:
  (\mathcal H_{\Sp})_{/S}
  \longrightarrow
  (\mathcal H_{\Sp})_{/S}
\]
assemble with substitution into an indexed accessible left-exact idempotent
reflection on the spectral slice doctrine. More precisely, for every
\(c:S'\to S\) there is a canonical equivalence
\[
  c^*\operatorname{dR}_S
  \xrightarrow{\simeq}
  \operatorname{dR}_{S'}c^*,
\]
These equivalences are represented by the commuting square of endofunctors
\[
\begin{tikzcd}[column sep=huge,row sep=large]
  (\mathcal H_{\Sp})_{/S}
    \arrow[r,"\operatorname{dR}_S"]
    \arrow[d,"c^*"']
  &
  (\mathcal H_{\Sp})_{/S}
    \arrow[d,"c^*"]
  \\
  (\mathcal H_{\Sp})_{/S'}
    \arrow[r,"\operatorname{dR}_{S'}"']
  &
  (\mathcal H_{\Sp})_{/S'}
    \arrow[ul,phantom,"\scriptstyle\simeq" description].
\end{tikzcd}
\]

with the pullback coherence fixed by the standing convention. The de Rham-local
objects over \(S\) therefore form a
reflective Cartesian subdoctrine of
\[
  S\longmapsto(\mathcal H_{\Sp})_{/S}.
\]

For a judgment \(X\to S\), the de Rham reflection rules are represented by
\[
\begin{array}{rcl}
\textup{formation} &:& X\longmapsto X_{\mathrm{dR}/S},\\[2mm]
\textup{introduction} &:& \eta_{X/S}^{\mathrm{dR}}:X\to X_{\mathrm{dR}/S},\\[2mm]
\textup{elimination} &:&
\Map_S(X_{\mathrm{dR}/S},Z)\simeq\Map_S(X,Z)
\quad(Z\text{ de Rham-local over }S),\\[2mm]
\textup{substitution} &:&
(X\times_SS')_{\mathrm{dR}/S'}
\simeq
X_{\mathrm{dR}/S}\times_SS'.
\end{array}
\]
\end{theorem}

\begin{proof}
\leavevmode
\begin{enumerate}[label=\textup{(\arabic*)}]
    \item \textbf{Identify the stated objects.}
    For each fixed \(S\), Theorem~\ref{thm:slice-de-Rham} constructs the accessible
    left-exact idempotent reflector and identifies its reflective full subcategory.
    \item \textbf{Apply the cited proposition.}
    Proposition~\ref{prop:relative-de-Rham-universal-mapping} supplies the displayed
    mapping property.
    \item \textbf{Apply the cited theorem.}
    Theorem~\ref{thm:relative-de-Rham-base-change} supplies the substitution
    equivalence.
    \item \textbf{Establish the locality claim.}
    These comparisons preserve the local essential images and assemble
    the fibrewise reflective subcategories into the asserted Cartesian subdoctrine.
    \item \textbf{Establish the next claim.}
    Formation and introduction are given by the relative de Rham object and its unit.
\end{enumerate}
\end{proof}


We have therefore turned the finite infinitesimal test geometry into an
indexed left-exact reflection. The reader may now retain two facts:
de Rham-local objects are exactly the objects invariant under finite
infinitesimal substitution, and relative reflection commutes with arbitrary
base change. What is still missing is an effective presentation of the
relative unit. That is a different problem, because a localization unit need
not itself be an effective epimorphism.
\section{Effective de Rham quotients and the infinitesimal relation}
\label{sec:de-Rham-factorization-equality}

Relative reflection gives the universal map from \(X\) to a de Rham-local object.
Effective descent is governed by the effective image of this unit: the unit can
contain points of \(X_{\mathrm{dR}/S}\) beyond those generated effectively by
points of \(X\). We therefore factor the relative unit before forming its
relation.
\[
  X\xrightarrow{e_{X/S}}Q_{X/S}
   \xrightarrow{m_{X/S}}X_{\mathrm{dR}/S}.
\]
\[
\begin{tikzcd}[column sep=huge,row sep=large]
  X
    \arrow[rr,"\eta_{X/S}^{\mathrm{dR}}"]
    \arrow[dr,two heads,"e_{X/S}"']
  &&
  X_{\mathrm{dR}/S}
  \\
  &
  Q_{X/S}
    \arrow[ur,hook,"m_{X/S}"']
  &
\end{tikzcd}
\]

The two factors play different roles. The monomorphism records how far the
effective image is from the full reflected object, while the effective
epimorphism is the map whose \v{C}ech nerve admits effective descent. This
section first controls that factorization under base change and relates its
right orthogonality to formal \(\acute{e}\)taleness. It then identifies the
kernel-pair nerve of \(X\to Q_{X/S}\) with the simplicial relation already
defined by fibre products over \(X_{\mathrm{dR}/S}\). Finally, de Rham
neighbourhoods express the same relation as an infinitesimal closure of the
diagonal.

The decisive comparison is the following pullback identity, proved below from
the monomorphism \(Q_{X/S}\hookrightarrow X_{\mathrm{dR}/S}\):
\[
\begin{tikzcd}[column sep=huge,row sep=large]
  X\times_{Q_{X/S}}X
      \arrow[r,"\simeq"]
      \arrow[d,shift left=1.1ex]
      \arrow[d,shift right=1.1ex]
  &
  X\times_{X_{\mathrm{dR}/S}}X
      \arrow[d,shift left=1.1ex]
      \arrow[d,shift right=1.1ex]
  \\
  X \arrow[r,equal] & X .
\end{tikzcd}
\]
It is this equivalence, rather than an identification
\(Q_{X/S}\simeq X_{\mathrm{dR}/S}\), that makes the de Rham relation an
effective groupoid in complete generality.
\subsection{The effective image of the relative unit}
\label{sec:effective-de-Rham-quotient}

Fix \(X\to S\).

\begin{definition}[Effective de Rham quotient]
Factor the relative de Rham unit by the effective-epimorphism--monomorphism
factorization in the slice \cite{LurieHigherToposTheory}:
\[
  X\xrightarrow{e_{X/S}}Q_{X/S}
   \xrightarrow{m_{X/S}}X_{\mathrm{dR}/S}.
\]
The object \(Q_{X/S}\) is the \emph{effective de Rham quotient} of \(X\)
over \(S\).

The factorization is a morphism in the slice:
\[
\begin{tikzcd}[column sep=large,row sep=large]
  X
    \arrow[r,two heads,"e_{X/S}"]
    \arrow[dr,"\eta_{X/S}^{\mathrm{dR}}"']
  &
  Q_{X/S}
    \arrow[d,hook,"m_{X/S}"]
  \\
  &
  X_{\mathrm{dR}/S}.
\end{tikzcd}
\]

\end{definition}

\begin{proposition}[Functoriality and base change of the effective quotient]
\label{prop:Q-base-change}
The assignment \(X\mapsto Q_{X/S}\) is functorial in the slice. For every
\(S'\to S\), with \(X'=X\times_SS'\), there is a canonical equivalence
\[
  Q_{X'/S'}\xrightarrow{\simeq}Q_{X/S}\times_SS',
\]
under arbitrary iterated base change.
\end{proposition}

\begin{proof}
\leavevmode
\begin{enumerate}[label=\textup{(\arabic*)}]
    \item \textbf{Establish the next claim.}
    The effective-image factorization is the functorial factorization attached to the
    \((\mathrm{effective\ epimorphism},\mathrm{monomorphism})\) factorization system.
    \item \textbf{Apply naturality.}
    Naturality of the
    relative unit gives functoriality.
    \item \textbf{Apply the cited theorem.}
    Theorem~\ref{thm:relative-de-Rham-base-change}
    identifies the base-changed target, and both effective epimorphisms and monomorphisms
    are pullback-stable.
    \item \textbf{Use pullback stability.}
    The pullback of the factorization is therefore the effective
    image factorization of the base-changed unit.
\end{enumerate}
\end{proof}

\subsubsection{Locality and lifting through the effective image}

To use \(Q_{X/S}\) in deformation-theoretic arguments, we need a criterion for
its locality that can be checked after passing to an effective cover. The
next lemmas isolate that criterion and identify the slice-local equivalences
with the strongly saturated class generated by square-zero substitutions.

\begin{lemma}[Target locality of de Rham equivalences over a base]
\label{lem:target-locality-de-Rham-equivalences}
Let \(f:A\to B\) be a morphism over \(S\), let \(e:B'\to B\) be an effective
epimorphism, and put \(A'=A\times_BB'\). If \(A'\to B'\) is a de Rham equivalence
over \(S\), then \(A\to B\) is a de Rham equivalence over \(S\).
\end{lemma}

\begin{proof}
\leavevmode
\begin{enumerate}[label=\textup{(\arabic*)}]
    \item \textbf{Fix the indicated data.}
    Let \(B'_\bullet\to B\) be the \v{C}ech nerve of \(e\).
    \item \textbf{Establish the indicated equivalence.}
    The local-equivalence
    class of the left-exact slice localization is pullback-stable, so every map
    \[
    A\times_BB'_n\longrightarrow B'_n
    \]
    is a de Rham equivalence over \(S\).
    \item \textbf{Use the stated construction.}
    For \(F\to S\) de Rham-local over \(S\),
    effective descent in the slice gives
    \[
    \Map_S(B,F)
    \simeq
    \operatorname*{Tot}_n\Map_S(B'_n,F)
    \]
    and
    \[
    \Map_S(A,F)
    \simeq
    \operatorname*{Tot}_n\Map_S(A\times_BB'_n,F).
    \]
    \item \textbf{Pass to totalization.}
    The degreewise comparison is an equivalence, hence so is its totalization.
    \item \textbf{Apply orthogonality.}
    This is
    the orthogonality criterion for \(A\to B\) to be a de Rham equivalence over \(S\).
\end{enumerate}
\end{proof}

\begin{lemma}[Relative image-locality criterion]
\label{lem:relative-image-locality}
Let \(m:Q\to Y\) be a monomorphism over \(S\), with \(Y\to S\) de Rham-local
over \(S\). Suppose that for every affine square-zero infinitesimal substitution \(V\to T\) over
\(S\) and every \(q_0:V\to Q\), the contractibly unique extension
\(a:T\to Y\) of \(mq_0\) factors spectral-Zariski-locally on \(T\) through
\(Q\). Then \(Q\to S\) is de Rham-local over \(S\).
\end{lemma}

\begin{proof}
\leavevmode
\begin{enumerate}[label=\textup{(\arabic*)}]
    \item \textbf{Establish the next claim.}
    For \(a:T\to Y\), write \(a_0\) for its restriction to \(V\) and define
    \[
    \Phi_T(a)
    :=
    \operatorname{fib}_a
    \bigl(\Map_S(T,Q)\to\Map_S(T,Y)\bigr),
    \]
    \[
    \Phi_V(a_0)
    :=
    \operatorname{fib}_{a_0}
    \bigl(\Map_S(V,Q)\to\Map_S(V,Y)\bigr).
    \]
    \item \textbf{Use the stated hypothesis.}
    Since \(m\) is a monomorphism, these fibres are \((-1)\)-truncated.
    \item \textbf{Establish the next claim.}
    A point of
    \(\Phi_V(a_0)\) is precisely a factorization \(q_0:V\to Q\).
    \item \textbf{Invoke the cited result.}
    By hypothesis the
    unique extension \(a:T\to Y\) factors locally through \(Q\).
    \item \textbf{Identify the stated objects.}
    On overlaps two local
    factorizations have the same composite to \(Y\), and the monomorphism property makes their
    identification space contractible.
    \item \textbf{Establish the next claim.}
    They therefore descend to a global factorization.
    \item \textbf{Conclude the stated claim.}
    Thus \(\Phi_T(a)\to\Phi_V(a_0)\) is an equivalence for every \(a\).
    \item \textbf{Construct the indicated map.}
    Together
    with locality of \(Y\), this makes
    \[
    \begin{tikzcd}
    \Map_S(T,Q) \ar[r] \ar[d] & \Map_S(T,Y) \ar[d,"\simeq"]\\
    \Map_S(V,Q) \ar[r] & \Map_S(V,Y)
    \end{tikzcd}
    \]
    \item \textbf{Establish the indicated equivalence.}
    Cartesian, so \(\Map_S(T,Q)\to\Map_S(V,Q)\) is an equivalence.
    \item \textbf{Apply the indicated result.}
    Apply
    Theorem~\ref{thm:slice-de-Rham}.
\end{enumerate}
\end{proof}

\begin{proposition}[De Rham equivalences over a base are square-zero generated]
\label{prop:relative-de-Rham-generated-squarezero}
For fixed \(S\), the local-equivalence class of \(L^{\mathrm{dR}}_S\) is the
strongly saturated class generated by a small set of affine square-zero
infinitesimal substitutions over \(S\).
\end{proposition}

\begin{proof}
\leavevmode
\begin{enumerate}[label=\textup{(\arabic*)}]
    \item \textbf{Use the stated construction.}
    Localizing the presentable slice at those arrows gives precisely the local objects of
    Theorem~\ref{thm:slice-de-Rham}.
    \item \textbf{Establish the indicated equivalence.}
    Accessible reflective localizations with the same
    local objects are canonically equivalent, and the local-equivalence class of a
    set-generated localization is its strongly saturated closure.
\end{enumerate}
\end{proof}

\subsubsection{Formal \texorpdfstring{\'etale}{etale} rigidity and smooth effectivity}

The locality criterion now converts the formal lifting notions of the
differentiated doctrine into statements about the relative de Rham unit.
Formal \(\acute{e}\)taleness gives Cartesianity and rigidity; formal smoothness
gives the complementary existence statement that makes the unit effective.

\begin{theorem}[Formal \texorpdfstring{\'etale}{etale} Cartesianity]
\label{thm:formal-etale-cartesianity}
Let \(u:U\to X\) be formally \(\acute{e}\)tale over \(S\) in the sense of
Chapter~\Ref{chap:differentiation-spectral-doctrine}. Then
\[
\begin{tikzcd}
U \ar[r,"u"] \ar[d,"\eta_{U/S}^{\mathrm{dR}}"'] & X \ar[d,"\eta_{X/S}^{\mathrm{dR}}"]\\
U_{\mathrm{dR}/S} \ar[r] & X_{\mathrm{dR}/S}
\end{tikzcd}
\]
is Cartesian. Equivalently,
\[
  U
  \simeq
  X\times_{X_{\mathrm{dR}/S}}U_{\mathrm{dR}/S}.
\]
\end{theorem}

\begin{proof}
\leavevmode
\begin{enumerate}[label=\textup{(\arabic*)}]
    \item \textbf{Invoke the cited result.}
    By one-stage detection in
    Chapter~\Ref{chap:differentiation-spectral-doctrine}, formal
    \(\acute{e}\)taleness is right orthogonality to the affine square-zero
    substitutions.
    \item \textbf{Construct the indicated map.}
    For a fixed right map, the class of left maps orthogonal to it is
    strongly saturated.
    \item \textbf{Apply orthogonality.}
    Proposition~\ref{prop:relative-de-Rham-generated-squarezero}
    therefore gives right orthogonality to every de Rham equivalence over \(S\).
    \item \textbf{Apply the indicated result.}
    Apply
    Lemma~\ref{lem:local-map-orthogonality} to the slice localization.
\end{enumerate}
\end{proof}

\begin{theorem}[Formal \texorpdfstring{\'etale}{etale} rigidity]
\label{thm:formal-etale-de-Rham}
If \(X\to S\) is formally \(\acute{e}\)tale, then
\[
  \eta_{X/S}^{\mathrm{dR}}:X\xrightarrow{\simeq}X_{\mathrm{dR}/S}.
\]
\end{theorem}

\begin{proof}
\leavevmode
\begin{enumerate}[label=\textup{(\arabic*)}]
    \item \textbf{Apply the indicated result.}
    Apply Theorem~\ref{thm:formal-etale-cartesianity} to
    \(X\to S\) regarded as a morphism in the slice over \(S\).
    \item \textbf{Use the stated hypothesis.}
    Since
    \[
    S_{\mathrm{dR}/S}
    =
    S\times_{S_{\mathrm{dR}}}S_{\mathrm{dR}}
    \simeq S,
    \]
    the resulting Cartesian square identifies
    \[
    X
    \simeq
    S\times_{S_{\mathrm{dR}/S}}X_{\mathrm{dR}/S}
    \simeq
    X_{\mathrm{dR}/S}.
    \]
    \item \textbf{Identify the stated objects.}
    Under this identification the comparison is the relative localization unit
    \(\eta_{X/S}^{\mathrm{dR}}\).
    \item \textbf{Establish the locality claim.}
    Equivalently, one may apply the locality criterion of
    Theorem~\ref{thm:slice-de-Rham}.
\end{enumerate}
\end{proof}

\begin{theorem}[Effectivity for formally smooth morphisms]
\label{thm:formal-smooth-effectivity}
If \(X\to S\) is formally smooth in the sense of
Chapter~\Ref{chap:differentiation-spectral-doctrine}, then
\[
  \eta_{X/S}^{\mathrm{dR}}:X\longrightarrow X_{\mathrm{dR}/S}
\]
is an effective epimorphism.
\end{theorem}

\begin{proof}
\leavevmode
\begin{enumerate}[label=\textup{(\arabic*)}]
    \item \textbf{Establish the next claim.}
    Factor the unit as
    \[
    X\xrightarrow{e}Q\xrightarrow{m}Y,
    \qquad
    Y:=X_{\mathrm{dR}/S},
    \]
    with \(e\) an effective epimorphism and \(m\) a monomorphism.
    \item \textbf{Establish the locality claim.}
    We first prove that \(Q\to S\) is
    de Rham-local over \(S\).
    \item \textbf{Fix the indicated data.}
    Let \(V\to T\) be an affine square-zero infinitesimal substitution over \(S\),
    and let \(q_0:V\to Q\).
    \item \textbf{Use the stated construction.}
    Pulling back \(e\) gives an effective epimorphism
    \[
    V\times_QX\longrightarrow V.
    \]
    \item \textbf{Invoke the cited result.}
    By Convention~\Ref{conv:spectral-effective-epimorphisms}, effective
    epimorphisms in the spectral Zariski topos are locally surjective on sections.
    \item \textbf{Conclude the stated claim.}
    Hence \(q_0\) lifts on a spectral Zariski covering sieve of \(V\).
    \item \textbf{Use the stated hypothesis.}
    Since
    \(V\) is affine, \Ref{prop:principal-opens-basis} and
    \Ref{prop:finite-principal-refinement} refine this sieve to
    a finite affine spectral Zariski cover
    \[
    \{V_\alpha\to V\}
    \]
    on which \(q_0\) lifts to \(X\).
    \item \textbf{Invoke the cited result.}
    By \Ref{prop:finite-infinitesimal-zariski-invariance}, pullback along
    \(V\to T\) identifies the small spectral Zariski sites and their finite affine
    covering families.
    \item \textbf{Conclude the stated claim.}
    Thus the chosen cover is transported, through a
    contractible space, to an affine spectral Zariski cover
    \[
    \{T_\alpha\to T\},
    \qquad
    V_\alpha\simeq V\times_TT_\alpha.
    \]
    \item \textbf{Establish the next claim.}
    Formal smoothness and one-stage detection then give, after a further spectral
    Zariski refinement, extensions \(T_{\alpha\beta}\to X\).
    \item \textbf{Establish the locality claim.}
    The object \(Y\) is local, so \(mq_0:V\to Y\) has a contractibly unique extension
    \(a:T\to Y\).
    \item \textbf{Identify the stated objects.}
    On each refined chart the composite through \(X\to Y\) is another
    extension of the same closed-stage map, hence is canonically identified with
    \(a\).
    \item \textbf{Conclude the stated claim.}
    Thus \(a\) locally factors through \(Q\).
    \item \textbf{Apply the cited lemma.}
    Lemma~\ref{lem:relative-image-locality}
    shows that \(Q\) is local.
    \item \textbf{Apply the indicated result.}
    Apply the relative de Rham endofunctor to
    \(X\xrightarrow eQ\xrightarrow mY\).
    \item \textbf{Use the stated hypothesis.}
    Since \(Q\) and \(Y\) are local, the image is
    \[
    \operatorname{dR}_S(X)
    \xrightarrow{\operatorname{dR}_S(e)}Q
    \xrightarrow{m}Y.
    \]
    \item \textbf{Establish the indicated equivalence.}
    The composite is the localization of the localization unit and is an equivalence.
    \item \textbf{Conclude the stated claim.}
    Hence \(m\) admits the section
    \[
    Y
    \xrightarrow{(m\operatorname{dR}_S(e))^{-1}}
    \operatorname{dR}_S(X)
    \xrightarrow{\operatorname{dR}_S(e)}Q.
    \]
    \item \textbf{Establish the indicated equivalence.}
    A monomorphism which is a split epimorphism is an equivalence.
    \item \textbf{Conclude the stated claim.}
    Therefore \(m\) is an
    equivalence and the original unit is the effective epimorphism \(e\).
\end{enumerate}
\end{proof}

\begin{corollary}[Effectivity domain]
If \(\eta_{X/S}^{\mathrm{dR}}\) is an effective epimorphism---in particular if \(X\to S\) is formally
smooth---then
\[
  Q_{X/S}\xrightarrow{\simeq}X_{\mathrm{dR}/S}.
\]
\end{corollary}

\subsubsection{Effectivity preservation and orthogonal factorization}

The formally smooth case shows when one particular localization unit is
effective. For descent we also need reflection itself to preserve effective
epimorphisms. Once that is established, the de Rham equivalences and formally
\(\acute{e}\)tale maps assemble into the orthogonal factorization system used
throughout the remainder of the chapter.

\begin{theorem}[Relative de Rham reflection preserves effective epimorphisms]
\label{thm:de-Rham-preserves-effective-epis}
Fix \(S\). If \(e:A\to B\) is an effective epimorphism in
\((\mathcal H_{\Sp})_{/S}\), then
\[
  A_{\mathrm{dR}/S}\longrightarrow B_{\mathrm{dR}/S}
\]
is an effective epimorphism in the ambient slice. This assertion is compatible with
arbitrary base change.
\end{theorem}

\begin{proof}
\leavevmode
\begin{enumerate}[label=\textup{(\arabic*)}]
    \item \textbf{Construct the indicated map.}
    Factor the reflected map as
    \[
    A_{\mathrm{dR}/S}\xrightarrow{e'}I\xrightarrow{m}B_{\mathrm{dR}/S},
    \]
    with \(e'\) an effective epimorphism and \(m\) a monomorphism.
    \item \textbf{Establish the locality claim.}
    We first show that \(I\) is
    local.
    \item \textbf{Fix the indicated data.}
    Let \(V\to T\) be an affine square-zero infinitesimal substitution and
    let \(V\to I\) be a map.
    \item \textbf{Use pullback stability.}
    The pullback of \(e'\) to \(V\) is an effective
    epimorphism.
    \item \textbf{Invoke the cited result.}
    By Convention~\Ref{conv:spectral-effective-epimorphisms}, the
    resulting section lifts on a spectral Zariski covering sieve of \(V\).
    \item \textbf{Establish the next claim.}
    The
    refinements of \Ref{prop:principal-opens-basis} and
    \Ref{prop:finite-principal-refinement} refine that sieve to a finite affine
    spectral Zariski cover of \(V\).
    \item \textbf{Establish the indicated equivalence.}
    Transport this cover to
    \(T\) by the small-site equivalence of
    \Ref{prop:finite-infinitesimal-zariski-invariance}.
    \item \textbf{Use the stated hypothesis.}
    Since
    \(A_{\mathrm{dR}/S}\) is local, the lifted maps extend across the corresponding
    square-zero charts.
    \item \textbf{Construct the indicated map.}
    Their composites in \(B_{\mathrm{dR}/S}\) agree with the
    unique extension there; the monomorphism property of \(m\) then descends the local
    factorizations through \(I\).
    \item \textbf{Apply the cited lemma.}
    Lemma~\ref{lem:relative-image-locality} makes
    \(I\) local.
    \item \textbf{Introduce the notation.}
    Put
    \[
    J:=B\times_{B_{\mathrm{dR}/S}}I.
    \]
    \item \textbf{Construct the indicated map.}
    The projection \(j:J\to B\) is a monomorphism.
    \item \textbf{Apply naturality.}
    Naturality of the relative unit and the
    factorization of the reflected map give a factorization
    \[
    A\longrightarrow J\xrightarrow{j}B
    \]
    of the original effective epimorphism \(A\to B\).
    \item \textbf{Establish the next claim.}
    In an infinity-topos the
    \((\mathrm{effective\ epimorphism},\mathrm{monomorphism})\) factorization system is orthogonal.
    \item \textbf{Apply orthogonality.}
    Applying this orthogonality to the square
    \[
    \begin{tikzcd}
    A \ar[r] \ar[d,two heads] & J \ar[d,hook,"j"]\\
    B \ar[r,"\id_B"'] & B
    \end{tikzcd}
    \]
    produces a section \(B\to J\) of \(j\).
    \item \textbf{Establish the indicated equivalence.}
    A monomorphism with a section is an
    equivalence, so \(J\simeq B\).
    \item \textbf{Conclude the stated claim.}
    Hence the unit
    \(B\to B_{\mathrm{dR}/S}\) factors through \(I\).
    \item \textbf{Use the stated identification.}
    Because \(I\) is local, the universal property of the unit extends this
    factorization through a contractible space to a map
    \[
    s:B_{\mathrm{dR}/S}\longrightarrow I.
    \]
    \item \textbf{Establish the next claim.}
    The composites \(ms\) and \(\id_{B_{\mathrm{dR}/S}}\) have the same restriction
    along the unit \(B\to B_{\mathrm{dR}/S}\).
    \item \textbf{Use the stated hypothesis.}
    Since the target is local, the unit
    mapping equivalence makes the space of such extensions contractible; hence
    \(ms\simeq\id\).
    \item \textbf{Conclude the stated claim.}
    Thus \(m\) is a monomorphism with a section and therefore an
    equivalence.
    \item \textbf{Construct the indicated map.}
    The reflected map is an effective epimorphism.
    \item \textbf{Use pullback stability.}
    Base-change compatibility
    follows from Theorem~\ref{thm:relative-de-Rham-base-change} and pullback stability
    of effective epimorphisms.
\end{enumerate}
\end{proof}

\begin{corollary}[\v{C}ech continuity of relative de Rham reflection]
\label{cor:de-Rham-cech-continuity}
If \(A\to B\) is an effective epimorphism over \(S\), with \v{C}ech nerve
\(A_\bullet\to B\), then
\[
  (A_n)_{\mathrm{dR}/S}
  \simeq
  \underbrace{
  A_{\mathrm{dR}/S}\times_{B_{\mathrm{dR}/S}}\cdots
  \times_{B_{\mathrm{dR}/S}}A_{\mathrm{dR}/S}
  }_{n+1\text{ factors}}
\]
compatibly with all simplicial operators, and
\[
  |(A_\bullet)_{\mathrm{dR}/S}|
  \simeq
  B_{\mathrm{dR}/S}.
\]
\end{corollary}

\begin{proof}
\leavevmode
\begin{enumerate}[label=\textup{(\arabic*)}]
    \item \textbf{Apply left exactness.}
    Left exactness gives the degreewise identification, and
    Theorem~\ref{thm:de-Rham-preserves-effective-epis} makes the reflected map an effective
    epimorphism.
    \item \textbf{Establish the next claim.}
    The displayed simplicial object is therefore its \v{C}ech nerve.
\end{enumerate}
\end{proof}

\begin{proposition}[Structured infinitesimal context extensions are de Rham equivalences]
\label{prop:global-square-zero-de-Rham}
Every structured infinitesimal context extension of spectral schemes constructed in
Chapter~\Ref{chap:differentiation-spectral-doctrine} is a de Rham equivalence.
The same holds for every finite composite and in every relative slice.
\end{proposition}

\begin{proof}
\leavevmode
\begin{enumerate}[label=\textup{(\arabic*)}]
    \item \textbf{Establish the next claim.}
    Choose an affine spectral Zariski cover of the target.
    \item \textbf{Use pullback stability.}
    The pullback of the
    structured infinitesimal context extension to each affine chart is an affine square-zero
    infinitesimal substitution and is therefore inverted.
    \item \textbf{Establish the indicated equivalence.}
    The coproduct of the
    affine charts is an effective epimorphism, so
    Lemma~\ref{lem:target-locality-de-Rham-equivalences} detects the global arrow as
    a de Rham equivalence.
    \item \textbf{Use the stated construction.}
    Closure under finite composition gives the remaining
    statement.
\end{enumerate}
\end{proof}

\begin{theorem}[De Rham--formally \texorpdfstring{\'etale}{etale} factorization system]
\label{thm:de-Rham-etale-factorization-system}
For every morphism
\[
  u:U\longrightarrow X
\]
in \(\mathcal H_{\Sp}\), put
\[
  P^{\mathrm{dR}}(u)
  :=
  X\times_{X_{\mathrm{dR}}}U_{\mathrm{dR}}.
\]
The defining pullback and the factorization may be displayed together as
\[
\begin{tikzcd}[column sep=large,row sep=large]
  U
    \arrow[r,"\ell_u"]
    \arrow[dr,"u"']
  &
  P^{\mathrm{dR}}(u)
    \arrow[r]
    \arrow[d,"r_u"]
  &
  U_{\mathrm{dR}}
    \arrow[d]
  \\
  & X \arrow[r,"\eta_X^{\mathrm{dR}}"'] & X_{\mathrm{dR}}.
\end{tikzcd}
\]

There is a canonical functorial factorization
\[
  U
  \xrightarrow{\ell_u}
  P^{\mathrm{dR}}(u)
  \xrightarrow{r_u}
  X
\]
in which \(\ell_u\) is a de Rham equivalence and \(r_u\) is formally
\(\acute{e}\)tale. The de Rham equivalences and formally
\(\acute{e}\)tale morphisms form a pullback-stable orthogonal factorization
system. The same factorization restricts in every slice over a fixed base \(S\),
where the left class is the class of de Rham equivalences over \(S\) and the right
class is the class of morphisms right orthogonal to the finite spectral infinitesimal
substitutions over \(S\).
\end{theorem}

\begin{proof}
\leavevmode
\begin{enumerate}[label=\textup{(\arabic*)}]
    \item \textbf{Apply naturality.}
    Naturality of the absolute de Rham unit gives the canonical comparison
    \[
    U\longrightarrow X\times_{X_{\mathrm{dR}}}U_{\mathrm{dR}},
    \]
    followed by the projection to \(X\).
    \item \textbf{Establish the next claim.}
    This is functorial in \(u\).
    \item \textbf{Apply left exactness.}
    Left exactness and idempotence give
    \[
    \begin{aligned}
    \operatorname{dR}\bigl(P^{\mathrm{dR}}(u)\bigr)
    &\simeq
    X_{\mathrm{dR}}
    \times_{X_{\mathrm{dR}}}
    U_{\mathrm{dR}}\\
    &\simeq U_{\mathrm{dR}}.
    \end{aligned}
    \]
    \item \textbf{Establish the indicated equivalence.}
    Under this equivalence \(\operatorname{dR}(\ell_u)\) is the identity of
    \(U_{\mathrm{dR}}\).
    \item \textbf{Conclude the stated claim.}
    Hence \(\ell_u\) is a de Rham equivalence.
    \item \textbf{Identify the stated objects.}
    The unit-naturality square for \(r_u\) is, after the preceding identification,
    the defining Cartesian square
    \[
    \begin{tikzcd}
    P^{\mathrm{dR}}(u) \ar[r] \ar[d,"r_u"']
    & U_{\mathrm{dR}} \ar[d]\\
    X \ar[r,"\eta_X^{\mathrm{dR}}"'] & X_{\mathrm{dR}}.
    \end{tikzcd}
    \]
    \item \textbf{Apply the cited lemma.}
    Lemma~\ref{lem:local-map-orthogonality} therefore makes \(r_u\) right
    orthogonal to every de Rham equivalence.
    \item \textbf{Use the stated hypothesis.}
    Since the de Rham class is generated
    by the affine square-zero substitutions, one-stage detection from
    Chapter~\Ref{chap:differentiation-spectral-doctrine} identifies this right
    orthogonality with formal \(\acute{e}\)taleness.
    \item \textbf{Prove the converse implication.}
    Conversely, Theorem~\ref{thm:formal-etale-cartesianity} identifies a formally
    \(\acute{e}\)tale morphism with a Cartesian de Rham-unit square, and
    Lemma~\ref{lem:local-map-orthogonality} then makes it right orthogonal to every
    de Rham equivalence.
    \item \textbf{Conclude the stated claim.}
    Thus the right class is precisely the formally
    \(\acute{e}\)tale class.
    \item \textbf{Apply orthogonality.}
    The functorial factorization above and the
    orthogonality just proved give the asserted orthogonal factorization system.
    \item \textbf{Establish the indicated equivalence.}
    Pullback stability follows from pullback stability of de Rham equivalences and
    from base-change stability of formal lifting.
    \item \textbf{Use the stated construction.}
    Applying the same argument in a
    fixed slice, using
    Proposition~\ref{prop:relative-de-Rham-generated-squarezero}, gives the relative
    statement.
\end{enumerate}
\end{proof}

\subsection{The effective de Rham relation groupoid}
\label{sec:effective-relation}

The map \(e_{X/S}:X\to Q_{X/S}\) is effective by construction, so its
\v{C}ech nerve is the correct object for descent. The relation, however, was
defined geometrically using fibre products over \(X_{\mathrm{dR}/S}\). We
must therefore show that these two simplicial objects agree. The
monomorphism \(Q_{X/S}\hookrightarrow X_{\mathrm{dR}/S}\) supplies exactly
that comparison.

\begin{definition}[De Rham relation groupoid]
For \(n\geq0\), define
\[
  \mathcal R^{\mathrm{dR}}_{X/S,n}
  :=
  \underbrace{
  X\times_{X_{\mathrm{dR}/S}}\cdots\times_{X_{\mathrm{dR}/S}}X
  }_{n+1\text{ factors}}.
\]
Deletion and repetition of factors define a simplicial object
\(\mathcal R^{\mathrm{dR}}_{X/S,\bullet}\) over \(S\). Let
\[
  v_n:\mathcal R^{\mathrm{dR}}_{X/S,n}\longrightarrow X
\]
denote the first-vertex projection. In degree one write
\[
  \mathcal R^{\mathrm{dR}}_{X/S}:=\mathcal R^{\mathrm{dR}}_{X/S,1},
  \qquad
  \mathsf s:=v_1=\operatorname{pr}_1,
  \qquad
  \mathsf t:=\operatorname{pr}_2.
\]
We call \(\mathsf s\) the source and \(\mathsf t\) the target.

In degree one the relation is displayed by
\[
\begin{tikzcd}[column sep=huge,row sep=large]
  &
  \mathcal R^{\mathrm{dR}}_{X/S}
    \arrow[dl,bend right=15,"\mathsf s"']
    \arrow[dr,bend left=15,"\mathsf t"]
  &
  \\
  X
    \arrow[dr,"e_{X/S}"']
  &&
  X
    \arrow[dl,"e_{X/S}"]
  \\
  &
  Q_{X/S}
  &
\end{tikzcd}
\qquad
e_{X/S}\mathsf s\simeq e_{X/S}\mathsf t.
\]
Indeed, by definition of the fibre product over \(X_{\mathrm{dR}/S}\),
\[
  \eta_{X/S}^{\mathrm{dR}}\mathsf s
  \simeq
  \eta_{X/S}^{\mathrm{dR}}\mathsf t.
\]
Since \(\eta_{X/S}^{\mathrm{dR}}=m_{X/S}e_{X/S}\) and \(m_{X/S}\) is a monomorphism, its
diagonal is an equivalence and cancellation along \(m_{X/S}\) gives the displayed
homotopy
\[
  e_{X/S}\mathsf s\simeq e_{X/S}\mathsf t.
\]

\end{definition}

\begin{theorem}[Effective presentation]
\label{thm:relation-effective-presentation}
There is a canonical equivalence of simplicial objects
\[
  \mathcal R^{\mathrm{dR}}_{X/S,\bullet}
  \simeq
  \check C_\bullet
  \bigl(e_{X/S}:X\to Q_{X/S}\bigr),
\]
and hence a canonical realization equivalence
\[
  |\mathcal R^{\mathrm{dR}}_{X/S,\bullet}|
  \simeq
  Q_{X/S}.
\]
\end{theorem}

\begin{proof}
\leavevmode
\begin{enumerate}[label=\textup{(\arabic*)}]
    \item \textbf{Use the stated hypothesis.}
    Since \(m_{X/S}:Q_{X/S}\to X_{\mathrm{dR}/S}\) is a monomorphism,
    \[
    Q_{X/S}\times_{X_{\mathrm{dR}/S}}Q_{X/S}
    \simeq
    Q_{X/S}.
    \]
    \item \textbf{Use the stated construction.}
    Finite-limit pasting gives
    \[
    X\times_{X_{\mathrm{dR}/S}}X
    \simeq
    X\times_{Q_{X/S}}X,
    \]
    and the same argument in every degree gives the equivalence of nerves.
    \item \textbf{Use the stated hypothesis.}
    Since
    \(e_{X/S}\) is an effective epimorphism, its \v{C}ech nerve realizes to
    \(Q_{X/S}\).
\end{enumerate}
\end{proof}

\begin{proposition}[Effective groupoid structure]
The simplicial object \(\mathcal R^{\mathrm{dR}}_{X/S,\bullet}\) is the effective
equivalence-relation groupoid of \(X\to Q_{X/S}\), as in the general effective
descent theory of \(\infty\)-topoi \cite{LurieHigherToposTheory}. In particular it carries
canonical unit, inversion,
and composition maps
\[
  X\longrightarrow \mathcal R^{\mathrm{dR}}_{X/S},
  \qquad
  \mathcal R^{\mathrm{dR}}_{X/S}\xrightarrow{\simeq}\mathcal R^{\mathrm{dR}}_{X/S},
  \qquad
  \mathcal R^{\mathrm{dR}}_{X/S}\times_{\mathsf t,X,\mathsf s}\mathcal R^{\mathrm{dR}}_{X/S}
  \longrightarrow \mathcal R^{\mathrm{dR}}_{X/S},
\]
with all coherences supplied by the simplicial identities.

In degree two, composition is encoded by the Segal diagram
\[
\begin{tikzcd}[column sep=huge,row sep=large]
  \mathcal R^{\mathrm{dR}}_{X/S,2}
    \arrow[r,"\sim","{(d_2,d_0)}"']
    \arrow[d,"d_1"']
  &
  \mathcal R^{\mathrm{dR}}_{X/S}
    \times_X
  \mathcal R^{\mathrm{dR}}_{X/S}
    \arrow[d,"\mathrm{comp}"]
  \\
  \mathcal R^{\mathrm{dR}}_{X/S}
    \arrow[r,equal]
  &
  \mathcal R^{\mathrm{dR}}_{X/S}.
\end{tikzcd}
\]
\end{proposition}

\subsection{De Rham neighbourhoods as infinitesimal closures}
\label{sec:de-Rham-neighbourhoods}

The relation groupoid admits a second description that is useful for
transport arguments. Each relation simplex is the de Rham closure of a
diagonal inside a Cartesian power of \(X\). Formulating this closure as a
relative de Rham neighbourhood makes factorization through finite
infinitesimal thickenings and formal \(\acute{e}\)tale invariance immediate
consequences of the relative reflector.

\begin{definition}[Relative de Rham neighbourhood]
For \(Z\to Y\to S\), define
\[
  \operatorname{Nbh}^{\mathrm{dR}}_{Z/S}(Y)
  :=
  Y\times_{Y_{\mathrm{dR}/S}}Z_{\mathrm{dR}/S}.
\]
It is an object over \(Y\), equipped with the canonical map from \(Z\).

Equivalently, it is defined by the Cartesian square
\[
\begin{tikzcd}[column sep=large,row sep=large]
  \operatorname{Nbh}^{\mathrm{dR}}_{Z/S}(Y)
    \arrow[r]
    \arrow[d]
  &
  Z_{\mathrm{dR}/S}
    \arrow[d]
  \\
  Y
    \arrow[r,"\eta_{Y/S}^{\mathrm{dR}}"']
  &
  Y_{\mathrm{dR}/S}.
\end{tikzcd}
\]

\end{definition}

\begin{proposition}[Base change]
\label{prop:de-Rham-neighbourhood-base-change}
For \(S'\to S\), with \(Y'=Y\times_SS'\) and \(Z'=Z\times_SS'\),
\[
  \operatorname{Nbh}^{\mathrm{dR}}_{Z'/S'}(Y')
  \xrightarrow{\simeq}
  \operatorname{Nbh}^{\mathrm{dR}}_{Z/S}(Y)\times_SS'.
\]
\end{proposition}

\begin{proof}
\leavevmode
\begin{enumerate}[label=\textup{(\arabic*)}]
    \item \textbf{Establish the indicated equivalence.}
    Substitute the coherent relative base-change equivalences of
    Theorem~\ref{thm:relative-de-Rham-base-change} into the defining fibre product.
\end{enumerate}
\end{proof}

\begin{proposition}[Factorization through a de Rham neighbourhood]
\label{prop:de-Rham-neighbourhood-factorization}
Suppose \(Z\to T\to Y\to S\) and that \(Z\to T\) is a de Rham equivalence
over \(S\). Then \(T\to Y\) factors canonically through
\[
  T\longrightarrow
  \operatorname{Nbh}^{\mathrm{dR}}_{Z/S}(Y)
  \longrightarrow Y.
\]
The factorization is natural in the diagram and compatible with arbitrary base
change.
\end{proposition}

\begin{proof}
\leavevmode
\begin{enumerate}[label=\textup{(\arabic*)}]
    \item \textbf{Establish the indicated equivalence.}
    The hypothesis makes
    \[
    L^{\mathrm{dR}}_S(Z)\longrightarrow L^{\mathrm{dR}}_S(T)
    \]
    an equivalence.
    \item \textbf{Identify the stated objects.}
    Precomposition with this equivalence identifies
    \[
    \Map_S\bigl(T_{\mathrm{dR}/S},Z_{\mathrm{dR}/S}\bigr)
    \xrightarrow{\simeq}
    \Map_S\bigl(Z_{\mathrm{dR}/S},Z_{\mathrm{dR}/S}\bigr),
    \]
    and the identity of \(Z_{\mathrm{dR}/S}\) therefore determines, through a
    contractible space, the inverse equivalence.
    \item \textbf{Establish the next claim.}
    Compose the unit
    \[
    T\longrightarrow T_{\mathrm{dR}/S}
    \]
    with that canonically determined inverse to obtain
    \(T\to Z_{\mathrm{dR}/S}\).
    \item \textbf{Apply naturality.}
    Naturality of the localization unit identifies the
    two composites to \(Y_{\mathrm{dR}/S}\), so the universal property of the fibre
    product produces the asserted factorization.
    \item \textbf{Apply naturality.}
    Naturality and base-change
    compatibility follow from the same mapping-space construction and the coherent
    base-change equivalences for relative de Rham reflection.
\end{enumerate}
\end{proof}

\begin{theorem}[De Rham relation simplices as diagonal neighbourhoods]
\label{thm:simplices-as-neighbourhoods}
Let
\[
  Y_n:=X^{\times_S(n+1)},
  \qquad
  \Delta_n:X\to Y_n
\]
be the diagonal. Then there is a canonical equivalence over \(Y_n\)
\[
  \mathcal R^{\mathrm{dR}}_{X/S,n}
  \xrightarrow{\simeq}
  \operatorname{Nbh}^{\mathrm{dR}}_{X/S}(Y_n).
\]
These equivalences are simplicially natural, functorial in \(X/S\), and compatible
with arbitrary base change.
\end{theorem}

\begin{proof}
\leavevmode
\begin{enumerate}[label=\textup{(\arabic*)}]
    \item \textbf{Apply left exactness.}
    Left exactness gives
    \[
    (Y_n)_{\mathrm{dR}/S}
    \simeq
    (X_{\mathrm{dR}/S})^{\times_S(n+1)}.
    \]
    \item \textbf{Establish the next claim.}
    The relative de Rham image of \(\Delta_n\) is the diagonal of
    \(X_{\mathrm{dR}/S}\).
    Therefore
    \[
    Y_n\times_{(Y_n)_{\mathrm{dR}/S}}X_{\mathrm{dR}/S}
    \simeq
    X\times_{X_{\mathrm{dR}/S}}\cdots\times_{X_{\mathrm{dR}/S}}X.
    \]
    \item \textbf{Use the stated construction.}
    Deleting and repeating factors gives the simplicial compatibility.
\end{enumerate}
\end{proof}

\begin{theorem}[Formal \texorpdfstring{\'etale}{etale} invariance of de Rham neighbourhoods]
\label{thm:formal-etale-neighbourhood-invariance}
Let \(u:U\to X\) be formally \(\acute{e}\)tale over \(S\), and let
\(Z\to U\). Regard \(Z\) also as an object over \(X\). Then there is a
canonical equivalence over \(X\)
\[
  \operatorname{Nbh}^{\mathrm{dR}}_{Z/S}(U)
  \xrightarrow{\simeq}
  \operatorname{Nbh}^{\mathrm{dR}}_{Z/S}(X),
\]
where the left-hand object is viewed over \(X\) through \(U\to X\). Under this
equivalence the right-hand de Rham neighbourhood factors canonically through \(U\).
Moreover
\[
  Q_{U/S}
  \simeq
  Q_{X/S}\times_{X_{\mathrm{dR}/S}}U_{\mathrm{dR}/S}
\]
and
\[
  \mathcal R^{\mathrm{dR}}_{U/S,n}
  \simeq
  U\times_{X,v_n}\mathcal R^{\mathrm{dR}}_{X/S,n}
  \simeq
  Q_{U/S}\times_{Q_{X/S}}\mathcal R^{\mathrm{dR}}_{X/S,n},
\]
compatibly with the simplicial structure.
\end{theorem}

\begin{proof}
\leavevmode
\begin{enumerate}[label=\textup{(\arabic*)}]
    \item \textbf{Use the stated construction.}
    Formal \(\acute{e}\)tale Cartesianity gives
    \[
    U
    \simeq
    X\times_{X_{\mathrm{dR}/S}}U_{\mathrm{dR}/S}.
    \]
    Hence
    \[
    \begin{aligned}
    \operatorname{Nbh}^{\mathrm{dR}}_{Z/S}(U)
    &=
    U\times_{U_{\mathrm{dR}/S}}Z_{\mathrm{dR}/S}\\
    &\simeq
    \bigl(X\times_{X_{\mathrm{dR}/S}}U_{\mathrm{dR}/S}\bigr)
    \times_{U_{\mathrm{dR}/S}}Z_{\mathrm{dR}/S}\\
    &\simeq
    X\times_{X_{\mathrm{dR}/S}}Z_{\mathrm{dR}/S}
    =
    \operatorname{Nbh}^{\mathrm{dR}}_{Z/S}(X).
    \end{aligned}
    \]
    \item \textbf{Use pullback stability.}
    The unit of \(U\) is the pullback of the unit of \(X\) along
    \(U_{\mathrm{dR}/S}\to X_{\mathrm{dR}/S}\).
    \item \textbf{Use pullback stability.}
    Pullback stability and uniqueness
    of effective-image factorization therefore give the formula for \(Q_{U/S}\).
    \item \textbf{Establish the next claim.}
    The
    de Rham relation formulas follow by taking the corresponding \v{C}ech nerves, or directly
    by finite-limit pasting.
\end{enumerate}
\end{proof}


At this point the unstable geometry required for descent is complete. The
relative reflector supplies the universal infinitesimal quotient, its
effective image \(Q_{X/S}\) supplies the object actually presented by
generalized points, and
\(\mathcal R^{\mathrm{dR}}_{X/S,\bullet}\) is the effective relation groupoid
presenting that image. The next task is therefore to construct a stable
scalar-valued semantics on these already-defined objects.
\section{Stable scalar semantics on de Rham geometry}
\label{sec:de-Rham-stable-hyperdoctrine}

The preceding section produced the effective groupoid on which descent must
be computed. To extract a de Rham complex from it, we need stable scalar
functions that can be pulled back, pushed forward by dependent product, and
descented along effective epimorphisms. We reuse the stabilized slice fibres
from formal reconstruction,
\[
  \mathcal D_{\Sp/S}
  =
  \operatorname{Sp}\bigl(((\mathcal H_{\Sp})_{/S})_*\bigr).
\]
and establish only the interface used later: dependent products,
Beck--Chevalley, internal function objects, effective descent, and a stable
realization of quasi-coherent coefficients. The reader may keep the intended
output in view throughout this section: for a map \(a:T\to S\), stable scalar
functions will be expressible both as \(\Pi_a a^*\mathcal O_S^{\mathrm{add}}\)
and as an internal function object, and both descriptions will commute with
base change and effective descent.

Stable realization connects the two coefficient languages already present in
the manuscript. Quasi-coherent modules remain the source of geometric
coefficients, while parameterized spectra provide the ambient category in
which functions on arbitrary spectral Zariski stacks can be formed. On an
affine point these descriptions are compared by ordinary global sections:
\[
\begin{tikzcd}[column sep=huge,row sep=large]
  \QCoh(S)
      \arrow[r,"\operatorname{Real}^{\mathrm{st}}_S"]
      \arrow[d,"a^*"']
  &
  \mathcal D_{\Sp/S}
      \arrow[d,"a^*"]
  \\
  \QCoh(T)
      \arrow[r,"\operatorname{Real}^{\mathrm{st}}_T"']
      \arrow[d,"{\Gamma(T,-)}"']
  &
  \mathcal D_{\Sp/T}
      \arrow[d,"{\Gamma^{\mathrm{st}}(T,-)}"]
  \\
  \Sp \arrow[r,equal] & \Sp .
\end{tikzcd}
\]
The bottom comparison is what later identifies represented stable scalar
functions with ordinary spectral functions and first-order scalar coefficients
with cotangent coefficients.
\subsection{Stable slices and dependent adjoints}
\label{sec:parameterized-stable}

The de Rham relation objects, effective quotients, and de Rham
neighbourhoods live in \(\mathcal H_{\Sp}\), so scalar functions on them must be
formed in stabilized slices. We first establish the dependent-adjoint calculus
for those slices. The affine-point presentation is then used only for the two
jobs that require it: detecting stable function objects and realizing
quasi-coherent coefficients.

\begin{proposition}[Affine-point presentation of a spectral slice]
\label{prop:affine-point-slice-site}
For \(S\in\mathcal H_{\Sp}\), equip \(\operatorname{AffPt}(S)\) with the topology induced by finite affine spectral
Zariski covers on the affine source. Restriction along the affine-point functor induces a
canonical equivalence
\[
  (\mathcal H_{\Sp})_{/S}
  \xrightarrow{\simeq}
  \Sh_{\tau_{\mathrm{Zar}}}\bigl(\operatorname{AffPt}(S)\bigr).
\]
Under this equivalence an object \(Y\to S\) is sent to the sheaf
\[
  (h_T\to S)\longmapsto\Map_S(h_T,Y).
\]
The statement applies to every \(S\in\mathcal H_{\Sp}\).
\end{proposition}

\begin{proof}
\leavevmode
\begin{enumerate}[label=\textup{(\arabic*)}]
    \item \textbf{Use the stated construction.}
    At the presheaf level, the category-of-elements construction gives the canonical
    equivalence
    \[
    \PSh(\Aff_{\Sp})_{/S}
    \simeq
    \PSh\bigl(\operatorname{AffPt}(S)\bigr).
    \]
    \item \textbf{Use pullback stability.}
    The spectral Zariski topology on the affine-point category is the pullback of the
    ambient topology: a covering family over \(h_T\to S\) is precisely a finite affine
    spectral Zariski covering family of \(T\), with the induced maps to \(S\).
    \item \textbf{Establish the locality claim.}
    Consequently the covering-sieve generators for the two sheaf localizations
    correspond.
    \item \textbf{Use the stated construction.}
    Passing to the associated accessible left-exact localizations gives the
    displayed equivalence.
    \item \textbf{Establish the next claim.}
    The pointwise formula is the Yoneda description of the
    slice presheaf.
    \item \textbf{Use the stated construction.}
    Equivalently,
    Lemma~\ref{lem:de-Rham-affine-density} supplies the affine density needed to recover
    every object of the slice from these tests.
\end{enumerate}
\end{proof}

\begin{construction}[Parameterized stable slice doctrine]
\label{def:parameterized-stable-coefficient-categories}
For every \(S\in\mathcal H_{\Sp}\), retain the stabilized slice category
\[
  \mathcal D_{\Sp/S}
  =
  \operatorname{Sp}\bigl(((\mathcal H_{\Sp})_{/S})_*\bigr)
\]
and the notation of Chapter~\Ref{chap:formal-reconstruction-spectral-doctrine}.
Write \(\mathbf1_S\) for its tensor unit. For an object \(V\to S\), let
\[
  \Sigma^\infty_{+,S}V\in\mathcal D_{\Sp/S}
\]
denote the suspension spectrum of the free pointed object \(V_+\) in the slice.
The cartesian symmetric monoidal structure on the slice stabilizes to the canonical
presentably symmetric monoidal structure on \(\mathcal D_{\Sp/S}\) \cite{LurieHigherAlgebra}; in particular
\[
  \Sigma^\infty_{+,S}U\otimes\Sigma^\infty_{+,S}V
  \simeq
  \Sigma^\infty_{+,S}(U\times_SV).
\]
As \(S\) varies, these already-fixed fibres and their pullback functors form the
parameterized stable slice doctrine used below; see \cite{LurieHigherAlgebra}.
\end{construction}

\begin{theorem}[Stable dependent adjoints and base change]
\label{thm:stable-dependent-adjoints}
For every morphism \(f:T\to S\), slice pullback induces an exact strong symmetric
monoidal functor
\[
  f^*:\mathcal D_{\Sp/S}\longrightarrow\mathcal D_{\Sp/T}
\]
with adjoints
\[
  \Sigma_f\dashv f^*\dashv\Pi_f.
\]
For every Cartesian square
\[
\begin{tikzcd}
T' \ar[r,"g'"] \ar[d,"f'"'] & T \ar[d,"f"]\\
S' \ar[r,"g"'] & S,
\end{tikzcd}
\]
the pullback functors themselves form the square
\[
\begin{tikzcd}[column sep=huge,row sep=large]
  \mathcal D_{\Sp/S} \arrow[r,"f^*"] \arrow[d,"g^*"']
  & \mathcal D_{\Sp/T} \arrow[d,"(g')^*"]
  \\
  \mathcal D_{\Sp/S'} \arrow[r,"(f')^*"']
  & \mathcal D_{\Sp/T'}
    \arrow[ul,phantom,"\scriptstyle\simeq" description],
\end{tikzcd}
\]
and there are canonical Beck--Chevalley equivalences
\[
  g^*\Sigma_f
  \xrightarrow{\simeq}
  \Sigma_{f'}(g')^*,
  \qquad
  g^*\Pi_f
  \xrightarrow{\simeq}
  \Pi_{f'}(g')^*.
\]
Moreover there is a canonical projection equivalence
\[
  \Sigma_f\bigl(f^*E\otimes F\bigr)
  \xrightarrow{\simeq}
  E\otimes\Sigma_fF,
\]
and \(\Pi_f\) carries the canonical lax symmetric monoidal structure right-adjoint
to the strong symmetric monoidal structure on \(f^*\). Their compatibility is determined by the slice pullback diagrams and the
calculus of mates.
\end{theorem}

\begin{proof}
\leavevmode
\begin{enumerate}[label=\textup{(\arabic*)}]
    \item \textbf{Use pullback stability.}
    Unstable slice pullback
    \[
    f^*:(\mathcal H_{\Sp})_{/S}\longrightarrow(\mathcal H_{\Sp})_{/T}
    \]
    preserves all limits and colimits because it has both the dependent sum and
    dependent product as adjoints.
    \item \textbf{Use pullback stability.}
    It preserves finite limits in particular, so
    functoriality of spectrum objects induces the displayed exact stable pullback.
    \item \textbf{Establish the next claim.}
    The induced functor is accessible and preserves all limits and colimits.
    \item \textbf{Establish the next claim.}
    The
    presentable adjoint functor theorem therefore constructs both adjoints
    \[
    \Sigma_f\dashv f^*\dashv\Pi_f.
    \]
    \item \textbf{Use pullback stability.}
    The strong symmetric monoidal structure on stable pullback is the stabilization
    of the cartesian pullback structure; see \cite{LurieHigherAlgebra}.
    \item \textbf{Fix the indicated data.}
    Let \(v:V\to T\).
    \item \textbf{Use the stated construction.}
    For every \(E\in\mathcal D_{\Sp/S}\), the
    suspension-spectrum/zeroth-space adjunction in the two slices gives
    \[
    \begin{aligned}
    \map_{\mathcal D_{\Sp/S}}
    \bigl(\Sigma_f\Sigma^\infty_{+,T}V,E\bigr)
    &\simeq
    \map_{\mathcal D_{\Sp/T}}
    \bigl(\Sigma^\infty_{+,T}V,f^*E\bigr)\\
    &\simeq
    \map_{\mathcal D_{\Sp/S}}
    \bigl(\Sigma^\infty_{+,S}V,E\bigr),
    \end{aligned}
    \]
    where on the last line \(V\) is regarded over \(S\) through \(f\circ v\).
    \item \textbf{Apply Yoneda density.}
    Yoneda therefore gives the canonical equivalence
    \[
    \Sigma_f\Sigma^\infty_{+,T}V
    \simeq
    \Sigma^\infty_{+,S}V.
    \]
    \item \textbf{Establish the next claim.}
    The canonical mate
    \[
    g^*\Sigma_f\longrightarrow\Sigma_{f'}(g')^*
    \]
    is a transformation between colimit-preserving exact functors.
    \item \textbf{Establish the next claim.}
    Suspension
    spectra of objects of the slice, together with their shifts, generate the
    stabilization.
    \item \textbf{Establish the next claim.}
    On a generator \(\Sigma^\infty_{+,T}V\), the two sides are
    canonically
    \[
    \Sigma^\infty_{+,S'}(V\times_SS')
    \qquad\text{and}\qquad
    \Sigma^\infty_{+,S'}(V\times_TT'),
    \]
    which agree because the square is Cartesian.
    \item \textbf{Conclude the stated claim.}
    Hence the mate is an equivalence.
    \item \textbf{Apply the indicated result.}
    Apply the left Beck--Chevalley equivalence of the preceding step to the same
    Cartesian square with the horizontal and vertical roles exchanged.
    \item \textbf{Use the stated construction.}
    This gives
    the canonical equivalence
    \[
    f^*\Sigma_g
    \xrightarrow{\simeq}
    \Sigma_{g'}(f')^*.
    \]
    \item \textbf{Identify the stated objects.}
    Passing to right adjoints identifies the right adjoint of the left-hand
    composite with \(g^*\Pi_f\), and the right adjoint of the right-hand composite
    with \(\Pi_{f'}(g')^*\).
    \item \textbf{Use the stated construction.}
    Uniqueness of right adjoints therefore gives the
    canonical mate equivalence
    \[
    g^*\Pi_f
    \xrightarrow{\simeq}
    \Pi_{f'}(g')^*.
    \]
    \item \textbf{Establish the next claim.}
    This comparison is obtained entirely by transposition and the calculus of
    adjoint mates, so its functoriality is inherited from the left Beck--Chevalley
    equivalence.
    \item \textbf{Establish the next claim.}
    Both sides of
    \[
    \Sigma_f(f^*E\otimes F)
    \longrightarrow
    E\otimes\Sigma_fF
    \]
    preserve colimits separately in \(E\) and \(F\).
    \item \textbf{Establish the next claim.}
    It therefore suffices to
    evaluate on
    \[
    E=\Sigma^\infty_{+,S}U,
    \qquad
    F=\Sigma^\infty_{+,T}V.
    \]
    \item \textbf{Identify the stated objects.}
    Strong symmetric monoidality of suspension spectra and the calculation of
    \(\Sigma_f\) above identify both sides canonically with
    \[
    \Sigma^\infty_{+,S}(U\times_SV).
    \]
    \item \textbf{Conclude the stated claim.}
    Hence the projection morphism is an equivalence for all \(E,F\).
    \item \textbf{Use the stated hypothesis.}
    Since \(f^*\) is strong symmetric monoidal, the right adjoint \(\Pi_f\)
    receives its canonical lax symmetric monoidal structure.
    \item \textbf{Use pullback stability.}
    All comparisons above
    are mates of the functorial pullback, product, and composition diagrams.
\end{enumerate}
\end{proof}

\begin{corollary}[Parameterized stable slice hyperdoctrine recognition]
\label{cor:stable-linear-hyperdoctrine-recognition}
The parameterized stable slice doctrine
\[
  S\longmapsto\mathcal D_{\Sp/S}
\]
with reindexing \(f^*\) is, in particular, a stable presentably symmetric monoidal
slice hyperdoctrine, in the sense of indexed categorical logic; compare
\cite{JacobsCategoricalLogic}. Its dependent sum and dependent product are
\[
  \Sigma_f\dashv f^*\dashv\Pi_f,
\]
the Beck--Chevalley comparisons are equivalences, and \(\Sigma_f\) satisfies the
projection formula. All substitution, dependent-adjoint, exchange, and projection
coherences are those constructed from the slice adjunctions and their mates.
\end{corollary}

\begin{proof}
\leavevmode
\begin{enumerate}[label=\textup{(\arabic*)}]
    \item \textbf{Apply the cited theorem.}
    Theorem~\ref{thm:stable-dependent-adjoints} provides the adjoint string
    \(\Sigma_f\dashv f^*\dashv\Pi_f\), the Beck--Chevalley equivalences, and the
    projection formula.
    \item \textbf{Establish the next claim.}
    These data give the
    asserted stable presentably symmetric monoidal slice hyperdoctrine; compare the
    terminology of indexed categorical logic in \cite{JacobsCategoricalLogic}.
\end{enumerate}
\end{proof}

\subsection{Affine-point semantics and parameterized function objects}

The dependent-adjoint calculus is now available abstractly. To turn it into
a usable function theory, we identify stable slices on affine points. At this
stage we need the resulting generating family, the behaviour of pullback,
relative suspension spectra, internal function objects, and effective descent.
The Postnikov and grouplike consequences of the same presentation will be
recalled only when the additive delooping construction needs them.

\begin{proposition}[Stable affine-point presentation]
\label{prop:stable-affine-point-presentation}
For every \(S\in\mathcal H_{\Sp}\), stabilization of
Proposition~\ref{prop:affine-point-slice-site} gives a canonical equivalence
\[
  \mathcal D_{\Sp/S}
  \xrightarrow{\simeq}
  \Sh_{\tau_{\mathrm{Zar}}}
  \bigl(\operatorname{AffPt}(S);\Sp\bigr).
\]
Here
\[
  \Sh_{\tau_{\mathrm{Zar}}}
  \bigl(\operatorname{AffPt}(S);\Sp\bigr)
\]
denotes the stabilization of the spectral Zariski \(\infty\)-topos
\(\Sh_{\tau_{\mathrm{Zar}}}(\operatorname{AffPt}(S))\). The notation uses the ambient spectral Zariski stack convention fixed in
Chapter~\Ref{chap:spectral-algebraic-geometry}.

Under this equivalence, \(E\in\mathcal D_{\Sp/S}\) is sent to the sheaf of spectra
\[
  (a:T\to S)
  \longmapsto
  \Gamma^{\mathrm{st}}(T,a^*E).
\]
Consequently \(\mathcal D_{\Sp/S}\) is generated under colimits by
\[
  \Sigma^n\Sigma^\infty_{+,S}T,
  \qquad n\in\mathbb Z,
\]
where \(T\to S\) ranges over affine spectral points, and a morphism in
\(\mathcal D_{\Sp/S}\) is an equivalence if and only if it induces equivalences on
mapping spectra from all these generators.
\end{proposition}

\begin{proof}
\leavevmode
\begin{enumerate}[label=\textup{(\arabic*)}]
    \item \textbf{Establish the next claim.}
    Spectrum objects are defined by finite-limit conditions.
    \item \textbf{Use the stated hypothesis.}
    Since the equivalence of
    Proposition~\ref{prop:affine-point-slice-site} is left exact, it induces an
    equivalence on stabilizations.
    \item \textbf{Invoke the cited result.}
    By the notation fixed in the statement,
    \[
    \Sh_{\tau_{\mathrm{Zar}}}
    \bigl(\operatorname{AffPt}(S);\Sp\bigr)
    :=
    \operatorname{Sp}\!\left(
    \Sh_{\tau_{\mathrm{Zar}}}
    \bigl(\operatorname{AffPt}(S)\bigr)_*
    \right).
    \]
    \item \textbf{Conclude the stated claim.}
    Hence stabilization of the affine-point presentation gives
    \[
    \operatorname{Sp}\bigl(((\mathcal H_{\Sp})_{/S})_*\bigr)
    \simeq
    \Sh_{\tau_{\mathrm{Zar}}}
    \bigl(\operatorname{AffPt}(S);\Sp\bigr).
    \]
    \item \textbf{Establish the next claim.}
    Equivalently, the right-hand side is the category of spectrum objects satisfying
    spectral Zariski descent on the same affine-point site.
    \item \textbf{Compute affine-point values.}
    The displayed value formula follows by applying the
    suspension-spectrum/zeroth-space adjunction in the slice over each affine point.
    \item \textbf{Use the stated construction.}
    Affine density gives a set of colimit generators for the unstable slice.
    \item \textbf{Establish the next claim.}
    Their
    relative suspension spectra and all integer shifts therefore generate the
    stabilization.
    \item \textbf{Establish the next claim.}
    The final assertion is the generation criterion of a generating
    set in a presentable stable infinity-category.
\end{enumerate}
\end{proof}

\begin{proposition}[Affine-point pullback]
\label{prop:stable-affine-point-pullback-t-structure}
Let \(f:T\to S\). Under
Proposition~\ref{prop:stable-affine-point-presentation}, stable pullback is
restriction along
\[
    \operatorname{AffPt}(T)\longrightarrow\operatorname{AffPt}(S),
    \qquad
    (U\to T)\longmapsto(U\to T\to S).
  \]
\end{proposition}

\begin{proof}
\leavevmode
\begin{enumerate}[label=\textup{(\arabic*)}]
    \item \textbf{Treat the stated case.}
    If \(u:U\to T\) is affine and \(E\in\mathcal D_{\Sp/S}\), the affine-point
    value of \(f^*E\) at \(u\) is the value of \(E\) at the composite
    \(fu:U\to S\).
    \item \textbf{Construct the indicated map.}
    This is the stabilized form of the Cartesian identity
    \[
    \Map_T(U,Y\times_ST)\simeq\Map_S(U,Y),
    \]
    natural in \(Y\to S\).
    \item \textbf{Conclude the stated claim.}
    Hence stable pullback is restriction on affine-point
    sheaves of spectra.
\end{enumerate}
\end{proof}

\begin{construction}[Relative suspension spectra]
Let \(a:T\to S\). Write \(T_+\) for the free pointed object on \(T\) in the
slice over \(S\), and define
\[
  \Sigma^\infty_{+,S}T
  :=
  \Sigma^\infty_ST_+\in\mathcal D_{\Sp/S}.
\]
The dependent-sum description gives
\[
  \Sigma^\infty_{+,S}T
  \simeq
  \Sigma_a\mathbf1_T.
\]
More generally, for \(U\to T\to S\),
\[
  \Sigma_a\Sigma^\infty_{+,T}U
  \simeq
  \Sigma^\infty_{+,S}U.
\]
\end{construction}

\begin{definition}[Stable sections]
For \(E\in\mathcal D_{\Sp/T}\), define
\[
  \Gamma^{\mathrm{st}}(T,E)
  :=
  \map_{\mathcal D_{\Sp/T}}(\mathbf1_T,E).
\]
\end{definition}

\begin{proposition}[Parameterized function object]
\label{prop:function-spectrum-identity}
For \(a:T\to S\) and \(E\in\mathcal D_{\Sp/S}\), there is a canonical equivalence of
objects of \(\mathcal D_{\Sp/S}\)
\[
  \Pi_a a^*E
  \simeq
  \HHom_{\mathcal D_{\Sp/S}}
  \bigl(\Sigma^\infty_{+,S}T,E\bigr).
\]
Consequently
\[
  \map_{\mathcal D_{\Sp/S}}(\Sigma^\infty_{+,S}T,E)
  \simeq
  \Gamma^{\mathrm{st}}(T,a^*E).
\]
\end{proposition}

\begin{proof}
\leavevmode
\begin{enumerate}[label=\textup{(\arabic*)}]
    \item \textbf{Construct the indicated map.}
    For \(B\in\mathcal D_{\Sp/S}\), the projection formula and the adjunction
    \(\Sigma_a\dashv a^*\) give
    \[
    \begin{aligned}
    \map_{\mathcal D_{\Sp/S}}
    \bigl(B,\HHom(\Sigma_a\mathbf1_T,E)\bigr)
    &\simeq
    \map_{\mathcal D_{\Sp/S}}(B\otimes\Sigma_a\mathbf1_T,E)\\
    &\simeq
    \map_{\mathcal D_{\Sp/T}}(a^*B,a^*E)\\
    &\simeq
    \map_{\mathcal D_{\Sp/S}}(B,\Pi_aa^*E).
    \end{aligned}
    \]
    \item \textbf{Apply Yoneda density.}
    Yoneda identifies the internal objects.
    \item \textbf{Take the indicated specialization.}
    Taking \(B=\mathbf1_S\) gives the second
    formula.
\end{enumerate}
\end{proof}

\begin{proposition}[Closed base change for parameterized internal Homs]
\label{prop:closed-base-change-internal-hom}
Let \(c:S'\to S\), and let \(E,F\in\mathcal D_{\Sp/S}\). There is a canonical
equivalence
\[
  c^*
  \HHom_{\mathcal D_{\Sp/S}}(E,F)
  \xrightarrow{\simeq}
  \HHom_{\mathcal D_{\Sp/S'}}
  \bigl(c^*E,c^*F\bigr),
\]
natural in \(c,E,F\) and compatible with the convolution algebra structures
induced by cocommutative coalgebra structures on \(E\) and commutative algebra
structures on \(F\).
\end{proposition}

\begin{proof}
\leavevmode
\begin{enumerate}[label=\textup{(\arabic*)}]
    \item \textbf{Fix the indicated data.}
    Let \(B\in\mathcal D_{\Sp/S'}\).
    \item \textbf{Use the stated construction.}
    Using
    \(\Sigma_c\dashv c^*\), the internal-Hom adjunction, the projection formula of
    Theorem~\ref{thm:stable-dependent-adjoints}, and the same adjunction again gives
    \[
    \begin{aligned}
    \map_{\mathcal D_{\Sp/S'}}
    \left(
    B,
    c^*\HHom_{\mathcal D_{\Sp/S}}(E,F)
    \right)
    &\simeq
    \map_{\mathcal D_{\Sp/S}}
    \left(
    \Sigma_cB,
    \HHom_{\mathcal D_{\Sp/S}}(E,F)
    \right)\\
    &\simeq
    \map_{\mathcal D_{\Sp/S}}
    \bigl(
    \Sigma_cB\otimes E,
    F
    \bigr)\\
    &\simeq
    \map_{\mathcal D_{\Sp/S}}
    \bigl(
    \Sigma_c(B\otimes c^*E),
    F
    \bigr)\\
    &\simeq
    \map_{\mathcal D_{\Sp/S'}}
    \bigl(
    B\otimes c^*E,
    c^*F
    \bigr)\\
    &\simeq
    \map_{\mathcal D_{\Sp/S'}}
    \left(
    B,
    \HHom_{\mathcal D_{\Sp/S'}}
    (c^*E,c^*F)
    \right).
    \end{aligned}
    \]
    \item \textbf{Apply Yoneda density.}
    Yoneda gives the displayed equivalence.
    \item \textbf{Establish the next claim.}
    Every comparison in the calculation
    is obtained functorially from the displayed adjunctions and projection formula.
    \item \textbf{Treat the stated case.}
    If \(E\) is cocommutative coalgebraic and \(F\) is commutative algebraic,
    convolution is obtained from the diagonal/comultiplication of \(E\), the
    multiplication of \(F\), and the closed symmetric monoidal structure.
    \item \textbf{Use the stated hypothesis.}
    Since \(c^*\) is strong symmetric monoidal, it transports those structural maps
    to the corresponding maps of \(c^*E\) and \(c^*F\).
    \item \textbf{Apply naturality.}
    Naturality of the
    preceding Yoneda comparison with respect to those maps identifies the two
    convolution multiplications and all higher commutative operations.
\end{enumerate}
\end{proof}

\begin{proposition}[Multiplicative parameterized function object]
\label{prop:multiplicative-function-spectrum-identity}
Let \(a:T\to S\), and let \(A\in\CAlg(\mathcal D_{\Sp/S})\). The diagonal
\(T\to T\times_ST\) and terminal map \(T\to S\) give
\(\Sigma^\infty_{+,S}T\) its canonical cocommutative coalgebra structure. Endow
\[
  \HHom_{\mathcal D_{\Sp/S}}
  \bigl(\Sigma^\infty_{+,S}T,A\bigr)
\]
with the resulting convolution commutative algebra structure. Then the equivalence
of Proposition~\ref{prop:function-spectrum-identity} refines canonically to an
equivalence in \(\CAlg(\mathcal D_{\Sp/S})\),
\[
  \Pi_aa^*A
  \xrightarrow{\simeq}
  \HHom_{\mathcal D_{\Sp/S}}
  \bigl(\Sigma^\infty_{+,S}T,A\bigr).
\]
The refinement is natural in \(a\) and \(A\) and is compatible with Cartesian base
change.
\end{proposition}

\begin{proof}
\leavevmode
\begin{enumerate}[label=\textup{(\arabic*)}]
    \item \textbf{Use the stated hypothesis.}
    Since \(a^*\) is strong symmetric monoidal, its right adjoint \(\Pi_a\) is
    canonically lax symmetric monoidal.
    \item \textbf{Establish multiplicative compatibility.}
    The multiplication on
    \(\Pi_aa^*A\) is therefore adjoint to
    \[
    a^*(\Pi_aa^*A)^{\otimes2}
    \longrightarrow
    a^*A^{\otimes2}
    \longrightarrow
    a^*A,
    \]
    where the first arrow is the tensor product of the two counits.
    \item \textbf{Establish the next claim.}
    Under
    \[
    \Sigma_a\mathbf1_T\simeq\Sigma^\infty_{+,S}T
    \]
    and the projection formula, the mate of this composite is obtained by
    precomposing two functions with the diagonal
    \[
    T\longrightarrow T\times_ST
    \]
    and then multiplying their values in \(A\).
    \item \textbf{Establish multiplicative compatibility.}
    This is exactly convolution against
    the comultiplication
    \[
    \Sigma^\infty_{+,S}T
    \longrightarrow
    \Sigma^\infty_{+,S}T\otimes\Sigma^\infty_{+,S}T.
    \]
    \item \textbf{Identify the stated objects.}
    The same mate calculation for the \(k\)-fold diagonal
    \(T\to T^{\times_Sk}\) identifies every \(k\)-ary operation of the lax symmetric
    monoidal structure on \(\Pi_a\) with the corresponding convolution operation.
    \item \textbf{Identify the stated objects.}
    The terminal map \(T\to S\) identifies the units.
    \item \textbf{Conclude the stated claim.}
    Thus the underlying
    equivalence of Proposition~\ref{prop:function-spectrum-identity} is an equivalence
    of commutative algebras.
    \item \textbf{Apply naturality.}
    Naturality and Cartesian base-change compatibility are
    inherited from the projection formula, Beck--Chevalley equivalences, and the
    functorial diagonal maps.
\end{enumerate}
\end{proof}

\begin{proposition}[Effective descent for parameterized spectra]
\label{thm:effective-descent-stable-coefficients-expanded}
If \(e:T\to S\) is an effective epimorphism with \v{C}ech nerve \(T_\bullet\),
pullback induces an equivalence of symmetric monoidal presentable stable
infinity-categories
\[
  \mathcal D_{\Sp/S}
  \xrightarrow{\simeq}
  \operatorname*{Tot}_{[n]\in\Delta}\mathcal D_{\Sp/T_n}.
\]
Under this equivalence tensor units, tensor products, pullbacks, and commutative
algebra objects are identified degreewise. Dependent adjoints and internal Homs are
then identified by their universal properties and the Beck--Chevalley equivalences
of Theorem~\ref{thm:stable-dependent-adjoints}. Pullback along \(e\) is
conservative on parameterized spectra.
\end{proposition}

\begin{proof}
\leavevmode
\begin{enumerate}[label=\textup{(\arabic*)}]
    \item \textbf{Apply effective descent.}
    Effective descent in the infinity-topos gives, as in \cite{LurieHigherToposTheory},
    \[
    (\mathcal H_{\Sp})_{/S}
    \simeq
    \operatorname*{Tot}_{[n]\in\Delta}
    (\mathcal H_{\Sp})_{/T_n}.
    \]
    \item \textbf{Take the indicated specialization.}
    Taking spectrum objects commutes with this limit because the spectrum-object
    conditions are finite-limit conditions.
    \item \textbf{Identify the stated objects.}
    The fiberwise symmetric monoidal structure
    is functorial for pullback, so the resulting equivalence is symmetric monoidal and
    identifies tensor units, tensor products, pullbacks, and commutative algebra objects
    degreewise.
    \item \textbf{Establish the indicated equivalence.}
    The dependent adjoints and internal Homs are characterized by their
    mapping universal properties and hence are transported by the same equivalence.
    \item \textbf{Establish the indicated equivalence.}
    To prove conservativity, let \(u:E\to F\) in \(\mathcal D_{\Sp/S}\) become an
    equivalence after pullback to \(T\).
    \item \textbf{Use pullback stability.}
    Every structural map
    \(T_n\to S\) of the \v{C}ech nerve factors through \(T\to S\), so the
    pullback of \(u\) to every \(T_n\) is an equivalence.
    \item \textbf{Pass to totalization.}
    Under the totalization
    equivalence the image of \(u\) is therefore degreewise an equivalence, hence an
    equivalence.
    \item \textbf{Conclude the stated claim.}
    Thus \(u\) itself is an equivalence.
\end{enumerate}
\end{proof}

\subsection{Stable realization of quasi-coherent coefficients}

The parameterized stable slice can now express functions on arbitrary
geometric objects, but the cotangent comparison later requires a bridge from
quasi-coherent coefficients to that stable slice. We construct this bridge by
taking quasi-coherent global sections on every affine point and descending the
resulting sheaf of spectra.

\begin{construction}[Section stacks over spectral Zariski stacks]
\label{con:section-stack-spectral-stack}
Let \(S\in\mathcal H_{\Sp}\) and \(\mathcal F\in\QCoh(S)\). Define
\[
  \operatorname{Sect}_S(\mathcal F)\longrightarrow S
\]
by the affine-point formula: over \(a:T\to S\), its fibre is
\[
  \Map_{\QCoh(T)}
  \bigl(\mathcal O_T,a^*\mathcal F\bigr).
\]
It is a spectral Zariski stack over \(S\), has a canonical zero section, and for
every \(f:T\to S\) there is a canonical Cartesian base-change equivalence
\[
  \operatorname{Sect}_S(\mathcal F)\times_ST
  \xrightarrow{\simeq}
  \operatorname{Sect}_T(f^*\mathcal F).
\]
More generally, for every \(g:Y\to S\),
\[
  \Map_S\bigl(Y,\operatorname{Sect}_S(\mathcal F)\bigr)
  \xrightarrow{\simeq}
  \Map_{\QCoh(Y)}\bigl(\mathcal O_Y,g^*\mathcal F\bigr).
\]
\end{construction}

\begin{proof}
\leavevmode
\begin{enumerate}[label=\textup{(\arabic*)}]
    \item \textbf{Establish the next claim.}
    The descent and base-change proof of
    \Ref{con:section-stack-qcoh} and
    \Ref{prop:section-stack-descent-base-change} uses only that the base is a spectral
    Zariski sheaf and that quasi-coherent coefficients on affine tests satisfy Zariski
    descent.
    \item \textbf{Construct the indicated map.}
    Those conditions hold for every \(S\in\mathcal H_{\Sp}\), so the
    same Grothendieck construction and mapping-space totalization prove the first two
    claims.
    \item \textbf{Establish the next claim.}
    For a general \(Y\to S\), apply affine density to \(Y\) and use the
    affine-point limit definition of \(\QCoh(Y)\).
    \item \textbf{Construct the indicated map.}
    Both sides become the same limit of
    the affine mapping spaces
    \[
    \Map_{\QCoh(U)}\bigl(\mathcal O_U,(g|_U)^*\mathcal F\bigr).
    \]
    \item \textbf{Conclude the stated claim.}
    Thus the construction extends the section-stack interface of
    Chapter~\Ref{chap:spectral-algebraic-geometry} to arbitrary bases in
    \(\mathcal H_{\Sp}\).
\end{enumerate}
\end{proof}

\begin{construction}[Stable realization of quasi-coherent coefficients]
\label{con:stable-realization-qcoh}
Let \(S\in\mathcal H_{\Sp}\) and \(\mathcal F\in\QCoh(S)\). On the
affine-point site consider the presheaf of spectra
\[
  (a:T\to S)
  \longmapsto
  \Gamma(T,a^*\mathcal F)
  =
  \map_{\QCoh(T)}(\mathcal O_T,a^*\mathcal F).
\]
For a morphism of affine points \(g:T'\to T\), restriction is the map on global
sections induced by quasi-coherent pullback.

The construction is governed by the affine-point triangle
\[
\begin{tikzcd}[column sep=huge,row sep=large]
  T' \arrow[r,"g"] \arrow[dr,"a'"']
  & T \arrow[d,"a"]
  \\
  & S,
\end{tikzcd}
\qquad
\Gamma(T,a^*\mathcal F)
\longrightarrow
\Gamma(T',(a')^*\mathcal F).
\]

This presheaf is a spectral Zariski sheaf. Indeed, if \(T_\bullet\to T\) is the
\v{C}ech nerve of an affine spectral Zariski cover, quasi-coherent descent gives a
symmetric monoidal equivalence
\[
  \QCoh(T)
  \simeq
  \operatorname*{Tot}_{[n]\in\Delta}\QCoh(T_n),
\]
under which \(\mathcal O_T\) and \(a^*\mathcal F\) correspond to their
restrictions. Mapping spectra in a categorical limit are computed by the
corresponding limit, so
\[
\begin{aligned}
  \Gamma(T,a^*\mathcal F)
  &\simeq
  \operatorname*{Tot}_{[n]\in\Delta}
  \Gamma(T_n,a_n^*\mathcal F).
\end{aligned}
\]
By Proposition~\ref{prop:stable-affine-point-presentation}, let
\[
  \operatorname{Real}^{\mathrm{st}}_S(\mathcal F)
  \in
  \mathcal D_{\Sp/S}
\]
denote the canonical object corresponding to this sheaf of spectra.
\end{construction}

\begin{proposition}[Properties and monoidality of stable realization]
\label{prop:stable-realization-qcoh}
The construction
\[
  \operatorname{Real}^{\mathrm{st}}_S:\QCoh(S)\longrightarrow\mathcal D_{\Sp/S}
\]
is exact and canonically lax symmetric monoidal.
Its affine-point evaluation is encoded by
\[
\begin{tikzcd}[column sep=huge,row sep=large]
  \QCoh(S)
    \arrow[r,"\operatorname{Real}^{\mathrm{st}}_S"]
    \arrow[d,"a^*"']
  &
  \mathcal D_{\Sp/S}
    \arrow[d,"a^*"]
  \\
  \QCoh(T)
    \arrow[r,"\operatorname{Real}^{\mathrm{st}}_T"']
  &
  \mathcal D_{\Sp/T} ,
\end{tikzcd}
\]
which commutes up to the canonical base-change equivalence. More precisely:
\begin{enumerate}
\item for every \(f:T\to S\), there is a coherent base-change equivalence
\[
  f^*\operatorname{Real}^{\mathrm{st}}_S(\mathcal F)
  \xrightarrow{\simeq}
  \operatorname{Real}^{\mathrm{st}}_T(f^*\mathcal F);
\]
\item stable global sections recover quasi-coherent global sections,
\[
  \Gamma^{\mathrm{st}}
  \bigl(S,\operatorname{Real}^{\mathrm{st}}_S(\mathcal F)\bigr)
  \xrightarrow{\simeq}
  \Gamma(S,\mathcal F);
\]
\item there are canonical natural maps
\[
  \operatorname{Real}^{\mathrm{st}}_S(\mathcal F)\otimes\operatorname{Real}^{\mathrm{st}}_S(\mathcal G)
  \longrightarrow
  \operatorname{Real}^{\mathrm{st}}_S(\mathcal F\otimes\mathcal G),
  \qquad
  \mathbf1_S\longrightarrow\operatorname{Real}^{\mathrm{st}}_S(\mathcal O_S),
\]
which satisfy the symmetric-monoidal associativity, symmetry, and unit coherences.
The base-change equivalences are compatible with this lax symmetric monoidal
structure.
\end{enumerate}

\end{proposition}

\begin{proof}
\leavevmode
\begin{enumerate}[label=\textup{(\arabic*)}]
    \item \textbf{Establish the next claim.}
    Exactness may be checked on the stable affine-point presentation.
    \item \textbf{Construct the indicated map.}
    For every affine
    point \(a:T\to S\), the functor
    \[
    \mathcal F\longmapsto
    \map_{\QCoh(T)}(\mathcal O_T,a^*\mathcal F)
    \]
    is exact: pullback is exact, and mapping spectra out of a fixed object preserve
    finite limits, hence finite colimits in a stable category.
    \item \textbf{Invoke the cited result.}
    By Proposition~\ref{prop:stable-affine-point-pullback-t-structure}, stable
    pullback is restriction of the affine-point sheaf of spectra.
    \item \textbf{Conclude the stated claim.}
    Hence for an affine
    point \(u:U\to T\),
    \[
    \begin{aligned}
    \Gamma^{\mathrm{st}}
    \bigl(U,u^*f^*\operatorname{Real}^{\mathrm{st}}_S(\mathcal F)\bigr)
    &\simeq
    \Gamma(U,(fu)^*\mathcal F)\\
    &\simeq
    \Gamma(U,u^*f^*\mathcal F),
    \end{aligned}
    \]
    which is exactly the affine-point value of
    \(\operatorname{Real}^{\mathrm{st}}_T(f^*\mathcal F)\).
    \item \textbf{Apply naturality.}
    Naturality in \(u\)
    and the stable affine-point presentation give the claimed equivalence.
    \item \textbf{Apply the indicated result.}
    Apply Lemma~\ref{lem:de-Rham-affine-density} to the terminal object
    \(S\to S\) of the slice:
    \[
    S
    \simeq
    \operatorname*{colim}_{(T,a)\in\operatorname{AffPt}(S)}T.
    \]
    \item \textbf{Establish the next claim.}
    The free-pointed functor on the slice and the relative suspension-spectrum
    functor are left adjoints, hence preserve this colimit.
    \item \textbf{Use the stated hypothesis.}
    Since
    \(\Sigma^\infty_{+,S}S\simeq\mathbf1_S\), we obtain canonically
    \[
    \mathbf1_S
    \simeq
    \operatorname*{colim}_{(T,a)\in\operatorname{AffPt}(S)}
    \Sigma^\infty_{+,S}T.
    \]
    Therefore
    \[
    \begin{aligned}
    \Gamma^{\mathrm{st}}
    \bigl(S,\operatorname{Real}^{\mathrm{st}}_S(\mathcal F)\bigr)
    &\simeq
    \operatorname*{lim}_{(T,a)\in\operatorname{AffPt}(S)^{\op}}
    \Gamma(T,a^*\mathcal F).
    \end{aligned}
    \]
    \item \textbf{Identify the stated objects.}
    The affine-point limit definition of \(\QCoh(S)\) and creation of mapping
    spectra in categorical limits identify the right side with
    \(\Gamma(S,\mathcal F)\).
    \item \textbf{Use the stated construction.}
    Finally, on an affine point \(a:T\to S\), tensoring two maps out of the tensor unit
    gives the canonical pairing
    \[
    \begin{aligned}
    &\map_{\QCoh(T)}(\mathcal O_T,a^*\mathcal F)
    \otimes
    \map_{\QCoh(T)}(\mathcal O_T,a^*\mathcal G)
    \\
    &\qquad\longrightarrow
    \map_{\QCoh(T)}
    \bigl(\mathcal O_T,a^*(\mathcal F\otimes\mathcal G)\bigr),
    \end{aligned}
    \]
    using \(\mathcal O_T\otimes\mathcal O_T\simeq\mathcal O_T\).
    \item \textbf{Use the stated construction.}
    The identity of
    \(\mathcal O_T\) gives the unit map.
    \item \textbf{Use pullback stability.}
    These pairings are natural under affine
    pullback because quasi-coherent pullback is strong symmetric monoidal, and their
    associativity, symmetry, and unit homotopies are precisely those of the symmetric
    monoidal coefficient categories.
    \item \textbf{Construct the indicated map.}
    On the affine-point site they therefore define a
    map from the objectwise smash-product presheaf to the target sheaf; by the universal
    property of spectral Zariski sheafification this factors canonically through the
    tensor product in the sheaf category.
    \item \textbf{Construct the indicated map.}
    The same construction for all arities and all
    operadic structure maps assembles to the displayed lax symmetric monoidal structure
    on \(\operatorname{Real}^{\mathrm{st}}_S\).
\end{enumerate}
\end{proof}

\begin{corollary}[Section-stack recovery from stable realization]
\label{cor:stable-realization-section-stack}
For every \(n\geq0\), there is a canonical equivalence over \(S\)
\[
  \Omega^\infty
  \bigl(\operatorname{Real}^{\mathrm{st}}_S(\mathcal F)[n]\bigr)
  \xrightarrow{\simeq}
  \operatorname{Sect}_S(\mathcal F[n]).
\]
Thus the shifted section stacks are the successive spaces of the spectrum object
\(\operatorname{Real}^{\mathrm{st}}_S(\mathcal F)\).
\end{corollary}

\begin{proof}
\leavevmode
\begin{enumerate}[label=\textup{(\arabic*)}]
    \item \textbf{Construct the indicated map.}
    On an affine point \(a:T\to S\), both sides have fibre
    \[
    \Map_{\QCoh(T)}
    \bigl(\mathcal O_T,a^*\mathcal F[n]\bigr).
    \]
    \item \textbf{Establish the indicated equivalence.}
    The affine-point presentation and Yoneda give the equivalence, with the loop
    compatibilities inherited from suspension in the same mapping spectra.
\end{enumerate}
\end{proof}

\subsection{The stable additive scalar}
\label{sec:additive-scalar}

\begin{definition}[Stable additive scalar]
The tensor unit \(\mathcal O_S\in\QCoh(S)\) carries its canonical commutative
algebra structure. Applying the lax symmetric monoidal functor of
Proposition~\ref{prop:stable-realization-qcoh} defines
\[
  \mathcal O_S^{\mathrm{add}}
  :=
  \operatorname{Real}^{\mathrm{st}}_S(\mathcal O_S)
  \in
  \CAlg(\mathcal D_{\Sp/S}).
\]
Equivalently, the scalar is obtained by applying stable realization to the tensor
unit:
\[
\begin{tikzcd}[column sep=huge]
  \mathcal O_S \arrow[r,mapsto,"\operatorname{Real}^{\mathrm{st}}_S"]
  & \mathcal O_S^{\mathrm{add}}.
\end{tikzcd}
\]

On an affine point \(a:T=\Spec A\to S\), its value is canonically the
\(E_\infty\)-ring spectrum
\[
  \Gamma(T,\mathcal O_T)\simeq A,
\]
and its multiplication is the multiplication of this affine coordinate
\(E_\infty\)-ring. The resulting coherences are inherited functorially from the
lax symmetric monoidal realization.
\end{definition}

\begin{proposition}[Existence and base change of the stable additive scalar]
\label{prop:additive-scalar-base-change}
For every \(f:T\to S\), there is a canonical equivalence of commutative algebra
objects
\[
  f^*\mathcal O_S^{\mathrm{add}}
  \xrightarrow{\simeq}
  \mathcal O_T^{\mathrm{add}},
\]
with functoriality supplied by the standing coherence convention.
\end{proposition}

\begin{proof}
\leavevmode
\begin{enumerate}[label=\textup{(\arabic*)}]
    \item \textbf{Invoke the cited result.}
    By Proposition~\ref{prop:stable-realization-qcoh} and pullback of the structure
    sheaf from Chapter~\Ref{chap:spectral-algebraic-geometry},
    \[
    \begin{aligned}
    f^*\mathcal O_S^{\mathrm{add}}
    &=f^*\operatorname{Real}^{\mathrm{st}}_S(\mathcal O_S)\\
    &\simeq\operatorname{Real}^{\mathrm{st}}_T(f^*\mathcal O_S)\\
    &\simeq\operatorname{Real}^{\mathrm{st}}_T(\mathcal O_T)
    =\mathcal O_T^{\mathrm{add}}.
    \end{aligned}
    \]
    \item \textbf{Establish the indicated equivalence.}
    The base-change equivalence is compatible with the lax symmetric monoidal structure
    of \(\operatorname{Real}^{\mathrm{st}}\), so it is multiplicative.
    \item \textbf{Use pullback stability.}
    Coherence is inherited from
    functorial quasi-coherent pullback.
\end{enumerate}
\end{proof}

\begin{proposition}[Recognition by the brave new affine line]
\label{prop:additive-scalar-affine-line-recognition}
If \(S\) is represented by a spectral scheme, the grouplike commutative group
object underlying \(\mathcal O_S^{\mathrm{add}}\) is canonically the additive
group object underlying the brave new affine line
\[
  \mathbb A^1_S
  =
  \Spec_S\operatorname{Sym}_{\mathcal O_S}(\mathcal O_S).
\]
Equivalently,
\[
  \Omega^\infty\mathcal O_S^{\mathrm{add}}
  \xrightarrow{\simeq}
  (\mathbb A^1_S,+)
\]
in the slice over \(S\).
\end{proposition}

\begin{proof}
\leavevmode
\begin{enumerate}[label=\textup{(\arabic*)}]
    \item \textbf{Identify the stated objects.}
    Corollary~\ref{cor:stable-realization-section-stack} identifies
    \(\Omega^\infty\mathcal O_S^{\mathrm{add}}\) with
    \(\operatorname{Sect}_S(\mathcal O_S)\).
    \item \textbf{Identify the stated objects.}
    The structure sheaf is finite locally
    free of rank one, so
    \Ref{prop:vector-bundle-total-space-sections} identifies this section stack with
    \[
    \mathbf V_S(\mathcal O_S)
    \simeq
    \mathbb L_S(\mathcal O_S)
    =
    \mathbb A^1_S,
    \]
    where the middle equivalence is the canonical self-duality of the rank-one free
    module.
    \item \textbf{Identify the stated objects.}
    The affine-point identification is compatible with addition and
    multiplication because both operations are induced by those of the affine
    coordinate \(E_\infty\)-rings.
\end{enumerate}
\end{proof}


\begin{proposition}[Stable scalar sections]
\label{prop:stable-scalar-sections-represented}
For every morphism \(y:Y\to S\) in \(\mathcal H_{\Sp}\), there is a
canonical natural equivalence of \(E_\infty\)-ring spectra
\[
  \Gamma^{\mathrm{st}}
  \bigl(Y,y^*\mathcal O_S^{\mathrm{add}}\bigr)
  \xrightarrow{\simeq}
  \Gamma(Y,\mathcal O_Y).
\]

\end{proposition}

\begin{proof}
\leavevmode
\begin{enumerate}[label=\textup{(\arabic*)}]
    \item \textbf{Apply the cited proposition.}
    Proposition~\ref{prop:additive-scalar-base-change} identifies
    \[
    y^*\mathcal O_S^{\mathrm{add}}
    \simeq
    \mathcal O_Y^{\mathrm{add}}
    =
    \operatorname{Real}^{\mathrm{st}}_Y(\mathcal O_Y).
    \]
    \item \textbf{Apply the cited proposition.}
    Proposition~\ref{prop:stable-realization-qcoh} gives the underlying spectrum
    equivalence.
    \item \textbf{Use the stated construction.}
    Its lax symmetric monoidal construction is defined pointwise by the
    same multiplication of structure-sheaf sections that gives the
    \(E_\infty\)-ring structure on \(\Gamma(Y,\mathcal O_Y)\); hence the equivalence
    is multiplicative.
\end{enumerate}
\end{proof}


The stable infrastructure now reduces to the promised interface. For every
map \(a:T\to S\), stable scalar functions may be written either as
\(\Pi_a a^*\mathcal O_S^{\mathrm{add}}\) or as the corresponding internal
function object; these descriptions are multiplicative, commute with base
change, and satisfy effective descent. Nothing else from the stable-slice
construction is needed to define de Rham cochains. The next section applies
this interface to the effective relation and returns to the main argument.
\section{De Rham cochains and effective infinitesimal descent}
\label{sec:de-Rham-scalar-descent}

We now have both ingredients required for de Rham descent: the effective
relation groupoid
\(\mathcal R^{\mathrm{dR}}_{X/S,\bullet}\) and the stable additive scalar
\(\mathcal O_S^{\mathrm{add}}\). The next construction is therefore forced.
We take stable scalar functions degreewise on the relation and use effective
descent along \(e_{X/S}:X\to Q_{X/S}\) to identify their totalization.

The geometric and stable descriptions are related by a descent square:
\[
\begin{tikzcd}[column sep=huge,row sep=large]
  \mathcal R^{\mathrm{dR}}_{X/S,\bullet}
      \arrow[r,"\text{stable scalar functions}"]
      \arrow[d,"\lvert-\rvert"']
  &
  C^\bullet_{\mathrm{dR}}(X/S)
      \arrow[d,"\operatorname*{Tot}"]
  \\
  Q_{X/S}
      \arrow[r,"\text{stable scalar functions}"']
  &
  \mathrm{DR}_{X/S}.
\end{tikzcd}
\]
The \v{C}ech--Amitsur comparison and effective descent identify the lower
horizontal object with the totalization of the upper-right cosimplicial
algebra. This is also the point
at which the distinction from the full reflected object becomes algebraic:
functions on \(X_{\mathrm{dR}/S}\) restrict to functions on \(Q_{X/S}\), and
the comparison is an equivalence precisely on the effectivity domain.
\subsection{Stable scalar functions on the de Rham relation}
\label{sec:infinitesimal-cochains}

We now apply the stable scalar interface degreewise to the effective relation.
For each simplex, dependent product along its structure map to \(S\) turns the
pulled-back additive scalar into a scalar-function object on the base. The
simplicial restriction maps then become the cofaces and codegeneracies of the
desired cosimplicial algebra.

Let \(a_n:\mathcal R^{\mathrm{dR}}_{X/S,n}\to S\) be the structure map.

Each cochain degree is attached to the triangle
\[
\begin{tikzcd}[column sep=huge,row sep=large]
  \mathcal R^{\mathrm{dR}}_{X/S,n}
    \arrow[rr,"a_n"]
    \arrow[dr,"v_n"']
  &&
  S
  \\
  &
  X
    \arrow[ur,"p"']
  &
\end{tikzcd}
\]
whenever \(p:X\to S\) denotes the structure morphism.

\begin{definition}[Geometric de Rham cochains]
Define
\[
  C^n_{\mathrm{dR}}(X/S)
  :=
  \Pi_{a_n}a_n^*\mathcal O_S^{\mathrm{add}}
  \in
  \CAlg(\mathcal D_{\Sp/S}).
\]
The simplicial operators of \(\mathcal R^{\mathrm{dR}}_{X/S,\bullet}\) induce a
cosimplicial commutative algebra
\[
  C^\bullet_{\mathrm{dR}}(X/S):
  \Delta\longrightarrow\CAlg(\mathcal D_{\Sp/S}).
\]
The passage from de Rham relation simplices to cochains is therefore
\[
\begin{tikzcd}[column sep=large]
  \mathcal R^{\mathrm{dR}}_{X/S,n}
    \arrow[r,"{\Sigma^\infty_{+,S}}"]
  &
  \mathsf{R}^{\mathrm{dR},\mathrm{st}}_n(X/S)
    \arrow[r,"{\HHom(-,\mathcal O_S^{\mathrm{add}})}"]
  &
  C^n_{\mathrm{dR}}(X/S).
\end{tikzcd}
\]
This de Rham relation-simplex formula is the definition; the Amitsur description below is a
presentation theorem.
\end{definition}

\begin{proposition}[Cosimplicial structure]
\label{prop:cosimplicial-commutative-algebra-expanded}
If \(\theta:[m]\to[n]\) is a morphism of \(\Delta\), with induced simplicial map
\(\theta^*:\mathcal R^{\mathrm{dR}}_{X/S,n}\to
\mathcal R^{\mathrm{dR}}_{X/S,m}\), the cosimplicial structure map
\[
  C^m_{\mathrm{dR}}(X/S)
  \longrightarrow
  C^n_{\mathrm{dR}}(X/S)
\]
is the scalar restriction map obtained from the unit of
\((\theta^*)^*\dashv\Pi_{\theta^*}\) and composition of dependent products. These
maps are morphisms of commutative algebras and satisfy the cosimplicial identities.
\end{proposition}

\begin{proof}
\leavevmode
\begin{enumerate}[label=\textup{(\arabic*)}]
    \item \textbf{Establish the next claim.}
    The triangle
    \[
    \mathcal R^{\mathrm{dR}}_{X/S,n}
    \xrightarrow{\theta^*}
    \mathcal R^{\mathrm{dR}}_{X/S,m}
    \xrightarrow{a_m}S
    \]
    satisfies \(a_n=a_m\theta^*\).
    \item \textbf{Use pullback stability.}
    Pullback is strong symmetric monoidal, dependent
    product is lax symmetric monoidal, and composition of dependent products is
    coherent.
    \item \textbf{Use the stated construction.}
    The adjunction unit gives the stated restriction map.
    \item \textbf{Apply functoriality.}
    Functoriality of
    the units and mates gives the cosimplicial identities.
\end{enumerate}
\end{proof}

\begin{construction}[Scalar restriction along a de Rham relation simplex]
Let
\[
  r:\mathcal R^{\mathrm{dR}}_{X/S,n}\longrightarrow \mathcal R^{\mathrm{dR}}_{X/S,m}
\]
be any simplicial structure map over \(S\), so \(a_n=a_mr\). Define the
corresponding scalar-restriction morphism
\[
  \operatorname{res}_r:
  C^m_{\mathrm{dR}}(X/S)
  \longrightarrow
  C^n_{\mathrm{dR}}(X/S)
\]
to be the composite
\[
\begin{aligned}
  \Pi_{a_m}a_m^*\mathcal O_S^{\mathrm{add}}
  &\longrightarrow
  \Pi_{a_m}\Pi_r r^*a_m^*\mathcal O_S^{\mathrm{add}}\\
  &\simeq
  \Pi_{a_n}a_n^*\mathcal O_S^{\mathrm{add}},
\end{aligned}
\]
where the first arrow is induced by the unit
\(\id\to\Pi_rr^*\). These are the cosimplicial structure maps of
Proposition~\ref{prop:cosimplicial-commutative-algebra-expanded}.
\end{construction}

\begin{convention}[Ordered degree-one faces and signs]
In degree one write
\[
  \mathcal R^{\mathrm{dR}}_{X/S,1}=X\times_{X_{\mathrm{dR}/S}}X,
\]
and order the two geometric faces by
\[
  \partial_0=\mathsf t=\operatorname{pr}_2,
  \qquad
  \partial_1=\mathsf s=\operatorname{pr}_1.
\]
The corresponding cofaces are therefore the typed scalar-restriction maps
\[
  d^0=\operatorname{res}_{\mathsf t},
  \qquad
  d^1=\operatorname{res}_{\mathsf s}.
\]
For normalized cosimplicial cochains we use the cohomological convention
\[
  d_N^q
  =
  \sum_{i=0}^{q+1}(-1)^i d^i.
\]
Thus the degree-zero alternating coface difference is
\[
  d^0-d^1
  =
  \operatorname{res}_{\mathsf t}-\operatorname{res}_{\mathsf s}.
\]
The stable Dold--Kan comparison and the complete-filtered/coherent-cochain
correspondence below are read with this convention.
\end{convention}

\begin{construction}[Suspension spectra of de Rham relation simplices]
Put
\[
  \mathsf{R}^{\mathrm{dR},\mathrm{st}}_n(X/S)
  :=
  \Sigma^\infty_{+,S}\mathcal R^{\mathrm{dR}}_{X/S,n}\in\mathcal D_{\Sp/S}.
\]
The de Rham relation groupoid makes \(\mathsf{R}^{\mathrm{dR},\mathrm{st}}_\bullet(X/S)\) a simplicial stable object. The
diagonal of each de Rham relation degree induces its canonical cocommutative comultiplication.
\end{construction}

\begin{theorem}[Internal-function presentation]
\label{thm:internal-function-cochains-expanded}
For every \(n\geq0\), there is a canonical multiplicative equivalence
\[
  C^n_{\mathrm{dR}}(X/S)
  \xrightarrow{\simeq}
  \HHom_{\mathcal D_{\Sp/S}}
  \bigl(\mathsf{R}^{\mathrm{dR},\mathrm{st}}_n(X/S),\mathcal O_S^{\mathrm{add}}\bigr).
\]
The equivalences assemble into an equivalence of cosimplicial commutative algebras.
Multiplication on the function object is convolution against the diagonal of
\(\mathsf{R}^{\mathrm{dR},\mathrm{st}}_n(X/S)\).
\end{theorem}

\begin{proof}
\leavevmode
\begin{enumerate}[label=\textup{(\arabic*)}]
    \item \textbf{Apply the cited proposition.}
    Proposition~\ref{prop:multiplicative-function-spectrum-identity}, applied to
    \(a_n\) and \(A=\mathcal O_S^{\mathrm{add}}\), gives the displayed equivalence
    directly in \(\CAlg(\mathcal D_{\Sp/S})\).
    \item \textbf{Establish multiplicative compatibility.}
    The convolution multiplication is induced
    by the diagonal of \(\mathcal R^{\mathrm{dR}}_{X/S,n}\).
    \item \textbf{Apply naturality.}
    Naturality under simplicial operators follows from Beck--Chevalley, functoriality
    of the diagonals, and the mate identities.
    \item \textbf{Identify the stated objects.}
    These compatibilities identify the
    degreewise equivalences as an equivalence of cosimplicial commutative algebras.
\end{enumerate}
\end{proof}

\begin{remark}[Cochains as stable scalar functions on the de Rham relation]
Under the internal-function presentation, a degree-\(n\) de Rham cochain is a
stable scalar function on an \(n\)-simplex of the de Rham relation. The
cofaces and codegeneracies are therefore restriction of scalar functions along the
structural maps of that relation groupoid. The totalization filtration of the resulting descent object supplies the de Rham
differential in the next section.
\end{remark}

\begin{corollary}[Mapping-spectrum formula]
For every \(n\),
\[
  \Gamma^{\mathrm{st}}
  \bigl(S,C^n_{\mathrm{dR}}(X/S)\bigr)
  \simeq
  \map_{\mathcal D_{\Sp/S}}
  \bigl(\mathsf{R}^{\mathrm{dR},\mathrm{st}}_n(X/S),\mathcal O_S^{\mathrm{add}}\bigr)
\]
and equivalently
\[
  \Gamma^{\mathrm{st}}
  \bigl(S,C^n_{\mathrm{dR}}(X/S)\bigr)
  \simeq
  \Gamma^{\mathrm{st}}
  \bigl(\mathcal R^{\mathrm{dR}}_{X/S,n},a_n^*\mathcal O_S^{\mathrm{add}}\bigr).
\]
\end{corollary}

\begin{proposition}[Represented function-spectrum formula]
\label{prop:represented-function-spectrum-detailed}
If \(\mathcal R^{\mathrm{dR}}_{X/S,n}\) is represented by a spectral scheme \(Y_n\), then
\[
  \Gamma^{\mathrm{st}}
  \bigl(S,C^n_{\mathrm{dR}}(X/S)\bigr)
  \simeq
  \Gamma(Y_n,\mathcal O_{Y_n})
\]
canonically as \(E_\infty\)-ring spectra, compatibly with every represented face and
degeneracy map.
\end{proposition}

\begin{proof}
\leavevmode
\begin{enumerate}[label=\textup{(\arabic*)}]
    \item \textbf{Combine the preceding comparisons.}
    Combine the mapping-spectrum formula with
    Proposition~\ref{prop:stable-scalar-sections-represented}.
\end{enumerate}
\end{proof}

\subsection{Effective descent and the \v{C}ech--Amitsur comparison}
\label{sec:infinitesimal-descent-monad}

The cosimplicial object has been defined geometrically. To prove that it
really computes descent, we now use the effective epimorphism
\(e:X\to Q_{X/S}\). Its pullback--dependent-product adjunction produces both
the descent comonad and the Amitsur monad; Beck--Chevalley identifies their
iterates with the \v{C}ech relation degrees.

Put \(Q:=Q_{X/S}\), \(e:=e_{X/S}:X\to Q\), and let
\(p_Q:Q\to S\). The effective epimorphism \(e\) determines the adjunction
\[
  e^*:\mathcal D_{\Sp/Q}
  \rightleftarrows
  \mathcal D_{\Sp/X}:\Pi_e.
\]
We display the adjunction together with its induced monad and comonad as
\[
\begin{tikzcd}[column sep=huge,row sep=large]
  \mathcal D_{\Sp/Q}
    \arrow[r,shift left=1.2ex,"e^*"]
    \arrow[loop left,"\mathbb T^{\mathrm{dR}}_{X/S}=\Pi_e e^*"']
  &
  \mathcal D_{\Sp/X}
    \arrow[l,shift left=1.2ex,"\Pi_e"]
    \arrow[loop right, swap, "\mathbb G^{\mathrm{dR}}_{X/S}=e^*\Pi_e"].
\end{tikzcd}
\]

Two canonical endofunctors arise from this single adjunction. The monad
\(\Pi_e e^*\) on \(\mathcal D_{\Sp/Q}\) produces the Amitsur cosimplicial object used
for scalar functions, while the comonad \(e^*\Pi_e\) on \(\mathcal D_{\Sp/X}\)
encodes the actual effective-descent coalgebra structure.

\begin{definition}[Stable de Rham Amitsur monad and descent comonad]
\label{def:de-Rham-amitsur-monad}
Define
\[
  \mathbb T^{\mathrm{dR}}_{X/S}
  :=
  \Pi_e e^*:
  \mathcal D_{\Sp/Q}\longrightarrow\mathcal D_{\Sp/Q}
\]
and
\[
  \mathbb G^{\mathrm{dR}}_{X/S}
  :=
  e^*\Pi_e:
  \mathcal D_{\Sp/X}\longrightarrow\mathcal D_{\Sp/X}.
\]
The adjunction \(e^*\dashv\Pi_e\) supplies the monad structure on
\(\mathbb T^{\mathrm{dR}}_{X/S}\) and the comonad structure on
\(\mathbb G^{\mathrm{dR}}_{X/S}\). Since \(e^*\) is strong symmetric
monoidal, \(\Pi_e\) is canonically lax symmetric monoidal, so the Amitsur monad
acts functorially on commutative algebra objects. The adjunction
\(e^*\dashv\Pi_e\) supplies both structures functorially.
\end{definition}

\begin{theorem}[Comonadic effective descent]
\label{thm:comonadic-effective-descent-de-Rham}
Pullback along \(e\) is comonadic. More precisely, there is a canonical
equivalence
\[
  \mathcal D_{\Sp/Q}
  \xrightarrow{\simeq}
  \operatorname{Coalg}_{\mathbb G^{\mathrm{dR}}_{X/S}}(\mathcal D_{\Sp/X}),
\]
under which an object over \(Q\) is recovered from its pullback to \(X\) together
with the descent coalgebra structure induced by the \v{C}ech relation.
\end{theorem}

\begin{proof}
\leavevmode
\begin{enumerate}[label=\textup{(\arabic*)}]
    \item \textbf{Invoke the cited result.}
    By Proposition~\ref{thm:effective-descent-stable-coefficients-expanded},
    \[
    \mathcal D_{\Sp/Q}
    \xrightarrow{\simeq}
    \operatorname*{Tot}_{[n]\in\Delta}\mathcal D_{\Sp/\check C_n(e)}.
    \]
    \item \textbf{Treat the stated case.}
    If a morphism becomes an equivalence after pullback along \(e\), then every
    further pullback to every \v{C}ech degree is an equivalence.
    \item \textbf{Pass to totalization.}
    The totalization
    equivalence therefore makes the original morphism an equivalence.
    \item \textbf{Conclude the stated claim.}
    Hence
    \(e^*\) is conservative.
    \item \textbf{Establish the next claim.}
    The functor \(e^*\) has the left adjoint \(\Sigma_e\), so it preserves all
    limits.
    \item \textbf{Establish the next claim.}
    In particular it preserves the totalizations required by the dual
    Barr--Beck theorem.
    \item \textbf{Use the stated construction.}
    Applying the comonadic form of Barr--Beck to
    \(e^*\dashv\Pi_e\) gives the displayed equivalence; compare
    \cite[Theorem~4.7.3.5]{LurieHigherAlgebra}.
\end{enumerate}
\end{proof}

\begin{remark}[Logical reading of effective descent]
The equivalence
\[
  \mathcal D_{\Sp/Q}
  \simeq
  \operatorname{Coalg}_{\mathbb G^{\mathrm{dR}}_{X/S}}(\mathcal D_{\Sp/X})
\]
identifies a stable predicate on the effective quotient with a stable predicate on
\(X\) equipped with the coherent transport data required along the de Rham
relation. The associated Amitsur object is the corresponding cosimplicial form of that
transport data, generated by the coalgebra structure of the adjunction.
\end{remark}

\begin{theorem}[\v{C}ech--Amitsur comparison]
\label{thm:amitsur-identification}
Let \(e_n:\mathcal R^{\mathrm{dR}}_{X/S,n}\to Q\) be the structural map of the
\v{C}ech nerve of \(e\). For every \(A\in\mathcal D_{\Sp/Q}\), there are canonical
equivalences
\[
  (\mathbb T^{\mathrm{dR}}_{X/S})^{n+1}(A)
  \xrightarrow{\simeq}
  \Pi_{e_n}e_n^*A
\]
which assemble into an equivalence between the Amitsur cosimplicial object of the
monad \(\mathbb T^{\mathrm{dR}}_{X/S}\) and the stable \v{C}ech
cosimplicial object.
Degreewise, the comparison has the form
\[
\begin{tikzcd}[column sep=huge]
  (\mathbb T^{\mathrm{dR}}_{X/S})^{n+1}(A)
    \arrow[r,"\simeq"]
  &
  \Pi_{e_n}e_n^*A.
\end{tikzcd}
\]
  If \(A\) is a commutative algebra, this is an equivalence of
cosimplicial commutative algebras. The monad unit gives the cofaces and the monad
multiplication gives the codegeneracies.
\end{theorem}

\begin{proof}
\leavevmode
\begin{enumerate}[label=\textup{(\arabic*)}]
    \item \textbf{Apply the cited theorem.}
    Theorem~\ref{thm:relation-effective-presentation} identifies
    \(\mathcal R^{\mathrm{dR}}_{X/S,\bullet}\) with the \v{C}ech nerve of
    \(e\).
    \item \textbf{Establish the next claim.}
    For the first nontrivial iterate, right Beck--Chevalley for
    \[
    \begin{tikzcd}
    X\times_QX \ar[r,"\operatorname{pr}_2"] \ar[d,"\operatorname{pr}_1"']
    & X \ar[d,"e"]\\
    X \ar[r,"e"'] & Q
    \end{tikzcd}
    \]
    gives
    \[
    e^*\Pi_e
    \simeq
    \Pi_{\operatorname{pr}_1}\operatorname{pr}_2^*.
    \]
    \item \textbf{Use the stated construction.}
    Substitution into \((\Pi_e e^*)^2\) and composition of dependent products
    gives
    \[
    (\mathbb T^{\mathrm{dR}}_{X/S})^2A
    \simeq
    \Pi_{e_1}e_1^*A.
    \]
    \item \textbf{Establish the next claim.}
    Iterating the same Cartesian comparison adds one \v{C}ech factor at each
    stage.
    \item \textbf{Apply naturality.}
    Naturality of adjunction units, counits, and their mates identifies the
    cofaces and codegeneracies and supplies all cosimplicial coherences.
\end{enumerate}
\end{proof}

\begin{corollary}[\v{C}ech--Amitsur totalization]
\label{cor:cech-amitsur-totalization}
For every \(A\in\mathcal D_{\Sp/Q}\), the canonical augmentation is a totalization
equivalence
\[
  A
  \xrightarrow{\simeq}
  \operatorname*{Tot}
  \operatorname{Am}^\bullet_{\mathbb T^{\mathrm{dR}}_{X/S}}(A).
\]
\end{corollary}

\begin{proof}
\leavevmode
\begin{enumerate}[label=\textup{(\arabic*)}]
    \item \textbf{Invoke the cited result.}
    By Theorem~\ref{thm:amitsur-identification}, the Amitsur object is the stable
    \v{C}ech object
    \[
    [n]\longmapsto\Pi_{e_n}e_n^*A.
    \]
    \item \textbf{Apply the cited proposition.}
    Proposition~\ref{thm:effective-descent-stable-coefficients-expanded} identifies
    its totalization with \(A\).
    \item \textbf{Pass to totalization.}
    This is the \v{C}ech totalization statement; the
    categorical descent theorem is the comonadic equivalence of
    Theorem~\ref{thm:comonadic-effective-descent-de-Rham}.
\end{enumerate}
\end{proof}

\begin{theorem}[Amitsur presentation of geometric de Rham cochains]
\label{thm:amitsur-presentation-geometric-cochains}
By Proposition~\ref{prop:additive-scalar-base-change},
\[
  p_Q^*\mathcal O_S^{\mathrm{add}}
  \xrightarrow{\simeq}
  \mathcal O_Q^{\mathrm{add}}
\]
canonically in \(\CAlg(\mathcal D_{\Sp/Q})\). There is a canonical equivalence of
cosimplicial commutative algebras
\[
  C^\bullet_{\mathrm{dR}}(X/S)
  \xrightarrow{\simeq}
  \Pi_{p_Q}
  \operatorname{Am}^\bullet_{\mathbb T^{\mathrm{dR}}_{X/S}}
  \bigl(\mathcal O_Q^{\mathrm{add}}\bigr).
\]
\end{theorem}

\begin{proof}
\leavevmode
\begin{enumerate}[label=\textup{(\arabic*)}]
    \item \textbf{Apply the cited theorem.}
    Theorem~\ref{thm:amitsur-identification} gives
    \[
    (\mathbb T^{\mathrm{dR}}_{X/S})^{n+1}
    \bigl(\mathcal O_Q^{\mathrm{add}}\bigr)
    \simeq
    \Pi_{e_n}e_n^*p_Q^*\mathcal O_S^{\mathrm{add}}.
    \]
    \item \textbf{Apply the indicated result.}
    Apply \(\Pi_{p_Q}\) and use \(a_n=p_Qe_n\).
    \item \textbf{Use pullback stability.}
    The strong symmetric monoidal
    structure on pullback and the canonical lax symmetric monoidal structure on
    dependent product make the comparison multiplicative in every degree, while the
    mate identities make it cosimplicial.
\end{enumerate}
\end{proof}

\subsection{Effective and full de Rham algebras}
\label{sec:full-effective-de-Rham}

Effective descent determines functions on \(Q_{X/S}\), not automatically on
the whole reflected object. We therefore name the two function algebras
separately before comparing them. This keeps the effectivity distinction
visible at the coefficient level.

Let
\[
  p_{\mathrm{dR}}:X_{\mathrm{dR}/S}\to S,
  \qquad
  p_Q:Q_{X/S}\to S.
\]

\begin{definition}[Full and effective de Rham algebras]
Define
\[
  \mathrm{DR}^{\mathrm{full}}_{X/S}
  :=
  \Pi_{p_{\mathrm{dR}}}p_{\mathrm{dR}}^*\mathcal O_S^{\mathrm{add}},
\]
and define the effective de Rham algebra by
\[
  \mathrm{DR}_{X/S}
  :=
  \operatorname*{Tot}C^\bullet_{\mathrm{dR}}(X/S).
\]
Both are commutative algebra objects of \(\mathcal D_{\Sp/S}\).
\end{definition}

\begin{theorem}[Effective de Rham algebra]
\label{thm:effective-DR-functions}
There are canonical multiplicative equivalences
\[
  \mathrm{DR}_{X/S}
  \simeq
  \Pi_{p_Q} p_Q^*\mathcal O_S^{\mathrm{add}}
\]
and
\[
  \mathrm{DR}_{X/S}
  \simeq
  \HHom_{\mathcal D_{\Sp/S}}
  \bigl(\Sigma^\infty_{+,S}Q_{X/S},\mathcal O_S^{\mathrm{add}}\bigr).
\]
Likewise
\[
  \mathrm{DR}^{\mathrm{full}}_{X/S}
  \simeq
  \HHom_{\mathcal D_{\Sp/S}}
  \bigl(\Sigma^\infty_{+,S}X_{\mathrm{dR}/S},\mathcal O_S^{\mathrm{add}}\bigr).
\]
The monomorphism \(m_{X/S}:Q_{X/S}\to X_{\mathrm{dR}/S}\) induces a
canonical morphism of commutative algebras
\[
  \rho_{X/S}:
  \mathrm{DR}^{\mathrm{full}}_{X/S}
  \longrightarrow
  \mathrm{DR}_{X/S}.
\]
It is induced contravariantly by the geometric inclusion
\[
\begin{tikzcd}[column sep=huge,row sep=large]
  Q_{X/S} \arrow[r,hook,"m_{X/S}"]
  & X_{\mathrm{dR}/S}
  \\
  \mathrm{DR}_{X/S}
  & \mathrm{DR}^{\mathrm{full}}_{X/S}
      \arrow[l,"\rho_{X/S}"'].
\end{tikzcd}
\]
The morphism \(\rho_{X/S}\) is an equivalence whenever \(m_{X/S}\) is an
equivalence; in particular this holds on the formally smooth effectivity domain.
\end{theorem}

\begin{proof}
\leavevmode
\begin{enumerate}[label=\textup{(\arabic*)}]
    \item \textbf{Construct the indicated map.}
    The de Rham relation groupoid is the \v{C}ech nerve of the effective epimorphism
    \(e_{X/S}:X\to Q_{X/S}\).
    \item \textbf{Use the stated construction.}
    Stable effective descent applied to
    \(p_Q^*\mathcal O_S^{\mathrm{add}}\) gives
    \[
    p_Q^*\mathcal O_S^{\mathrm{add}}
    \xrightarrow{\simeq}
    \operatorname*{Tot}_n
    \Pi_{e_n}e_n^*p_Q^*\mathcal O_S^{\mathrm{add}}.
    \]
    \item \textbf{Apply the indicated result.}
    Apply \(\Pi_{p_Q}\) and use composition of dependent products to obtain the first
    formula.
    \item \textbf{Apply the indicated result.}
    Apply Proposition~\ref{prop:function-spectrum-identity} directly to
    \(p_Q:Q_{X/S}\to S\) to obtain
    \[
    \Pi_{p_Q} p_Q^*\mathcal O_S^{\mathrm{add}}
    \simeq
    \HHom_{\mathcal D_{\Sp/S}}
    \bigl(\Sigma^\infty_{+,S}Q_{X/S},\mathcal O_S^{\mathrm{add}}\bigr).
    \]
    \item \textbf{Establish the next claim.}
    The full formula is the same proposition applied to \(p_{\mathrm{dR}}\).
    \item \textbf{Use the stated construction.}
    Restriction along \(m_{X/S}\) gives \(\rho\), and its multiplicativity follows from
    naturality of the diagonal comultiplications.
    \item \textbf{Treat the stated case.}
    If \(m_{X/S}\) is an equivalence, so
    is the induced restriction map.
\end{enumerate}
\end{proof}


\begin{proposition}[Augmentation to degree zero]
\label{prop:de-Rham-augmentation-expanded}
If \(p:X\to S\), evaluation of the totalization cone in cosimplicial degree zero
defines a canonical morphism
\[
  \varepsilon_{\mathrm{dR}}:
  \mathrm{DR}_{X/S}
  \longrightarrow
  \Pi_pp^*\mathcal O_S^{\mathrm{add}}.
\]
Under the effective-quotient description it is \(\Pi_{p_Q}\) applied to the monad unit
\[
  p_Q^*\mathcal O_S^{\mathrm{add}}
  \longrightarrow
  \Pi_e e^*p_Q^*\mathcal O_S^{\mathrm{add}}.
\]
\end{proposition}


The effective descent problem has now been solved at the level of scalar
functions. The total object \(\mathrm{DR}_{X/S}\) is the algebra of stable
scalar functions on \(Q_{X/S}\), while
\(\mathrm{DR}^{\mathrm{full}}_{X/S}\) records functions on the full relative
de Rham reflection. What remains is to resolve the totalization itself.
Its partial totalizations and normalized layers will expose the successive
coherence conditions of infinitesimal descent and thereby produce forms and
the de Rham differential.
\section{Totalization: forms, differential, products, and closedness}
\label{sec:de-Rham-coherence-forms}
\label{sec:de-Rham-functoriality}

The previous section identified \(\mathrm{DR}_{X/S}\) as the totalization of a
cosimplicial stable scalar algebra. The remaining question is how the
coherent descent conditions contained in that totalization should be read
degree by degree. Stable Dold--Kan supplies the answer: normalization isolates
the new datum appearing in each cosimplicial degree, while the
partial-totalization tower records the successive extension conditions that
assemble those data into a coherent descent object.

The two descriptions are the layer and tower views of the same object:
\[
\begin{tikzcd}[column sep=huge,row sep=large]
  C^\bullet_{\mathrm{dR}}(X/S)
      \arrow[r,"{\operatorname*{Tot}}"]
      \arrow[d,"{N^q}"']
  &
  \mathrm{DR}_{X/S}
      \arrow[d,"{\operatorname{Gr}^q_{\mathrm{Tot}}[-]}"]
  \\
  \Omega^q_{\mathrm{dR}}(X/S)
      \arrow[r,"{\simeq}"']
  &
  \operatorname{Gr}^q_{\mathrm{Tot}}\mathrm{DR}_{X/S}[q].
\end{tikzcd}
\]
The connecting maps between adjacent layers give the homotopy-coherent de
Rham differential. Because the original cochains form a cosimplicial
commutative algebra, the same construction also supplies products and the
Leibniz rule. Keeping the full filtration stage rather than only its
associated graded records closedness, exact-form lifts, and the higher
coherences carried by a closed form.
\subsection{Totalization layers and intrinsic de Rham forms}
\label{sec:complete-totalization}


For each degree \(q\), the partial totalization through degree \(q-1\)
records all descent conditions already imposed below that degree. Its fibre
inside the complete totalization measures the remaining coherence, while the
normalized degree-\(q\) term is the next associated-graded layer. We now
make this tower explicit.


For \(q\geq0\), let \(\Delta_{<q}\subset\Delta\) be the full subcategory on
\([0],\ldots,[q-1]\), and define
\[
  \mathrm{DR}^{<q}_{X/S}
  :=
  \operatorname*{lim}_{[n]\in\Delta_{<q}}
  C^n_{\mathrm{dR}}(X/S).
\]
For \(q=0\), this is the zero object in the underlying stable category. Define the
decreasing totalization filtration
\[
  F^q_{\mathrm{Tot}}\mathrm{DR}_{X/S}
  :=
  \operatorname{fib}
  \bigl(
    \mathrm{DR}_{X/S}
    \longrightarrow
    \mathrm{DR}^{<q}_{X/S}
  \bigr).
\]
The first stages form the inverse tower
\[
\begin{tikzcd}[column sep=large]
  \mathrm{DR}_{X/S}
    \arrow[r]
  & \cdots \arrow[r]
  & \mathrm{DR}^{<q+1}_{X/S}
    \arrow[r]
  & \mathrm{DR}^{<q}_{X/S}
    \arrow[r]
  & \cdots \arrow[r]
  & \mathrm{DR}^{<1}_{X/S},
\end{tikzcd}
\]
whose fibre tower is \(F^\bullet_{\mathrm{Tot}}\mathrm{DR}_{X/S}\).


\begin{definition}[Intrinsic de Rham forms]
For \(q\geq0\), define \(\Omega^q_{\mathrm{dR}}(X/S)\) to be the stable
Dold--Kan normalization of the cosimplicial object
\(C^\bullet_{\mathrm{dR}}(X/S)\) in degree \(q\). Equivalently,
\(\Omega^q_{\mathrm{dR}}\) is the total fibre of the finite codegeneracy matching cube
formed from all proper surjections \([q]\twoheadrightarrow[r]\).
\end{definition}

\begin{proposition}[Explicit normalization]
\label{prop:explicit-normalization-expanded}
There is a canonical equivalence
\[
  \Omega^q_{\mathrm{dR}}(X/S)
  \simeq
  \operatorname{tfib}
  \Bigl(
  C^q_{\mathrm{dR}}(X/S)
  \longrightarrow
  \{C^r_{\mathrm{dR}}(X/S)\}_{[q]\twoheadrightarrow[r],\,r<q}
  \Bigr).
\]
In particular,
\[
  \Omega^0_{\mathrm{dR}}(X/S)=C^0_{\mathrm{dR}}(X/S)
\]
and
\[
  \Omega^1_{\mathrm{dR}}(X/S)
  \simeq
  \operatorname{fib}
  \bigl(C^1_{\mathrm{dR}}(X/S)\to C^0_{\mathrm{dR}}(X/S)\bigr),
\]
where the latter map is induced by the diagonal \(X\to \mathcal R^{\mathrm{dR}}_{X/S,1}\).
The degree-one normalization is displayed by the fibre square
\[
\begin{tikzcd}[column sep=huge,row sep=large]
  \Omega^1_{\mathrm{dR}}(X/S) \arrow[r] \arrow[d]
  & C^1_{\mathrm{dR}}(X/S) \arrow[d]
  \\
  0 \arrow[r]
  & C^0_{\mathrm{dR}}(X/S).
\end{tikzcd}
\]
\end{proposition}

\begin{proof}
\leavevmode
\begin{enumerate}[label=\textup{(\arabic*)}]
    \item \textbf{Establish the next claim.}
    This is the matching-object description of stable Dold--Kan normalization; see
    \cite{LurieHigherAlgebra}.
    \item \textbf{Identify the relevant fibre.}
    The matching diagram is finite, so its fibre is equivalently the
    total fibre of the codegeneracy cube.
\end{enumerate}
\end{proof}

\begin{theorem}[Stable Dold--Kan layers and completeness]
\label{thm:complete-filtration}
For every \(q\geq0\), there is a canonical fibre/cofiber sequence
\[
  F^{q+1}_{\mathrm{Tot}}\mathrm{DR}_{X/S}
  \longrightarrow
  F^q_{\mathrm{Tot}}\mathrm{DR}_{X/S}
  \longrightarrow
  \Omega^q_{\mathrm{dR}}(X/S)[-q].
\]
Consequently
\[
  \operatorname{Gr}^q_{\mathrm{Tot}}\mathrm{DR}_{X/S}
  \simeq
  \Omega^q_{\mathrm{dR}}(X/S)[-q].
\]
The filtration is already complete:
\[
  \varprojlim\nolimits_qF^q_{\mathrm{Tot}}\mathrm{DR}_{X/S}
  \simeq0,
\]
and
\[
  \mathrm{DR}_{X/S}
  \xrightarrow{\simeq}
  \varprojlim\nolimits_q\mathrm{DR}^{<q}_{X/S}.
\]
\end{theorem}

\begin{proof}
\leavevmode
\begin{enumerate}[label=\textup{(\arabic*)}]
    \item \textbf{Invoke the cited result.}
    By stable Dold--Kan, equivalently the Bousfield--Kan partial-totalization tower of
    a cosimplicial object in a stable infinity-category, the successive fibres of the
    partial-totalization tower are the shifted normalized terms; see
    \cite{LurieHigherAlgebra}.
    \item \textbf{Pass to totalization.}
    Passing to the fibre tower from the complete totalization gives
    the displayed sequence.
    \item \textbf{Invoke the cited result.}
    By definition, the totalization is the limit of its finite partial totalizations.
    \item \textbf{Take the indicated specialization.}
    Taking the inverse limit of the defining fibre sequences gives
    \[
    \varprojlim\nolimits_qF^q_{\mathrm{Tot}}\mathrm{DR}_{X/S}
    \simeq
    \operatorname{fib}
    \left(
    \mathrm{DR}_{X/S}
    \longrightarrow
    \varprojlim\nolimits_q\mathrm{DR}^{<q}_{X/S}
    \right)
    \simeq0.
    \]
    \item \textbf{Conclude the stated claim.}
    Thus the totalization filtration is complete.
\end{enumerate}
\end{proof}

Extend the filtration by the constant value \(\mathrm{DR}_{X/S}\) in
nonpositive filtration degrees. Since \(\mathcal D_{\Sp/S}\) is presentable and
stable, it admits sequential limits. The complete-filtered/coherent-cochain
correspondence of \cite[Theorem~4.7 and Remark~4.9]{AriottaCoherentCochains}
therefore applies directly to the preceding complete filtration.

\begin{lemma}[Stable Dold--Kan identification of totalization boundaries]
\label{lem:totalization-normalization-sign}
Let \(C^\bullet\) be a cosimplicial object in a stable infinity-category, and use
the cohomological normalization
\[
  d_N^q
  =
  \sum_{i=0}^{q+1}(-1)^id^i
  :
  N^qC\longrightarrow N^{q+1}C.
\]
Put
\[
  \operatorname{Tot}^{<q}C
  :=
  \operatorname*{lim}_{[n]\in\Delta_{<q}}C^n,
  \qquad
  F^q
  :=
  \operatorname{fib}
  \bigl(
    \operatorname{Tot}C
    \longrightarrow
    \operatorname{Tot}^{<q}C
  \bigr).
\]
Under the canonical stable Dold--Kan layer equivalences
\[
  \operatorname{cofib}(F^{q+1}\to F^q)
  \xrightarrow{\simeq}
  N^qC[-q],
\]
the shifted connecting morphism
\[
  N^qC
  \simeq
  \operatorname{Gr}^q_F[q]
  \longrightarrow
  \operatorname{Gr}^{q+1}_F[q+1]
  \simeq
  N^{q+1}C
\]
is represented in the homotopy category by the normalized alternating coface
differential \(d_N^q\). The identification is natural in \(C^\bullet\).
\end{lemma}

\begin{proof}
\leavevmode
\begin{enumerate}[label=\textup{(\arabic*)}]
    \item \textbf{Identify the stated objects.}
    Stable Dold--Kan identifies a cosimplicial object in a stable
    \(\infty\)-category, equivalently a simplicial object in the opposite stable
    \(\infty\)-category, with the filtration obtained from its skeletal tower.
    \item \textbf{Identify the stated objects.}
    Dualizing the skeletal filtration identifies the successive fibres of the
    partial-totalization tower with \(N^qC[-q]\).
    \item \textbf{Establish the next claim.}
    Under the comparison with the normalized complex in the homotopy category, the
    boundary is the cohomological differential
    \[
    \sum_{i=0}^{q+1}(-1)^id^i.
    \]
    \item \textbf{Establish the next claim.}
    This is the compatibility of the stable Dold--Kan filtration differential with
    normalization; see \cite[Section~1.2.4]{LurieHigherAlgebra}.
    \item \textbf{Apply naturality.}
    Naturality follows
    from functoriality of the Dold--Kan construction.
\end{enumerate}
\end{proof}

\begin{theorem}[Homotopy-coherent de Rham differential]
\label{thm:complete-filtrations-coherent-cochains}
The complete totalization filtration canonically determines a coherent nonnegative
cochain object
\[
  \Omega^0_{\mathrm{dR}}(X/S)
  \xrightarrow{d_{\mathrm{dR}}}
  \Omega^1_{\mathrm{dR}}(X/S)
  \xrightarrow{d_{\mathrm{dR}}}
  \Omega^2_{\mathrm{dR}}(X/S)
  \longrightarrow\cdots.
\]
Its first two stages may be read together with the totalization layers as
\[
\begin{tikzcd}[column sep=large,row sep=large]
  F^0_{\mathrm{Tot}} \arrow[d]
  & F^1_{\mathrm{Tot}} \arrow[l] \arrow[d]
  & F^2_{\mathrm{Tot}} \arrow[l] \arrow[d]
  \\
  \Omega^0_{\mathrm{dR}}
    \arrow[r,"d_{\mathrm{dR}}"']
  & \Omega^1_{\mathrm{dR}}
    \arrow[r,"d_{\mathrm{dR}}"']
  & \Omega^2_{\mathrm{dR}},
\end{tikzcd}
\]
where the vertical arrows denote the associated layer maps with the shifts
specified in Theorem~\ref{thm:complete-filtration}.

The construction retains specified nullhomotopies
\[
  d^{q+1}_{\mathrm{dR}}d^q_{\mathrm{dR}}\simeq0
\]
and the full higher coherence supplied by the complete filtration. With the
cohomological normalization convention fixed above, the image of
\(d^q_{\mathrm{dR}}\) in the homotopy category is
\[
  d^q_{\mathrm{dR}}
  =
  \sum_{i=0}^{q+1}(-1)^id^i
  :
  \Omega^q_{\mathrm{dR}}(X/S)
  \longrightarrow
  \Omega^{q+1}_{\mathrm{dR}}(X/S).
\]
\end{theorem}

\begin{proof}
\leavevmode
\begin{enumerate}[label=\textup{(\arabic*)}]
    \item \textbf{Establish the next claim.}
    Extend the complete decreasing totalization filtration by the constant value
    \(\mathrm{DR}_{X/S}\) in nonpositive filtration degrees.
    \item \textbf{Establish the next claim.}
    The
    complete-filtered/coherent-cochain correspondence of
    \cite[Theorem~4.7 and Remark~4.9]{AriottaCoherentCochains} sends this filtration to a coherent nonnegative
    cochain object whose degree-\(q\) term is
    \[
    \operatorname{Gr}^q_{\mathrm{Tot}}\mathrm{DR}_{X/S}[q].
    \]
    \item \textbf{Invoke the cited result.}
    By Theorem~\ref{thm:complete-filtration},
    \[
    \operatorname{Gr}^q_{\mathrm{Tot}}\mathrm{DR}_{X/S}[q]
    \simeq
    \Omega^q_{\mathrm{dR}}(X/S).
    \]
    \item \textbf{Conclude the stated claim.}
    Thus the terms of the resulting coherent cochain object are precisely the intrinsic
    de Rham forms.
    \item \textbf{Construct the indicated map.}
    Under the complete-filtered/coherent-cochain correspondence, the first structure
    map from degree \(q\) to degree \(q+1\) is the shifted connecting morphism between
    the adjacent filtration layers.
    \item \textbf{Invoke the cited result.}
    By
    Lemma~\ref{lem:totalization-normalization-sign}, its image in the homotopy category
    is exactly
    \[
    \sum_{i=0}^{q+1}(-1)^id^i.
    \]
    \item \textbf{Construct the indicated map.}
    We denote the corresponding cochain morphism by \(d^q_{\mathrm{dR}}\).
    \item \textbf{Record the higher cochain coherence.}
    The coherent-cochain object supplied by the equivalence contains, as part of the
    same functorial construction, the specified nullhomotopies
    \[
    d^{q+1}_{\mathrm{dR}}d^q_{\mathrm{dR}}\simeq0
    \]
    and all higher compatibility among them.
\end{enumerate}
\end{proof}

\begin{proposition}[Degree-zero de Rham differential]
\label{prop:degree-zero-infinitesimal-differential}
With the ordered degree-one faces
\[
  \partial_0=\mathsf t,
  \qquad
  \partial_1=\mathsf s,
\]
the first differential
\[
  d^0_{\mathrm{dR}}:
  \Omega^0_{\mathrm{dR}}(X/S)
  \longrightarrow
  \Omega^1_{\mathrm{dR}}(X/S)
\]
is induced by the scalar-restriction difference
\[
  \operatorname{res}_{\mathsf t}-\operatorname{res}_{\mathsf s}:
  C^0_{\mathrm{dR}}(X/S)
  \longrightarrow
  C^1_{\mathrm{dR}}(X/S).
\]
\end{proposition}

\begin{proof}
\leavevmode
\begin{enumerate}[label=\textup{(\arabic*)}]
    \item \textbf{Invoke the cited result.}
    By
    Lemma~\ref{lem:totalization-normalization-sign}, the degree-zero homotopy-coherent
    differential is represented in the homotopy category by the normalized
    cohomological alternating coface sum
    \[
    d_N^0=d^0-d^1.
    \]
    \item \textbf{Establish the next claim.}
    Under the ordered degree-one face convention,
    \[
    d^0=\operatorname{res}_{\mathsf t},
    \qquad
    d^1=\operatorname{res}_{\mathsf s},
    \]
    so
    \[
    d_N^0
    =
    \operatorname{res}_{\mathsf t}-\operatorname{res}_{\mathsf s}.
    \]
    \item \textbf{Establish the next claim.}
    The degree-one codegeneracy is induced by the diagonal
    \[
    X\longrightarrow \mathcal R^{\mathrm{dR}}_{X/S,1}.
    \]
    \item \textbf{Establish the next claim.}
    Both projections \(\mathsf t\) and \(\mathsf s\) restrict along the diagonal to the identity of
    \(X\).
    \item \textbf{Establish the next claim.}
    Consequently the composite of
    \[
    \operatorname{res}_{\mathsf t}-\operatorname{res}_{\mathsf s}
    \]
    with the degree-one codegeneracy is canonically null.
    \item \textbf{Identify the relevant fibre.}
    The stable fibre universal
    property therefore factors the alternating difference through
    \[
    \Omega^1_{\mathrm{dR}}(X/S)
    \simeq
    \operatorname{fib}
    \bigl(
    C^1_{\mathrm{dR}}(X/S)
    \longrightarrow
    C^0_{\mathrm{dR}}(X/S)
    \bigr),
    \]
    and that factor is \(d^0_{\mathrm{dR}}\).
\end{enumerate}
\end{proof}

\begin{corollary}[Associated-graded detection]
A morphism between two complete totalization filtrations of the form above is an
equivalence if it is an equivalence on every normalized associated-graded term.
\end{corollary}

\begin{proof}
\leavevmode
\begin{enumerate}[label=\textup{(\arabic*)}]
    \item \textbf{Construct the indicated map.}
    The fibre of such a morphism is complete.
    \item \textbf{Treat the stated case.}
    If all associated-graded terms vanish,
    every transition in the fibre filtration is an equivalence; completeness then forces
    the constant value to vanish.
\end{enumerate}
\end{proof}

\begin{remark}[Higher-order information]
The totalization filtration is defined from the de Rham relation itself. Its
normalized terms can retain finite higher-order infinitesimal information beyond the
first structured infinitesimal neighbourhood. Section~\ref{sec:first-order-linearization}
constructs the canonical first-order comparison with the cotangent complex.
\end{remark}

\subsubsection{Skeletal presentation}
\label{subsec:skeletal-presentation}

Let \(\mathsf{R}^{\mathrm{dR},\mathrm{st}}_\bullet(X/S)\) be the simplicial
stable de Rham relation object constructed above. Define
\(\mathsf{R}^{\mathrm{dR},\mathrm{st}}_{\geq q}(X/S)\) by the cofiber sequence
\[
  |\operatorname{sk}_{q-1}\mathsf{R}^{\mathrm{dR},\mathrm{st}}_\bullet(X/S)|
  \longrightarrow
  |\mathsf{R}^{\mathrm{dR},\mathrm{st}}_\bullet(X/S)|
  \longrightarrow
  \mathsf{R}^{\mathrm{dR},\mathrm{st}}_{\geq q}(X/S).
\]

\begin{theorem}[Skeletal totalization presentation]
\label{thm:skeletal-totalization-presentation}
There are canonical equivalences
\[
  \mathrm{DR}_{X/S}
  \simeq
  \HHom
  \bigl(|\mathsf{R}^{\mathrm{dR},\mathrm{st}}_\bullet(X/S)|,\mathcal O_S^{\mathrm{add}}\bigr),
\]
\[
  \mathrm{DR}^{<q}_{X/S}
  \simeq
  \HHom
  \bigl(|\operatorname{sk}_{q-1}\mathsf{R}^{\mathrm{dR},\mathrm{st}}_\bullet(X/S)|,
        \mathcal O_S^{\mathrm{add}}\bigr),
\]
and
\[
  F^q_{\mathrm{Tot}}\mathrm{DR}_{X/S}
  \simeq
  \HHom
  \bigl(\mathsf{R}^{\mathrm{dR},\mathrm{st}}_{\geq q}(X/S),\mathcal O_S^{\mathrm{add}}\bigr).
\]
They are natural in \(X/S\) and compatible with arbitrary base change.
\end{theorem}

\begin{proof}
\leavevmode
\begin{enumerate}[label=\textup{(\arabic*)}]
    \item \textbf{Establish the next claim.}
    The degreewise internal-function theorem turns geometric realization in the
    simplicial source into totalization of the cosimplicial function object: internal
    Hom into the fixed scalar converts colimits in its first variable into limits.
    \item \textbf{Pass to totalization.}
    Restricting to the \((q-1)\)-skeleton gives the partial totalization, and applying
    internal Hom to the skeletal cofiber sequence gives the fibre sequence defining the
    filtration stage.
    \item \textbf{Apply naturality.}
    Naturality in \(X/S\) follows degreewise from naturality of the de Rham
    relation nerve and scalar restriction.
    Let \(c:S'\to S\).
    \item \textbf{Identify the stated objects.}
    Relative de Rham
    base change identifies every relation degree with its pullback to \(S'\),
    while Proposition~\ref{prop:additive-scalar-base-change} identifies the scalar
    coefficient.
    \item \textbf{Establish multiplicative compatibility.}
    The functor \(c^*\) has both adjoints and therefore preserves all
    limits and colimits, so it commutes with geometric realizations, skeleta, and
    cofibres.
    \item \textbf{Use the stated construction.}
    The internal-function terms are transported by the closed
    base-change equivalence of
    Proposition~\ref{prop:closed-base-change-internal-hom}; its convolution
    compatibility gives the multiplicative comparison.
    \item \textbf{Conclude the stated claim.}
    Thus all three displayed function-object formulas commute with arbitrary base
    change.
\end{enumerate}
\end{proof}

\subsection{Multiplication and the normalized Leibniz rule}
\label{sec:de-Rham-multiplication}

Every partial totalization and the complete totalization are limits of a cosimplicial
commutative algebra and hence carry canonical \(E_\infty\)-algebra structures.
The normalized terms carry the ordered Alexander--Whitney product used below.

\begin{lemma}[Normalized Alexander--Whitney factorization]
\label{lem:consecutive-block-factorization}
Let \(C^\bullet,D^\bullet,E^\bullet\) be cosimplicial objects in a stable symmetric
monoidal infinity-category whose tensor product is exact separately in each
variable, and let
\[
  \mu^\bullet:C^\bullet\otimes D^\bullet\to E^\bullet
\]
be a natural pairing. For \(r,s\geq0\), the consecutive-block map in degree
\(r+s\) canonically factors to an actual morphism of normalized objects
\[
  N^rC\otimes N^sD
  \longrightarrow
  N^{r+s}E.
\]
In the homotopy category it is represented by
\[
  x\boxtimes_{\mathrm{AW}}y
  =
  \mu^{r+s}
  \bigl(d^{r+s}\cdots d^{r+1}x,(d^0)^ry\bigr).
\]
\end{lemma}

\begin{proof}
\leavevmode
\begin{enumerate}[label=\textup{(\arabic*)}]
    \item \textbf{Identify the relevant fibre.}
    For \(k\geq0\), let \(K_k(C)\) be the finite codegeneracy cube whose total fibre
    is \(N^kC\).
    \item \textbf{Establish the next claim.}
    The decomposition of the degeneracy directions
    \[
    \{0,\ldots,r+s-1\}
    =
    \{0,\ldots,r-1\}
    \amalg
    \{r,\ldots,r+s-1\}
    \]
    produces, by the front and back consecutive-block inclusions, a morphism
    \[
    K_r(C)\boxtimes K_s(D)
    \longrightarrow
    K_{r+s}(E).
    \]
    \item \textbf{Apply naturality.}
    Naturality of \(\mu\) supplies all cube compatibilities.
    \item \textbf{Use the stated hypothesis.}
    Since tensor product is
    exact separately in each variable, taking the iterated total fibre of the source
    identifies it canonically with
    \[
    N^rC\otimes N^sD.
    \]
    \item \textbf{Take the indicated specialization.}
    Taking total fibres therefore gives the required normalized map.
    \item \textbf{Establish the next claim.}
    The initial
    vertex is the displayed consecutive-block formula.
\end{enumerate}
\end{proof}

\begin{theorem}[Ordered product and Leibniz rule]
\label{thm:aw-product-expanded}
For \(r,s\geq0\), the multiplication of \(C^\bullet_{\mathrm{dR}}(X/S)\) induces canonical maps
\[
  \smile_{\mathrm{AW}}:
  \Omega^r_{\mathrm{dR}}(X/S)\otimes\Omega^s_{\mathrm{dR}}(X/S)
  \longrightarrow
  \Omega^{r+s}_{\mathrm{dR}}(X/S).
\]
Their images in the homotopy category are associative and unital, and
\[
  d_{\mathrm{dR}}(x\smile_{\mathrm{AW}}y)
  =
  d_{\mathrm{dR}}x\smile_{\mathrm{AW}}y
  +(-1)^r x\smile_{\mathrm{AW}}d_{\mathrm{dR}}y
\]
there.
\end{theorem}

\begin{proof}
\leavevmode
\begin{enumerate}[label=\textup{(\arabic*)}]
    \item \textbf{Apply the indicated result.}
    Apply Lemma~\ref{lem:consecutive-block-factorization} to the multiplication pairing.
    \item \textbf{Establish multiplicative compatibility.}
    For three inputs, both parenthesizations are induced by multiplication along the same
    three consecutive vertex blocks, so associativity follows from the cosimplicial
    identities and associativity of multiplication.
    \item \textbf{Establish the next claim.}
    The unit is the degree-zero unit.
    \item \textbf{Establish the next claim.}
    For the differential, expand the alternating coface sum.
    \item \textbf{Establish the next claim.}
    Cofaces before the break
    produce \(d x\smile y\), those after the break produce
    \((-1)^rx\smile dy\), and the two contributions at the break agree by the ordinal
    identities and occur with opposite signs.
    \item \textbf{Establish the next claim.}
    They cancel, giving the Leibniz formula.
\end{enumerate}
\end{proof}


\begin{theorem}[Commutative algebra structures on totalizations]
\label{thm:einfty-totalizations}
Every \(\mathrm{DR}^{<q}_{X/S}\) and \(\mathrm{DR}_{X/S}\) carries its canonical
commutative algebra structure, the transition maps form a tower in
\(\CAlg(\mathcal D_{\Sp/S})\), and
\[
  \mathrm{DR}_{X/S}
  \xrightarrow{\simeq}
  \varprojlim\nolimits_q\mathrm{DR}^{<q}_{X/S}
\]
is an equivalence of commutative algebras.
\end{theorem}

\begin{proof}
\leavevmode
\begin{enumerate}[label=\textup{(\arabic*)}]
    \item \textbf{Establish multiplicative compatibility.}
    Limits of commutative algebra objects are created in the underlying stable category; see \cite{LurieHigherAlgebra}.
    \item \textbf{Apply the indicated result.}
    Apply this to the finite and complete totalizations of
    \(C^\bullet_{\mathrm{dR}}(X/S)\).
\end{enumerate}
\end{proof}

\subsection{External products}
\label{subsec:external-products-detail}

\begin{proposition}[Products of relative de Rham objects and effective images]
\label{prop:de-Rham-product-compatibility}
For \(X,Y\to S\), there are canonical equivalences
\[
  (X\times_SY)_{\mathrm{dR}/S}
  \xrightarrow{\simeq}
  X_{\mathrm{dR}/S}\times_SY_{\mathrm{dR}/S},
\]
\[
  Q_{X\times_SY/S}
  \xrightarrow{\simeq}
  Q_{X/S}\times_SQ_{Y/S},
\]
and, for every \(n\geq0\),
\[
  \mathcal R^{\mathrm{dR}}_{X\times_SY/S,n}
  \xrightarrow{\simeq}
  \mathcal R^{\mathrm{dR}}_{X/S,n}
  \times_S
  \mathcal R^{\mathrm{dR}}_{Y/S,n}.
\]
These equivalences are natural in both variables and compatible with iterated
products and arbitrary base change.
\end{proposition}

\begin{proof}
\leavevmode
\begin{enumerate}[label=\textup{(\arabic*)}]
    \item \textbf{Establish the next claim.}
    The endofunctor \(\operatorname{dR}_S\) is left exact, so it preserves the
    binary product in the slice:
    \[
    \operatorname{dR}_S(X\times_SY)
    \simeq
    \operatorname{dR}_S(X)\times_S\operatorname{dR}_S(Y).
    \]
    \item \textbf{Establish the indicated equivalence.}
    This is the first equivalence, and under it the unit of \(X\times_SY\) is the
    product of the two relative units.
    \item \textbf{Introduce the notation.}
    Put \(Q_X:=Q_{X/S}\) and \(Q_Y:=Q_{Y/S}\), and write the two effective-image
    factorizations as
    \[
    X\twoheadrightarrow Q_X\lhook\joinrel\longrightarrow X_{\mathrm{dR}/S},
    \qquad
    Y\twoheadrightarrow Q_Y\lhook\joinrel\longrightarrow Y_{\mathrm{dR}/S}.
    \]
    \item \textbf{Construct the indicated map.}
    The product
    \[
    X\times_SY\longrightarrow Q_X\times_SQ_Y
    \]
    is an effective epimorphism: it is the composite of
    \[
    X\times_SY\longrightarrow Q_X\times_SY
    \]
    and
    \[
    Q_X\times_SY\longrightarrow Q_X\times_SQ_Y,
    \]
    each a pullback of one of the two effective epimorphisms.
    \item \textbf{Construct the indicated map.}
    The product of the two
    monomorphisms is a monomorphism.
    \item \textbf{Conclude the stated claim.}
    Thus
    \[
    X\times_SY
    \twoheadrightarrow
    Q_X\times_SQ_Y
    \lhook\joinrel\longrightarrow
    X_{\mathrm{dR}/S}\times_SY_{\mathrm{dR}/S}
    \]
    is the effective-epimorphism--monomorphism factorization of the product unit.
    \item \textbf{Use the stated construction.}
    Uniqueness of the factorization gives
    \[
    Q_{X\times_SY/S}\simeq Q_X\times_SQ_Y.
    \]
    \item \textbf{Establish the indicated equivalence.}
    Substitute the first equivalence into the iterated fibre-product definition of
    the relation nerve.
    \item \textbf{Use the stated construction.}
    Finite-limit pasting gives, degreewise,
    \[
    \mathcal R^{\mathrm{dR}}_{X\times_SY/S,n}
    \simeq
    \mathcal R^{\mathrm{dR}}_{X/S,n}
    \times_S
    \mathcal R^{\mathrm{dR}}_{Y/S,n}.
    \]
    \item \textbf{Apply naturality.}
    Naturality follows from finite products, pullback-stable effective-image
    factorization, and relative de Rham base change.
\end{enumerate}
\end{proof}

\begin{proposition}[External products]
For \(r,s\geq0\), the scalar multiplication and ordered Alexander--Whitney construction induce
natural products
\[
  \Omega^r_{\mathrm{dR}}(X/S)
  \otimes
  \Omega^s_{\mathrm{dR}}(Y/S)
  \longrightarrow
  \Omega^{r+s}_{\mathrm{dR}}(X\times_SY/S).
\]
They are natural in both variables and under arbitrary base change. The normalized
differential satisfies the signed Leibniz rule in the homotopy category.
\end{proposition}

\begin{proof}
\leavevmode
\begin{enumerate}[label=\textup{(\arabic*)}]
    \item \textbf{Establish the next claim.}
    For each \(n\), strong symmetric monoidality of relative suspension spectra and the
    product formula for de Rham relation degrees give
    \[
    \mathsf{R}^{\mathrm{dR},\mathrm{st}}_n(X\times_SY/S)
    \simeq
    \mathsf{R}^{\mathrm{dR},\mathrm{st}}_n(X/S)
    \otimes
    \mathsf{R}^{\mathrm{dR},\mathrm{st}}_n(Y/S).
    \]
    \item \textbf{Use the stated construction.}
    Evaluation followed by multiplication of the stable scalar therefore gives a
    canonical pairing
    \[
    \begin{aligned}
    &\HHom
    \bigl(\mathsf{R}^{\mathrm{dR},\mathrm{st}}_n(X/S),\mathcal O_S^{\mathrm{add}}\bigr)
    \otimes
    \HHom
    \bigl(\mathsf{R}^{\mathrm{dR},\mathrm{st}}_n(Y/S),\mathcal O_S^{\mathrm{add}}\bigr)
    \\
    &\qquad\longrightarrow
    \HHom
    \bigl(\mathsf{R}^{\mathrm{dR},\mathrm{st}}_n(X\times_SY/S),\mathcal O_S^{\mathrm{add}}\bigr).
    \end{aligned}
    \]
    \item \textbf{Invoke the cited result.}
    By Theorem~\ref{thm:internal-function-cochains-expanded}, these maps form a natural
    pairing of cosimplicial de Rham cochain objects.
    \item \textbf{Apply the indicated result.}
    Apply
    Theorem~\ref{thm:aw-product-expanded} to obtain the normalized product.
    \item \textbf{Apply naturality.}
    Naturality
    of the de Rham relation product, scalar multiplication, and the consecutive-block
    construction gives functoriality and arbitrary base-change compatibility, while
    the signed Leibniz rule is the one already proved for normalized Alexander--Whitney
    products.
\end{enumerate}
\end{proof}

\subsection{Closed forms, exact forms, and de Rham primitives}
\label{sec:closed-forms}

A normalized \(q\)-form records the \(q\)-th associated-graded layer, but
closedness requires more: it is a lift of that form to the entire \(q\)-th
stage of the complete totalization filtration. That lift retains a
nullhomotopy of the next boundary together with all higher compatibility
data. We package this retained coherence directly as the closed-form
object.

\begin{definition}[Closed de Rham forms]
For \(q\geq0\) and \(n\in\mathbb Z\), define the stable object
\[
  \Omega^{q,\mathrm{closed}}_{\mathrm{dR}}(X/S;n)
  :=
  F^q_{\mathrm{Tot}}\mathrm{DR}_{X/S}[n+q].
\]
The associated-graded quotient gives the underlying-form map
\[
  u_q:
  \Omega^{q,\mathrm{closed}}_{\mathrm{dR}}(X/S;n)
  \longrightarrow
  \Omega^q_{\mathrm{dR}}(X/S)[n].
\]
\end{definition}

\begin{remark}[Closed forms and higher lift data]
A global closed de Rham form is a point of
\[
  \Omega^\infty
  \Gamma^{\mathrm{st}}
  \bigl(S,\Omega^{q,\mathrm{closed}}_{\mathrm{dR}}(X/S;n)\bigr).
\]
The complete filtration identifies such a point with an underlying de Rham
\(q\)-form, a nullhomotopy of its boundary, and the higher compatibility data
encoded by the associated coherent cochain object.
\end{remark}

\begin{definition}[Exact-form lift]
For \(q\geq1\), the connecting morphism of
\[
  F^q_{\mathrm{Tot}}\mathrm{DR}
  \longrightarrow
  F^{q-1}_{\mathrm{Tot}}\mathrm{DR}
  \longrightarrow
  \Omega^{q-1}_{\mathrm{dR}}[1-q]
\]
defines, after shifting, a canonical morphism
\[
  d^{\mathrm{closed}}_{\mathrm{dR}}:
  \Omega^{q-1}_{\mathrm{dR}}(X/S)[n]
  \longrightarrow
  \Omega^{q,\mathrm{closed}}_{\mathrm{dR}}(X/S;n).
\]
\end{definition}

\begin{proposition}[Underlying boundary of the exact-form lift]
\label{prop:exact-classifier-boundary}
The composite
\[
  \Omega^{q-1}_{\mathrm{dR}}(X/S)[n]
  \xrightarrow{d^{\mathrm{closed}}_{\mathrm{dR}}}
  \Omega^{q,\mathrm{closed}}_{\mathrm{dR}}(X/S;n)
  \xrightarrow{u_q}
  \Omega^q_{\mathrm{dR}}(X/S)[n]
\]
is the homotopy-coherent de Rham differential \(d^{q-1}_{\mathrm{dR}}\). The lift to the
closed stage carries the canonical nullhomotopy of the next boundary and its higher
coherences.
\end{proposition}

\begin{proof}
\leavevmode
\begin{enumerate}[label=\textup{(\arabic*)}]
    \item \textbf{Construct the indicated map.}
    Under the complete-filtered/coherent-cochain correspondence, the shifted connecting
    morphism between adjacent filtration layers is precisely the cochain differential.
    \item \textbf{Construct the indicated map.}
    The fact that the connecting map lands in the entire stage
    \(F^q_{\mathrm{Tot}}[q+n]\), before projection to associated graded, is exactly the
    retained closed-form coherence of an exact form.
\end{enumerate}
\end{proof}

\begin{definition}[De Rham primitive spaces]
Let \(q\geq1\). Fix a global closed de Rham form
\[
  \omega\in
  \Omega^\infty
  \Gamma^{\mathrm{st}}
  \bigl(S,\Omega^{q,\mathrm{closed}}_{\mathrm{dR}}(X/S;n)\bigr).
\]
Define
\[
  \operatorname{PrimDiff}^q_{\mathrm{dR}}(X/S;n)
  :=
  \operatorname{fib}
  \Bigl(
  \Omega^{q-1}_{\mathrm{dR}}(X/S)[n]
  \xrightarrow{d^{\mathrm{closed}}_{\mathrm{dR}}}
  \Omega^{q,\mathrm{closed}}_{\mathrm{dR}}(X/S;n)
  \Bigr),
\]
and define the space of de Rham primitives of \(\omega\) by
\[
  \operatorname{PrimSpace}(\omega)
  :=
  \operatorname{hofib}_{\omega}
  \Bigl(
  \Omega^\infty\Gamma^{\mathrm{st}}
  \bigl(S,\Omega^{q-1}_{\mathrm{dR}}(X/S)[n]\bigr)
  \longrightarrow
  \Omega^\infty\Gamma^{\mathrm{st}}
  \bigl(S,\Omega^{q,\mathrm{closed}}_{\mathrm{dR}}(X/S;n)\bigr)
  \Bigr).
\]
\end{definition}

\begin{theorem}[De Rham primitive torsor theorem]
\label{thm:primitive-torsor-expanded}
If \(\operatorname{PrimSpace}(\omega)\) is nonempty, it is canonically a torsor
for the grouplike \(E_\infty\)-space
\[
  \Omega^\infty
  \Gamma^{\mathrm{st}}
  \bigl(S,\operatorname{PrimDiff}^q_{\mathrm{dR}}(X/S;n)\bigr).
\]
A chosen de Rham primitive identifies the torsor with this group object by translation. For
the zero closed form, the zero de Rham primitive makes the identification canonical.
\end{theorem}

\begin{proof}
\leavevmode
\begin{enumerate}[label=\textup{(\arabic*)}]
    \item \textbf{Construct the indicated map.}
    Stable sections preserve fibres because they are mapping spectra out of the tensor
    unit.
    \item \textbf{Construct the indicated map.}
    The exact-form lift therefore induces a morphism of grouplike
    \(E_\infty\)-spaces whose fibre over the identity is the displayed infinite-loop
    space.
    \item \textbf{Construct the indicated map.}
    Every nonempty homotopy fibre of a morphism of grouplike spaces is a torsor
    under the identity fibre.
\end{enumerate}
\end{proof}

\subsubsection{Additive delooping from the skeletal quotient}
\label{subsec:skeletal-delooping-closed-forms}

The additive classifier uses one feature of the stable affine-point
presentation that has not yet been needed: its canonical Postnikov structure
and the recognition of connective spectrum objects as commutative group
objects. We record that feature here, immediately before it becomes active.

\begin{proposition}[Postnikov structure and grouplike recognition]
\label{prop:stable-slice-postnikov-grouplike}
For every \(S\in\mathcal H_{\Sp}\), the stable slice
\(\mathcal D_{\Sp/S}\) carries its canonical Postnikov \(t\)-structure. Every
relative suspension spectrum \(\Sigma^\infty_{+,S}U\) is connective, and
\((\mathcal D_{\Sp/S})_{\geq0}\) is closed under small colimits and cofibres.
Moreover, zeroth space induces a canonical equivalence
\[
    (\mathcal D_{\Sp/S})_{\geq0}
    \xrightarrow{\simeq}
    \operatorname{CGrp}\bigl((\mathcal H_{\Sp})_{/S}\bigr).
  \]
These statements take place in the ordinary spectral Zariski \(\infty\)-topos.
\end{proposition}

\begin{proof}
\leavevmode
\begin{enumerate}[label=\textup{(\arabic*)}]
    \item \textbf{Establish the next claim.}
    Stabilization of an infinity-topos carries the canonical accessible Postnikov
    \(t\)-structure.
    \item \textbf{Establish the next claim.}
    Suspension spectra of pointed objects are connective, and
    the connective aisle is closed under small colimits and extensions, hence
    under cofibres; see \cite{LurieHigherAlgebra,LurieSAG}.
    \item \textbf{Establish the indicated equivalence.}
    The equivalence between
    connective spectrum objects and grouplike commutative monoid objects is natural
    in finite-limit-preserving functors.
    \item \textbf{Use the stated construction.}
    Applied to the spectral slice it gives
    \[
    (\mathcal D_{\Sp/S})_{\geq0}
    \simeq
    \operatorname{CGrp}\bigl((\mathcal H_{\Sp})_{/S}\bigr).
    \]
    \item \textbf{Establish the next claim.}
    This is the claimed grouplike recognition.
\end{enumerate}
\end{proof}

The stable skeletal formula itself is the primary classifier:
\[
  \Gamma^{\mathrm{st}}
  \bigl(S,\Omega^{q,\mathrm{closed}}_{\mathrm{dR}}(X/S;n)\bigr)
  \simeq
  \map_{\mathcal D_{\Sp/S}}
  \bigl(\mathsf{R}^{\mathrm{dR},\mathrm{st}}_{\geq q}(X/S),\mathcal O_S^{\mathrm{add}}[n+q]\bigr).
\]
By Proposition~\ref{prop:stable-slice-postnikov-grouplike},
\(\mathcal D_{\Sp/S}\) carries its canonical Postnikov \(t\)-structure, every relative
suspension spectrum \(\Sigma^\infty_{+,S}T\) is connective, and the connective
subcategory \((\mathcal D_{\Sp/S})_{\geq0}\) is closed under colimits and cofibres.
Hence
\[
  |\operatorname{sk}_{q-1}\mathsf{R}^{\mathrm{dR},\mathrm{st}}_\bullet(X/S)|
  \qquad\text{and}\qquad
  |\mathsf{R}^{\mathrm{dR},\mathrm{st}}_\bullet(X/S)|
\]
are connective, and the defining cofiber sequence
\[
  |\operatorname{sk}_{q-1}\mathsf{R}^{\mathrm{dR},\mathrm{st}}_\bullet(X/S)|
  \longrightarrow
  |\mathsf{R}^{\mathrm{dR},\mathrm{st}}_\bullet(X/S)|
  \longrightarrow
  \mathsf{R}^{\mathrm{dR},\mathrm{st}}_{\geq q}(X/S)
\]
shows that \(\mathsf{R}^{\mathrm{dR},\mathrm{st}}_{\geq q}(X/S)\) is connective.

\begin{construction}[Additive delooping object]
Define the commutative group object
\[
  \mathsf{R}^{\mathrm{dR},\langle q\rangle}_{X/S}
  :=
  \Omega^\infty \mathsf{R}^{\mathrm{dR},\mathrm{st}}_{\geq q}(X/S)
  \in
  \operatorname{CGrp}
  \bigl((\mathcal H_{\Sp})_{/S}\bigr).
\]
This is the grouplike object corresponding to the connective spectrum
\(\mathsf{R}^{\mathrm{dR},\mathrm{st}}_{\geq q}(X/S)\).
\end{construction}

\begin{theorem}[Additive-delooping classifier]
\label{thm:additive-delooping-expanded}
For every \(q\geq0\) and \(n\in\mathbb Z\), there is a canonical equivalence of
spaces
\[
\begin{aligned}
  &\Omega^\infty
  \Gamma^{\mathrm{st}}
  \bigl(S,\Omega^{q,\mathrm{closed}}_{\mathrm{dR}}(X/S;n)\bigr)
  \\
  &\qquad\simeq
  \Map_{\operatorname{CGrp}((\mathcal H_{\Sp})_{/S})}
  \Bigl(
     \mathsf{R}^{\mathrm{dR},\langle q\rangle}_{X/S},
     \Omega^\infty\bigl(\mathcal O_S^{\mathrm{add}}[n+q]\bigr)
  \Bigr).
\end{aligned}
\]
For \(q=0\), the stable source is
\[
  \mathsf{R}^{\mathrm{dR},\mathrm{st}}_{\geq0}(X/S)
  \simeq
  \Sigma^\infty_{+,S}Q_{X/S},
\]
so the stable formula reduces to additive scalar functions on the effective quotient.
\end{theorem}

\begin{proof}
\leavevmode
\begin{enumerate}[label=\textup{(\arabic*)}]
    \item \textbf{Establish the next claim.}
    The stable formula is Theorem~\ref{thm:skeletal-totalization-presentation}.
    \item \textbf{Invoke the cited result.}
    By the
    connectivity calculation above, the source \(\mathsf{R}^{\mathrm{dR},\mathrm{st}}_{\geq q}(X/S)\) belongs to
    \((\mathcal D_{\Sp/S})_{\geq0}\).
    \item \textbf{Conclude the stated claim.}
    Hence, for every \(E\in\mathcal D_{\Sp/S}\), the canonical
    connective-cover map induces an equivalence
    \[
    \Map_{\mathcal D_{\Sp/S}}
    \bigl(\mathsf{R}^{\mathrm{dR},\mathrm{st}}_{\geq q}(X/S),\tau_{\geq0}E\bigr)
    \xrightarrow{\simeq}
    \Map_{\mathcal D_{\Sp/S}}
    \bigl(\mathsf{R}^{\mathrm{dR},\mathrm{st}}_{\geq q}(X/S),E\bigr),
    \]
    because maps from a connective object to an object of
    \((\mathcal D_{\Sp/S})_{\leq-1}\) are contractible.
    \item \textbf{Apply the cited proposition.}
    Proposition~\ref{prop:stable-slice-postnikov-grouplike} identifies
    \((\mathcal D_{\Sp/S})_{\geq0}\) with the commutative group objects in the slice.
    \item \textbf{Apply the indicated result.}
    Apply the preceding mapping equivalence with
    \[
    E=\mathcal O_S^{\mathrm{add}}[n+q].
    \]
    \item \textbf{Establish the indicated equivalence.}
    Under the grouplike-recognition equivalence, the connective spectrum
    \(\mathsf{R}^{\mathrm{dR},\mathrm{st}}_{\geq q}(X/S)\) corresponds to the grouplike object
    \[
    \mathsf{R}^{\mathrm{dR},\langle q\rangle}_{X/S}
    =
    \Omega^\infty \mathsf{R}^{\mathrm{dR},\mathrm{st}}_{\geq q}(X/S),
    \]
    and the connective cover of the target corresponds to the grouplike object
    \[
    \Omega^\infty
    \bigl(\mathcal O_S^{\mathrm{add}}[n+q]\bigr).
    \]
    \item \textbf{Use the stated construction.}
    This gives the displayed mapping-space equivalence.
    \item \textbf{Identify the degree-zero case.}
    For \(q=0\), the \((-1)\)-skeleton is initial, so
    \[
    \mathsf{R}^{\mathrm{dR},\mathrm{st}}_{\geq0}(X/S)
    \simeq
    |\mathsf{R}^{\mathrm{dR},\mathrm{st}}_\bullet(X/S)|
    \simeq
    \Sigma^\infty_{+,S}Q_{X/S},
    \]
    using
    \(|\mathcal R^{\mathrm{dR}}_{X/S,\bullet}|\simeq Q_{X/S}\) and preservation of
    geometric realizations by relative suspension spectra.
\end{enumerate}
\end{proof}

\begin{remark}[Skeletal source in positive degree]
For \(q>0\), the classifier is the additive delooping
\(\mathsf{R}^{\mathrm{dR},\langle q\rangle}_{X/S}\) of the normalized skeletal quotient. The removal of
lower-dimensional and degenerate simplices is therefore part of the positive-degree
classification.
\end{remark}


All coefficient objects of the de Rham calculus have now been constructed, so
their structural behaviour can be stated without referring to objects that
have not yet appeared. Contravariance in \(X\) comes from functoriality of the
relative unit and effective-image factorization; substitution in the base
comes from relative de Rham base change and right Beck--Chevalley. The basic
geometric situation is the following Cartesian square:
\[
\begin{tikzcd}[column sep=huge,row sep=large]
  X' \arrow[r] \arrow[d]
  & X \arrow[d]
  \\
  S' \arrow[r]
  & S,
\end{tikzcd}
\qquad
X'_{\mathrm{dR}/S'}
\simeq
X_{\mathrm{dR}/S}\times_SS'.
\]
Because every later construction is obtained from the degreewise cochain
diagram by limits, fibres, shifts, and the complete-filtered/coherent-cochain
correspondence, the degreewise comparison propagates through the entire
calculus.
\subsection{Substitution and base change of the completed calculus}
\subsubsection{Contravariance in the geometric variable}
\label{subsec:contravariance-geometric-variable}

Let \(f:X\to Y\) be a morphism over \(S\). Naturality of the relative de Rham
unit gives
\[
\begin{tikzcd}
X \ar[r,"f"] \ar[d,"\eta_{X/S}^{\mathrm{dR}}"']
  & Y \ar[d,"\eta_{Y/S}^{\mathrm{dR}}"]\\
X_{\mathrm{dR}/S} \ar[r]
  & Y_{\mathrm{dR}/S}.
\end{tikzcd}
\]
Functorial effective-image factorization gives
\(Q_{X/S}\to Q_{Y/S}\), and finite limits give a simplicial map
\[
  \mathcal R^{\mathrm{dR}}_{X/S,\bullet}
  \longrightarrow
  \mathcal R^{\mathrm{dR}}_{Y/S,\bullet}.
\]
Scalar functions are contravariant, so these maps induce
\[
  f^*_{\mathrm{dR}}:
  C^\bullet_{\mathrm{dR}}(Y/S)
  \longrightarrow
  C^\bullet_{\mathrm{dR}}(X/S)
\]
as a morphism of cosimplicial commutative algebras.

\begin{theorem}[Contravariant functoriality]
\label{thm:contravariant-functoriality-expanded}
The preceding assignment is contravariant in \(X\). The induced maps
commute with the Amitsur presentation, full-to-effective comparison, partial and
complete totalizations, normalized forms, homotopy-coherent differentials, closed stages,
exact-form lifts, and Alexander--Whitney products. For composable
\(X\xrightarrow fY\xrightarrow gZ\),
\[
  (gf)^*_{\mathrm{dR}}
  \simeq
  f^*_{\mathrm{dR}}g^*_{\mathrm{dR}}
\]
by functorial image factorization, finite limits, and adjoint mates.
\end{theorem}

\begin{proof}
\leavevmode
\begin{enumerate}[label=\textup{(\arabic*)}]
    \item \textbf{Establish the next claim.}
    The geometric terms are functorial constructions from the natural transformation
    \(\eta_{-/S}^{\mathrm{dR}}\), the functorial
    \((\mathrm{effective\ epimorphism},\mathrm{monomorphism})\)-factorization, and finite limits.
    \item \textbf{Construct the indicated map.}
    The
    cochain maps are the resulting scalar restriction maps.
    \item \textbf{Construct the indicated map.}
    All subsequent objects are
    limits, fibres, shifts, or natural maps of the same diagrams, and the homotopy-coherent
    differential is functorial under exact maps of complete filtrations.
\end{enumerate}
\end{proof}

\subsubsection{Arbitrary base change in the base}

\begin{theorem}[Parameterized base change]
\label{thm:parameterized-de-Rham-base-change}
Let \(c:S'\to S\) and \(X'=X\times_SS'\). There are coherent equivalences
\[
  c^*C^n_{\mathrm{dR}}(X/S)
  \xrightarrow{\simeq}
  C^n_{\mathrm{dR}}(X'/S')
\]
for all \(n\), assembling into an equivalence of cosimplicial commutative algebras.
Consequently
\[
  c^*\mathrm{DR}_{X/S}
  \simeq
  \mathrm{DR}_{X'/S'},
  \qquad
  c^*\mathrm{DR}^{<q}_{X/S}
  \simeq
  \mathrm{DR}^{<q}_{X'/S'},
\]
\[
  c^*\Omega^q_{\mathrm{dR}}(X/S)
  \simeq
  \Omega^q_{\mathrm{dR}}(X'/S'),
\]
and
\[
  c^*\Omega^{q,\mathrm{closed}}_{\mathrm{dR}}(X/S;n)
  \simeq
  \Omega^{q,\mathrm{closed}}_{\mathrm{dR}}(X'/S';n),
\]
compatibly with the homotopy-coherent differential, exact-form lifts, and normalized
Alexander--Whitney products. For morphisms of spectral schemes, the same statement
holds for the first-order cotangent comparison.
\end{theorem}

\begin{proof}
\leavevmode
\begin{enumerate}[label=\textup{(\arabic*)}]
    \item \textbf{Establish multiplicative compatibility.}
    The de Rham relation degrees commute with base change by
    Theorem~\ref{thm:relative-de-Rham-base-change} and finite-limit pasting:
    \[
    \mathcal R^{\mathrm{dR}}_{X'/S',n}
    \simeq
    \mathcal R^{\mathrm{dR}}_{X/S,n}\times_SS'.
    \]
    \item \textbf{Establish the next claim.}
    The scalar coefficient satisfies
    \(c^*\mathcal O_S^{\mathrm{add}}\simeq\mathcal O_{S'}^{\mathrm{add}}\).
    \item \textbf{Use the stated construction.}
    Right Beck--Chevalley from
    Theorem~\ref{thm:stable-dependent-adjoints} therefore gives the degreewise cochain
    equivalence.
    \item \textbf{Use pullback stability.}
    The pullback functor has both adjoints between stable categories and
    hence preserves all limits and colimits; it consequently commutes with partial and
    complete totalizations, matching fibres, filtration fibres, shifts, geometric
    realizations, and skeletal quotients.
    \item \textbf{Apply naturality.}
    Naturality gives compatibility with every
    structure map.
    \item \textbf{Establish multiplicative compatibility.}
    The first structured infinitesimal neighbourhood and cotangent complex
    commute with Cartesian base change of spectral schemes by
    Chapter~\Ref{chap:differentiation-spectral-doctrine}.
\end{enumerate}
\end{proof}

\begin{construction}[Restriction on global section spectra]
The coefficient equivalences above induce natural restriction morphisms
\[
  \Gamma^{\mathrm{st}}
  \bigl(S,\Omega^q_{\mathrm{dR}}(X/S)[n]\bigr)
  \longrightarrow
  \Gamma^{\mathrm{st}}
  \bigl(S',\Omega^q_{\mathrm{dR}}(X'/S')[n]\bigr),
\]
and similarly for closed forms and de Rham algebras, by pulling a section along
\(c^*\) and then using the base-change equivalence. These maps are the natural restriction morphisms on global section spectra.
\end{construction}


The complete descent calculus is now stable under the two substitutions that
matter later: contravariant change of the geometric object and Cartesian
change of base. Its multiplication is likewise compatible with products.
The remaining global issue is descent of these form objects with respect to
covers of the geometric context. We address that next, together with the
affine-descent step required for comprehension.
\section{Formal \texorpdfstring{\'etale}{etale} descent and comprehension of de Rham forms}
\label{sec:moduli-forms}

The intrinsic form coefficients are defined for every object of the spectral
slice, and comprehension requires one further step: their stable-section
semantics on affine tests must satisfy spectral Zariski descent. Establishing
that descent requires an additional argument beyond base change. We therefore
first control the de Rham relation under formally \(\acute{e}\)tale effective
epimorphisms, the class that contains the finite Zariski covering morphisms used
below.

Only two semantic levels will be distinguished in this section. The
\emph{stable-section semantics} is obtained directly from the intrinsic stable
coefficient object. The \emph{comprehended semantics} is its affine-descent
extension represented by an object of the slice. Formally \(\acute{e}\)tale
transport and finite-coproduct compatibility provide the descent bridge;
effective descent reconstructs the cochain tower degreewise; affine density
then produces the representing comprehension objects.

The logic of the last step is expressed by the following square. The lower
functor is the right Kan extension of the affine stable-section semantics:
\[
\begin{tikzcd}[column sep=huge,row sep=large]
  \operatorname{AffPt}(S)^{\op}
      \arrow[r,"{\mathfrak{Form}^{q,\mathrm{stsec}}_{\mathrm{dR},S}}"]
      \arrow[d,"j_S^{\op}"']
  &
  \mathcal S
  \\
  ((\mathcal H_{\Sp})_{/S})^{\op}
      \arrow[r,"{\operatorname{Ran}_{j_S^{\op}}}"']
  &
  \mathcal S
      \arrow[u,equal].
\end{tikzcd}
\]
The representing object obtained from the affine-point sheaf therefore
comprehends exactly this affine-descended extension.
\subsection{Formally \texorpdfstring{\'etale}{etale} transport}

Theorem~\ref{thm:formal-etale-neighbourhood-invariance} gives the precise geometric
invariance under a formally \(\acute{e}\)tale morphism \(u:U\to X\):
\[
  \mathcal R^{\mathrm{dR}}_{U/S,n}
  \simeq
  U\times_{X,v_n}\mathcal R^{\mathrm{dR}}_{X/S,n}.
\]
This Cartesian formula controls the geometric transport of the de Rham relation.

\begin{proposition}[Degreewise centred coefficient base change]
\label{prop:centered-degreewise-etale-base-change}
For \(n\geq0\), let \(v_n:\mathcal R^{\mathrm{dR}}_{X/S,n}\to X\) be the first-vertex projection and
put
\[
  C^{n,\mathrm{ctr}}_{\mathrm{dR}}(X/S)
  :=
  \Pi_{v_n}v_n^*\mathcal O_X^{\mathrm{add}}
  \in\mathcal D_{\Sp/X}.
\]
If \(u:U\to X\) is formally \(\acute{e}\)tale over \(S\), then
\[
  u^*C^{n,\mathrm{ctr}}_{\mathrm{dR}}(X/S)
  \xrightarrow{\simeq}
  C^{n,\mathrm{ctr}}_{\mathrm{dR}}(U/S)
\]
canonically in each degree.
\end{proposition}

\begin{proof}
\leavevmode
\begin{enumerate}[label=\textup{(\arabic*)}]
    \item \textbf{Establish the next claim.}
    Use
    \(\mathcal R^{\mathrm{dR}}_{U/S,n}\simeq
    U\times_X\mathcal R^{\mathrm{dR}}_{X/S,n}\) and right Beck--Chevalley for
    the Cartesian square defined by the first-vertex projection, together with
    \(u^*\mathcal O_X^{\mathrm{add}}\simeq\mathcal O_U^{\mathrm{add}}\).
\end{enumerate}
\end{proof}

\begin{remark}[First-vertex indexing]
The objects \(C^{n,\mathrm{ctr}}_{\mathrm{dR}}(X/S)\) are indexed by the distinguished first
vertex. The face deleting that vertex changes the chosen centre, so the complete
cosimplicial algebra is formed after dependent product to the base:
\[
  C^\bullet_{\mathrm{dR}}(X/S)\in\CAlg(\mathcal D_{\Sp/S}).
\]
A formally \(\acute{e}\)tale map \(U\to X\) induces the corresponding
contravariant morphism of cosimplicial de Rham algebras.
\end{remark}

\begin{theorem}[Formal \texorpdfstring{\'etale}{etale} change of intermediate base]
\label{thm:formal-etale-intermediate-base}
Let \(g:S'\to S\) be formally \(\acute{e}\)tale and let \(V\to S'\). Regard
\(V\) also as an object over \(S\). Then
\[
  V_{\mathrm{dR}/S'}
  \xrightarrow{\simeq}
  V_{\mathrm{dR}/S},
\]
with the two relative units identified. Consequently
\[
  Q_{V/S'}\xrightarrow{\simeq}Q_{V/S},
  \qquad
  \mathcal R^{\mathrm{dR}}_{V/S',\bullet}\xrightarrow{\simeq}\mathcal R^{\mathrm{dR}}_{V/S,\bullet},
\]
and for every \(Z\to Y\to S'\),
\[
  \operatorname{Nbh}^{\mathrm{dR}}_{Z/S'}(Y)
  \xrightarrow{\simeq}
  \operatorname{Nbh}^{\mathrm{dR}}_{Z/S}(Y).
\]
For coefficient objects the correctly typed comparison is
\[
  \Pi_g C^\bullet_{\mathrm{dR}}(V/S')
  \xrightarrow{\simeq}
  C^\bullet_{\mathrm{dR}}(V/S).
\]
It induces corresponding equivalences
\[
  \Pi_g\mathrm{DR}_{V/S'}
  \simeq
  \mathrm{DR}_{V/S},
  \qquad
  \Pi_g\Omega^q_{\mathrm{dR}}(V/S')
  \simeq
  \Omega^q_{\mathrm{dR}}(V/S),
\]
and similarly for partial totalizations, closed stages, and exact-form lifts.
On stable global sections, the two intermediate-base constructions therefore agree
canonically.
\end{theorem}

\begin{proof}
\leavevmode
\begin{enumerate}[label=\textup{(\arabic*)}]
    \item \textbf{Apply the indicated result.}
    Apply Theorem~\ref{thm:formal-etale-cartesianity} to \(g:S'\to S\) over
    the terminal object.
    \item \textbf{Use the stated construction.}
    It gives
    \[
    S'
    \simeq
    S\times_{S_{\mathrm{dR}}}S'_{\mathrm{dR}}.
    \]
    Hence
    \[
    \begin{aligned}
    V_{\mathrm{dR}/S'}
    &=
    S'\times_{S'_{\mathrm{dR}}}V_{\mathrm{dR}}\\
    &\simeq
    \bigl(S\times_{S_{\mathrm{dR}}}S'_{\mathrm{dR}}\bigr)
    \times_{S'_{\mathrm{dR}}}V_{\mathrm{dR}}\\
    &\simeq
    S\times_{S_{\mathrm{dR}}}V_{\mathrm{dR}}
    =
    V_{\mathrm{dR}/S}.
    \end{aligned}
    \]
    \item \textbf{Establish the next claim.}
    The unit in both cases is induced by the same absolute unit \(V\to V_{\mathrm{dR}}\).
    \item \textbf{Establish the next claim.}
    Uniqueness of effective-image factorization and finite-limit pasting give the quotient,
    de Rham relation, and de Rham neighbourhood formulas.
    \item \textbf{Introduce the notation.}
    Write
    \[
    r_n:\mathcal R^{\mathrm{dR}}_{V/S',n}\longrightarrow S',
    \qquad
    a_n:=gr_n:\mathcal R^{\mathrm{dR}}_{V/S',n}\longrightarrow S.
    \]
    Since
    \(\mathcal O_{S'}^{\mathrm{add}}\simeq g^*\mathcal O_S^{\mathrm{add}}\),
    \[
    \begin{aligned}
    \Pi_gC^n_{\mathrm{dR}}(V/S')
    &\simeq
    \Pi_g\Pi_{r_n}r_n^*\mathcal O_{S'}^{\mathrm{add}}\\
    &\simeq
    \Pi_{a_n}a_n^*\mathcal O_S^{\mathrm{add}}\\
    &=
    C^n_{\mathrm{dR}}(V/S).
    \end{aligned}
    \]
    \item \textbf{Identify the stated objects.}
    Composition of dependent products makes this identification cosimplicial and
    multiplicative.
    \item \textbf{Establish multiplicative compatibility.}
    The functor \(\Pi_g\) preserves all limits and is exact, so it
    commutes with totalization, normalization, and filtration fibres.
    \item \textbf{Establish the next claim.}
    Finally
    \[
    \Gamma^{\mathrm{st}}(S,\Pi_gE)
    \simeq
    \Gamma^{\mathrm{st}}(S',E)
    \]
    by adjunction.
\end{enumerate}
\end{proof}

\subsection{Finite disjoint coproducts}

Finite affine Zariski covers enter the slice through a coproduct of their
members. To use formal \(\acute{e}\)tale descent for such a covering map, we
must know that maps into a finite coproduct decompose compatibly across an
infinitesimal substitution. The next lemma transports the corresponding
clopen decomposition from the centre to the thickening.

\begin{lemma}[Finite clopen decompositions under infinitesimal substitution]
\label{lem:finite-clopen-decomposition-infinitesimal}
Let
\[
  j:V\longrightarrow T
\]
be a finite spectral infinitesimal substitution between affine spectral schemes, and
let \(\{F_i\to S\}_{i\in I}\) be a finite family of objects of the slice
\((\mathcal H_{\Sp})_{/S}\).

Every morphism
\[
  q:V\longrightarrow\coprod_{i\in I}F_i
\]
over \(S\) canonically determines a finite disjoint open-and-closed decomposition
\[
  V\simeq\coprod_{i\in I}V_i
\]
and morphisms
\[
  q_i:V_i\longrightarrow F_i.
\]
Moreover this decomposition is the pullback, through a contractible space of
choices, of a unique finite disjoint open-and-closed decomposition
\[
  T\simeq\coprod_{i\in I}T_i
\]
such that
\[
  V_i\simeq V\times_TT_i.
\]
Every \(T_i\) and \(V_i\) is affine, and
\[
  V_i\longrightarrow T_i
\]
is again a finite spectral infinitesimal substitution. If \(j\) is a one-stage
square-zero infinitesimal substitution, then every \(V_i\to T_i\) is a one-stage
square-zero infinitesimal substitution.
\end{lemma}

\begin{proof}
\leavevmode
\begin{enumerate}[label=\textup{(\arabic*)}]
    \item \textbf{Fix the indicated data.}
    Let
    \[
    I_S:=\coprod_{i\in I}S
    \]
    be the finite constant object of the slice.
    \item \textbf{Construct the indicated map.}
    The structure maps \(F_i\to S\)
    assemble to a canonical morphism
    \[
    \coprod_{i\in I}F_i\longrightarrow I_S
    \]
    which on the \(i\)-th summand lands in the \(i\)-th copy of \(S\).
    \item \textbf{Establish the next claim.}
    Compose with
    \(q\) to obtain
    \[
    \bar q:V\longrightarrow I_S.
    \]
    \item \textbf{Use the stated identification.}
    Because finite coproducts in an infinity-topos are disjoint and universal, pulling
    back the \(i\)-th summand \(S\to I_S\) along \(\bar q\) gives pairwise disjoint
    complemented subobjects
    \[
    V_i\longrightarrow V
    \]
    whose coproduct is \(V\).
    \item \textbf{Construct the indicated map.}
    The original map \(q\) restricts canonically to
    \[
    q_i:V_i\longrightarrow F_i.
    \]
    \item \textbf{Establish the next claim.}
    On the spectral Zariski site of the affine scheme \(V\), a section of the finite
    constant sheaf \(I\) is equivalently a finite locally constant \(I\)-valued function,
    and hence equivalently a finite disjoint clopen decomposition of \(V\).
    \item \textbf{Conclude the stated claim.}
    Thus the
    subobjects \(V_i\) are represented by open-and-closed spectral subschemes.
    \item \textbf{Use the stated hypothesis.}
    Since
    \(V\) is affine, each \(V_i\) is affine: its clopen support is represented by the
    spectral localization at the corresponding idempotent in
    \(\pi_0\mathcal O(V)\).
    \item \textbf{Invoke the cited result.}
    By
    \Ref{prop:finite-infinitesimal-zariski-invariance}, pullback along \(j\) induces an
    equivalence between the posets of open spectral subschemes of \(T\) and \(V\),
    preserving finite intersections and covering families.
    \item \textbf{Conclude the stated claim.}
    Hence each \(V_i\) is the
    pullback of a contractibly unique open spectral subscheme
    \[
    T_i\longrightarrow T.
    \]
    \item \textbf{Establish the indicated equivalence.}
    The equalities
    \[
    V_i\cap V_k=\varnothing
    \qquad(i\neq k)
    \]
    and the covering relation
    \[
    \coprod_iV_i\longrightarrow V
    \]
    transport through this equivalence to
    \[
    T_i\cap T_k=\varnothing
    \qquad(i\neq k)
    \]
    and to a covering family \(\{T_i\to T\}\).
    Thus
    \[
    T\simeq\coprod_iT_i.
    \]
    \item \textbf{Establish the next claim.}
    Each \(T_i\) is open and its complement is the union of the remaining \(T_k\), so it
    is open and closed.
    \item \textbf{Use the stated hypothesis.}
    Since \(T\) is affine, every \(T_i\) is affine by the same
    idempotent-localization argument.
    \item \textbf{Establish the next claim.}
    Finally,
    \[
    V_i\simeq V\times_TT_i
    \]
    by construction.
    \item \textbf{Use the stated construction.}
    Base-change stability of finite infinitesimal substitutions gives
    that \(V_i\to T_i\) is finite infinitesimal.
    \item \textbf{Use the stated construction.}
    In the one-stage case,
    base-change stability of square-zero infinitesimal substitutions gives the final
    assertion.
\end{enumerate}
\end{proof}

\begin{proposition}[Finite-coproduct compatibility]
\label{prop:finite-coproduct-compatibility}
For a finite disjoint family \(\{X_i\to S\}\),
\[
  \left(\coprod_iX_i\right)_{\mathrm{dR}/S}
  \simeq
  \coprod_i(X_i)_{\mathrm{dR}/S},
\]
\[
  Q_{(\coprod_iX_i)/S}
  \simeq
  \coprod_iQ_{X_i/S},
\]
and
\[
  \mathcal R^{\mathrm{dR}}_{(\coprod_iX_i)/S,n}
  \simeq
  \coprod_i\mathcal R^{\mathrm{dR}}_{X_i/S,n}.
\]
de Rham scalar functions send these finite disjoint coproducts to finite products of
coefficient objects, compatibly with normalization and totalization.
\end{proposition}

\begin{proof}
\leavevmode
\begin{enumerate}[label=\textup{(\arabic*)}]
    \item \textbf{Prove locality of a finite coproduct.}
    We first prove that a finite disjoint coproduct of objects de Rham-local over \(S\) is
    again de Rham-local over \(S\).
    \item \textbf{Fix the indicated data.}
    Let
    \[
    F:=\coprod_iF_i\longrightarrow S
    \]
    with every \(F_i\to S\) de Rham-local over \(S\).
    \item \textbf{Fix the indicated data.}
    Let
    \[
    j:V\longrightarrow T
    \]
    be an affine square-zero infinitesimal substitution over \(S\).
    \item \textbf{Establish the indicated equivalence.}
    We show that
    \[
    \Map_S(T,F)\longrightarrow\Map_S(V,F)
    \]
    is an equivalence.
    \item \textbf{Fix the indicated data.}
    Fix a point
    \[
    q:V\longrightarrow F.
    \]
    \item \textbf{Invoke the cited result.}
    By
    Lemma~\ref{lem:finite-clopen-decomposition-infinitesimal}, \(q\) determines a
    finite disjoint open-and-closed decomposition
    \[
    V\simeq\coprod_iV_i
    \]
    together with maps
    \[
    q_i:V_i\longrightarrow F_i.
    \]
    \item \textbf{Establish the next claim.}
    The same lemma transports this decomposition through \(j\) to a contractibly unique
    finite disjoint open-and-closed decomposition
    \[
    T\simeq\coprod_iT_i
    \]
    such that
    \[
    V_i\simeq V\times_TT_i.
    \]
    \item \textbf{Establish the next claim.}
    Every \(T_i\) is affine, and each
    \[
    V_i\longrightarrow T_i
    \]
    is an affine square-zero infinitesimal substitution.
    \item \textbf{Use the stated construction.}
    De Rham locality of \(F_i\) over \(S\) therefore gives a contractibly unique extension
    \[
    \widetilde q_i:T_i\longrightarrow F_i
    \]
    of \(q_i\).
    \item \textbf{Construct the indicated map.}
    The maps \(\widetilde q_i\) assemble uniquely to an extension
    \[
    \widetilde q:T\simeq\coprod_iT_i
    \longrightarrow
    \coprod_iF_i=F.
    \]
    \item \textbf{Establish the next claim.}
    The transported clopen decomposition is contractibly unique, and each component
    extension space is contractible.
    \item \textbf{Conclude the stated claim.}
    Hence the homotopy fibre over \(q\) of
    \[
    \Map_S(T,F)\longrightarrow\Map_S(V,F)
    \]
    is contractible.
    \item \textbf{Use the stated hypothesis.}
    Since this holds for every \(q\), the restriction map is an
    equivalence.
    \item \textbf{Invoke the cited result.}
    By Theorem~\ref{thm:slice-de-Rham}, \(F\to S\) is de Rham-local
    over \(S\).
    \item \textbf{Apply the indicated result.}
    Apply this to
    \[
    F_i=(X_i)_{\mathrm{dR}/S}.
    \]
    \item \textbf{Use the stated construction.}
    The coproduct of the localization units gives
    \[
    \coprod_iX_i
    \longrightarrow
    \coprod_i(X_i)_{\mathrm{dR}/S}.
    \]
    \item \textbf{Construct the indicated map.}
    For every \(Z\to S\) de Rham-local over \(S\), the induced mapping-space comparison is
    the finite product of the componentwise localization equivalences:
    \[
    \begin{aligned}
    \Map_S\left(\coprod_i(X_i)_{\mathrm{dR}/S},Z\right)
    &\simeq
    \prod_i
    \Map_S\bigl((X_i)_{\mathrm{dR}/S},Z\bigr)
    \\
    &\simeq
    \prod_i\Map_S(X_i,Z)
    \\
    &\simeq
    \Map_S\left(\coprod_iX_i,Z\right).
    \end{aligned}
    \]
    \item \textbf{Conclude the stated claim.}
    Thus the coproduct map is a de Rham equivalence over \(S\) with target de Rham-local
    over \(S\).
    \item \textbf{Establish the locality claim.}
    It therefore exhibits the relative de Rham localization:
    \[
    \left(\coprod_iX_i\right)_{\mathrm{dR}/S}
    \simeq
    \coprod_i(X_i)_{\mathrm{dR}/S}.
    \]
    \item \textbf{Establish the next claim.}
    For each \(i\), write the effective-image factorization as
    \[
    X_i
    \xrightarrow{e_i}
    Q_{X_i/S}
    \xrightarrow{m_i}
    (X_i)_{\mathrm{dR}/S}.
    \]
    \item \textbf{Construct the indicated map.}
    Finite coproducts of effective epimorphisms are effective epimorphisms in an
    infinity-topos.
    \item \textbf{Invoke the cited result.}
    By disjointness and universality of finite coproducts, the
    coproduct of the monomorphisms \(m_i\) is again a monomorphism.
    \item \textbf{Conclude the stated claim.}
    Hence
    \[
    \coprod_iX_i
    \longrightarrow
    \coprod_iQ_{X_i/S}
    \longrightarrow
    \coprod_i(X_i)_{\mathrm{dR}/S}
    \]
    is the effective-epimorphism--monomorphism factorization of the relative de Rham
    unit of \(\coprod_iX_i\).
    \item \textbf{Use the stated construction.}
    Functorial uniqueness of effective images therefore gives
    \[
    Q_{(\coprod_iX_i)/S}
    \simeq
    \coprod_iQ_{X_i/S}.
    \]
    \item \textbf{Compute the relation degrees.}
    Using the relative de Rham coproduct formula and disjointness of the summands,
    disjointness eliminates mixed components from the iterated fibre products. Therefore
    \[
    \begin{aligned}
    \mathcal R^{\mathrm{dR}}_{(\coprod_iX_i)/S,n}
    &=
    \underbrace{
    \left(\coprod_iX_i\right)
    \times_{(\coprod_iX_i)_{\mathrm{dR}/S}}
    \cdots
    \times_{(\coprod_iX_i)_{\mathrm{dR}/S}}
    \left(\coprod_iX_i\right)
    }_{n+1\text{ factors}}
    \\
    &\simeq
    \coprod_i
    \underbrace{
    X_i\times_{(X_i)_{\mathrm{dR}/S}}
    \cdots
    \times_{(X_i)_{\mathrm{dR}/S}}X_i
    }_{n+1\text{ factors}}
    \\
    &=
    \coprod_i\mathcal R^{\mathrm{dR}}_{X_i/S,n}.
    \end{aligned}
    \]
    \item \textbf{Establish the next claim.}
    Finally, dependent product of the stable additive scalar along a finite disjoint
    coproduct is the finite product of the corresponding dependent products.
    \item \textbf{Conclude the stated claim.}
    Hence
    de Rham scalar functions carry the displayed de Rham relation decomposition to finite
    products of coefficient objects.
    \item \textbf{Identify the stated objects.}
    Stable normalization, partial totalization, and
    complete totalization are limit constructions, so these identifications commute
    with normalization and totalization.
\end{enumerate}
\end{proof}

\subsection{Effective formal \texorpdfstring{\'etale}{etale} descent}
\label{subsec:effective-formally-etale-descent}

The preceding coproduct analysis shows that finite Zariski covering morphisms
are formally \(\acute{e}\)tale. We can now combine formal
\(\acute{e}\)tale transport with effective descent. The result reconstructs
the effective quotient and every relation degree from the \v{C}ech nerve of
the cover, and hence reconstructs the entire cochain tower.

\begin{lemma}[Spectral Zariski covers are formally \(\acute{e}\)tale]
\label{lem:zariski-open-formally-etale}
Every spectral Zariski open immersion is formally \(\acute{e}\)tale with respect to
the finite spectral infinitesimal substitutions of
Chapter~\Ref{chap:differentiation-spectral-doctrine}. More generally, if
\(\{U_i\to X\}_{i\in I}\) is a finite spectral Zariski covering family, then the
covering morphism
\[
  \coprod_{i\in I}U_i\longrightarrow X
\]
is formally \(\acute{e}\)tale.
\end{lemma}

\begin{proof}
\leavevmode
\begin{enumerate}[label=\textup{(\arabic*)}]
    \item \textbf{Invoke the cited result.}
    By
    \Ref{prop:formal-lifting-one-stage-detection}, it is enough to test
    contractible lifting against affine one-stage square-zero substitutions
    \(j:V\to T\).
    \item \textbf{Fix the indicated data.}
    Let \(U\hookrightarrow X\) be a spectral Zariski open immersion
    and consider such a lifting problem.
    \item \textbf{Use pullback stability.}
    The inverse image
    \[
    T\times_XU\longrightarrow T
    \]
    is an open subobject whose pullback to \(V\) is all of \(V\).
    \item \textbf{Invoke the cited result.}
    By
    \Ref{prop:square-zero-substitution-permanence}, pullback along \(j\) identifies
    the small spectral Zariski sites of \(T\) and \(V\).
    \item \textbf{Conclude the stated claim.}
    Hence the displayed open
    subobject is all of \(T\).
    \item \textbf{Establish the next claim.}
    The required factorization through \(U\) therefore
    exists, and its space is contractible because an open immersion is a
    monomorphism.
    \item \textbf{Establish the next claim.}
    One-stage detection proves that \(U\to X\) is formally
    \(\acute{e}\)tale.
    \item \textbf{Fix the indicated data.}
    Let \(\{U_i\to X\}_{i\in I}\) be a finite spectral Zariski covering family.
    \item \textbf{Establish the next claim.}
    Again let \(j:V\to T\) be an affine one-stage square-zero substitution.
    \item \textbf{Construct the indicated map.}
    A map
    \[
    V\longrightarrow\coprod_iU_i
    \]
    determines, by
    Lemma~\ref{lem:finite-clopen-decomposition-infinitesimal}, a finite disjoint
    open-and-closed decomposition
    \[
    V\simeq\coprod_iV_i
    \]
    together with maps \(V_i\to U_i\).
    \item \textbf{Establish the next claim.}
    The same lemma transports this through
    \(j\) to a contractibly unique decomposition
    \[
    T\simeq\coprod_iT_i,
    \qquad
    V_i\simeq V\times_TT_i.
    \]
    \item \textbf{Use pullback stability.}
    For each \(i\), the inverse image of \(U_i\) in \(T_i\) is an open subobject
    whose pullback to \(V_i\) is all of \(V_i\).
    \item \textbf{Establish the indicated equivalence.}
    The one-stage site equivalence
    therefore makes that inverse image all of \(T_i\), so \(T_i\to X\) factors
    through \(U_i\).
    \item \textbf{Establish the next claim.}
    These factorizations assemble to a lift
    \[
    T\longrightarrow\coprod_iU_i.
    \]
    \item \textbf{Establish the next claim.}
    The transported clopen decomposition and the component factorizations have
    contractible spaces of choices, so the total filler space is contractible.
    \item \textbf{Construct the indicated map.}
    One-stage detection proves that the covering morphism is formally
    \(\acute{e}\)tale.
\end{enumerate}
\end{proof}

\begin{theorem}[Descent of effective de Rham quotients]
\label{thm:effective-formally-etale-quotient-descent}
Let \(u:U\twoheadrightarrow X\) be an effective epimorphism over \(S\) which is
formally \(\acute{e}\)tale relative to \(X\), and let \(U_\bullet\to X\) be its
\v{C}ech nerve. Then
\[
  Q_{U/S}\longrightarrow Q_{X/S}
\]
is an effective epimorphism, the simplicial object \(Q_{U_\bullet/S}\) is its
\v{C}ech nerve, and
\[
  Q_{X/S}
  \simeq
  |Q_{U_\bullet/S}|.
\]
For each de Rham relation degree \(n\), the simplicial object in the cover degree
\[
  k\longmapsto \mathcal R^{\mathrm{dR}}_{U_k/S,n}
\]
is the \v{C}ech nerve of the effective epimorphism
\[
  U\times_X\mathcal R^{\mathrm{dR}}_{X/S,n}
  \longrightarrow
  \mathcal R^{\mathrm{dR}}_{X/S,n}.
\]
\end{theorem}

\begin{proof}
\leavevmode
\begin{enumerate}[label=\textup{(\arabic*)}]
    \item \textbf{Establish the next claim.}
    Formal-etale Cartesianity and
    Theorem~\ref{thm:formal-etale-neighbourhood-invariance} give
    \[
    U
    \simeq
    X\times_{Q_{X/S}}Q_{U/S}.
    \]
    \item \textbf{Construct the indicated map.}
    The composite \(U\to X\to Q_{X/S}\) is an effective epimorphism and factors through
    \(Q_{U/S}\to Q_{X/S}\).
    \item \textbf{Construct the indicated map.}
    Factor the latter map as
    \[
    Q_{U/S}\twoheadrightarrow I\lhook\joinrel\longrightarrow Q_{X/S}.
    \]
    \item \textbf{Construct the indicated map.}
    Then the effective epimorphism \(U\to Q_{X/S}\) factors through the monomorphism
    \(I\hookrightarrow Q_{X/S}\).
    \item \textbf{Apply orthogonality.}
    Orthogonality of effective epimorphisms and
    monomorphisms gives a section of this monomorphism, hence it is an equivalence.
    \item \textbf{Conclude the stated claim.}
    Therefore \(Q_{U/S}\to Q_{X/S}\) is an effective epimorphism.
    \item \textbf{Identify the stated objects.}
    Applying the same Cartesian
    formula to the iterated fibre products \(U_k\) identifies
    \[
    Q_{U_k/S}
    \simeq
    \underbrace{
    Q_{U/S}\times_{Q_{X/S}}\cdots\times_{Q_{X/S}}Q_{U/S}
    }_{k+1\text{ factors}}.
    \]
    \item \textbf{Establish the next claim.}
    This is the \v{C}ech nerve statement.
    \item \textbf{Establish the next claim.}
    The de Rham relation formula follows from
    \[
    \mathcal R^{\mathrm{dR}}_{U_k/S,n}
    \simeq
    U_k\times_X\mathcal R^{\mathrm{dR}}_{X/S,n}.
    \]
\end{enumerate}
\end{proof}

\begin{theorem}[Effective descent for the intrinsic cochain tower]
\label{thm:effective-descent-cochain-tower-expanded}
Under the hypotheses of
Theorem~\ref{thm:effective-formally-etale-quotient-descent}, for every de Rham relation degree
\(n\),
\[
  C^n_{\mathrm{dR}}(X/S)
  \xrightarrow{\simeq}
  \operatorname*{Tot}_{[k]\in\Delta}
  C^n_{\mathrm{dR}}(U_k/S).
\]
These equivalences assemble into an equivalence of cosimplicial commutative algebras.
Consequently
\[
  \mathrm{DR}_{X/S}
  \simeq
  \operatorname*{Tot}_k\mathrm{DR}_{U_k/S},
\]
\[
  \Omega^q_{\mathrm{dR}}(X/S)
  \simeq
  \operatorname*{Tot}_k\Omega^q_{\mathrm{dR}}(U_k/S),
\]
and
\[
  F^q_{\mathrm{Tot}}\mathrm{DR}_{X/S}
  \simeq
  \operatorname*{Tot}_k
  F^q_{\mathrm{Tot}}\mathrm{DR}_{U_k/S}.
\]
The same holds for partial totalizations, closed de Rham forms, exact-form
lifts, and the normalized Alexander--Whitney products.
\end{theorem}

\begin{proof}
\leavevmode
\begin{enumerate}[label=\textup{(\arabic*)}]
    \item \textbf{Identify the stated objects.}
    For fixed \(n\), Theorem~\ref{thm:effective-formally-etale-quotient-descent}
    identifies \(\mathcal R^{\mathrm{dR}}_{U_k/S,n}\) as the \v{C}ech nerve of an
    effective epimorphism over \(\mathcal R^{\mathrm{dR}}_{X/S,n}\).
    \item \textbf{Apply effective descent.}
    Effective descent
    for parameterized spectra applied to
    \(a_n^*\mathcal O_S^{\mathrm{add}}\) gives, with
    \[
    c_{k,n}:\mathcal R^{\mathrm{dR}}_{U_k/S,n}
    \longrightarrow
    \mathcal R^{\mathrm{dR}}_{X/S,n}
    \]
    the cover-stage structure map,
    \[
    a_n^*\mathcal O_S^{\mathrm{add}}
    \simeq
    \operatorname*{Tot}_k
    \Pi_{c_{k,n}}c_{k,n}^*a_n^*\mathcal O_S^{\mathrm{add}}.
    \]
    \item \textbf{Apply the indicated result.}
    Apply \(\Pi_{a_n}\), which preserves limits, and compose dependent products.
    \item \textbf{Use the stated construction.}
    This
    gives the degreewise cochain equivalence.
    \item \textbf{Apply naturality.}
    Naturality in \(n\) makes it
    cosimplicial and multiplicative.
    \item \textbf{Identify the relevant fibre.}
    All remaining constructions are obtained by
    limits, fibres, and shifts, which commute with the cover totalization.
\end{enumerate}
\end{proof}

\begin{theorem}[Effective descent for full de Rham functions]
\label{thm:effective-descent-full-de-Rham}
Under the same hypotheses,
\[
  \mathrm{DR}^{\mathrm{full}}_{X/S}
  \xrightarrow{\simeq}
  \operatorname*{Tot}_k
  \mathrm{DR}^{\mathrm{full}}_{U_k/S}
\]
as commutative algebras, and the descended full-to-effective comparison is
\(\rho_{X/S}\).
\end{theorem}

\begin{proof}
\leavevmode
\begin{enumerate}[label=\textup{(\arabic*)}]
    \item \textbf{Apply the cited theorem.}
    Theorem~\ref{thm:de-Rham-preserves-effective-epis} makes
    \(U_{\mathrm{dR}/S}\to X_{\mathrm{dR}/S}\) an effective epimorphism, while left exactness
    identifies the de Rham objects of the \v{C}ech stages with its \v{C}ech nerve.
    \item \textbf{Apply the indicated result.}
    Apply effective descent for parameterized spectra to
    \(p_{\mathrm{dR}}^*\mathcal O_S^{\mathrm{add}}\) and then apply
    \(\Pi_{p_{\mathrm{dR}}}\).
    \item \textbf{Apply naturality.}
    Naturality of restriction along
    \(Q_{-/S}\hookrightarrow(-)_{\mathrm{dR}/S}\) gives compatibility with \(\rho\).
\end{enumerate}
\end{proof}

\subsection{Affine descent and comprehension}

With effective formal \(\acute{e}\)tale descent available, the remaining
argument is representability inside the slice. We first form the stable-section
space-valued functors and verify Zariski descent on affine tests. The
affine-point presentation then constructs the objects representing their
comprehended semantics on arbitrary contexts.

\begin{definition}[Stable-section form-space functors]
\label{def:intrinsic-form-space-functors}
For an arbitrary object \(t:T\to S\), define
\[
  \mathfrak{Form}^{q,\mathrm{stsec}}_{\mathrm{dR},S}(T)
  :=
  \Omega^\infty
  \Gamma^{\mathrm{st}}
  \bigl(S,\Omega^q_{\mathrm{dR}}(T/S)\bigr),
\]
and
\[
  \mathfrak{Form}^{q,\mathrm{closed},\mathrm{stsec}}_{\mathrm{dR},S}(T)
  :=
  \Omega^\infty
  \Gamma^{\mathrm{st}}
  \bigl(S,\Omega^{q,\mathrm{closed}}_{\mathrm{dR}}(T/S;0)\bigr).
\]
Contravariant functoriality in \(T\) makes these presheaves of spaces on
\((\mathcal H_{\Sp})_{/S}\). Write
\[
  j_S:\operatorname{AffPt}(S)\longrightarrow(\mathcal H_{\Sp})_{/S}
\]
for the affine-point inclusion.
\end{definition}

\begin{proposition}[Zariski descent of stable-section form spaces on affine tests]
\label{prop:affine-test-descent-form-moduli}
The restrictions
\[
  j_S^*\mathfrak{Form}^{q,\mathrm{stsec}}_{\mathrm{dR},S},
  \qquad
  j_S^*\mathfrak{Form}^{q,\mathrm{closed},\mathrm{stsec}}_{\mathrm{dR},S}
\]
are spectral Zariski sheaves on \(\operatorname{AffPt}(S)\). If
\(\{T_i\to T\}\) is a finite affine spectral Zariski cover with \v{C}ech nerve
\(T_\bullet\) formed in the slice, then
\[
  \mathfrak{Form}^{q,\mathrm{stsec}}_{\mathrm{dR},S}(T)
  \xrightarrow{\simeq}
  \operatorname*{Tot}_{[k]\in\Delta}
  \mathfrak{Form}^{q,\mathrm{stsec}}_{\mathrm{dR},S}(T_k),
\]
and similarly for closed forms.
\end{proposition}

\begin{proof}
\leavevmode
\begin{enumerate}[label=\textup{(\arabic*)}]
    \item \textbf{Construct the indicated map.}
    The covering morphism
    \[
    \coprod_iT_i\longrightarrow T
    \]
    is an effective epimorphism and is formally \(\acute{e}\)tale by
    Lemma~\ref{lem:zariski-open-formally-etale}.
    \item \textbf{Apply the indicated result.}
    Apply
    Theorem~\ref{thm:effective-descent-cochain-tower-expanded} to its \v{C}ech
    nerve.
    \item \textbf{Establish the next claim.}
    This reconstructs the normalized and closed coefficient objects from
    the full \v{C}ech diagram, including its nonaffine iterated fibre products.
    \item \textbf{Construct the indicated map.}
    Stable global sections are mapping spectra out of the tensor unit and therefore
    preserve limits, and \(\Omega^\infty\) preserves limits.
    \item \textbf{Establish the indicated equivalence.}
    The displayed
    equivalence follows.
    \item \textbf{Establish the next claim.}
    On the affine-point site this is precisely the spectral
    Zariski sheaf condition.
\end{enumerate}
\end{proof}

\begin{construction}[Affine-descent comprehension objects]
\label{con:affine-descent-form-comprehension}
By Proposition~\ref{prop:affine-point-slice-site}, the two affine-point sheaves of
Proposition~\ref{prop:affine-test-descent-form-moduli} determine canonical objects
\[
  \mathbf{Form}^q_{\mathrm{dR},S},
  \qquad
  \mathbf{Form}^{q,\mathrm{closed}}_{\mathrm{dR},S}
  \in
  (\mathcal H_{\Sp})_{/S}.
\]
They are characterized on every affine point \(U\to S\) by
\[
  \Map_S(U,\mathbf{Form}^q_{\mathrm{dR},S})
  \simeq
  \mathfrak{Form}^{q,\mathrm{stsec}}_{\mathrm{dR},S}(U),
\]
and
\[
  \Map_S(U,\mathbf{Form}^{q,\mathrm{closed}}_{\mathrm{dR},S})
  \simeq
  \mathfrak{Form}^{q,\mathrm{closed},\mathrm{stsec}}_{\mathrm{dR},S}(U).
\]
For an arbitrary \(T\to S\), define the \emph{comprehended de Rham form
judgments} by
\[
  \mathfrak{Form}^{q,\mathrm{comp}}_{\mathrm{dR},S}(T)
  :=
  \Map_S(T,\mathbf{Form}^q_{\mathrm{dR},S}),
\]
and
\[
  \mathfrak{Form}^{q,\mathrm{closed},\mathrm{comp}}_{\mathrm{dR},S}(T)
  :=
  \Map_S(T,\mathbf{Form}^{q,\mathrm{closed}}_{\mathrm{dR},S}).
\]
Thus the comprehended semantics on arbitrary contexts is the affine-descent
extension of the stable-section semantics on affine tests.
\end{construction}

\begin{proposition}[Right-Kan description of comprehended form judgments]
\label{prop:form-comprehension-right-kan}
There are canonical equivalences of presheaves on the full slice
\[
  \mathfrak{Form}^{q,\mathrm{comp}}_{\mathrm{dR},S}
  \simeq
  \operatorname{Ran}_{j_S^{\op}}
  \bigl(
    j_S^*\mathfrak{Form}^{q,\mathrm{stsec}}_{\mathrm{dR},S}
  \bigr),
\]
and
\[
  \mathfrak{Form}^{q,\mathrm{closed},\mathrm{comp}}_{\mathrm{dR},S}
  \simeq
  \operatorname{Ran}_{j_S^{\op}}
  \bigl(
    j_S^*
    \mathfrak{Form}^{q,\mathrm{closed},\mathrm{stsec}}_{\mathrm{dR},S}
  \bigr).
\]
Equivalently, for every \(T\to S\),
\[
\begin{aligned}
  \mathfrak{Form}^{q,\mathrm{comp}}_{\mathrm{dR},S}(T)
  &\simeq
  \operatorname*{lim}_{(U\to T)\in\operatorname{AffPt}(T)^{\op}}
  \mathfrak{Form}^{q,\mathrm{stsec}}_{\mathrm{dR},S}(U),\\
  \mathfrak{Form}^{q,\mathrm{closed},\mathrm{comp}}_{\mathrm{dR},S}(T)
  &\simeq
  \operatorname*{lim}_{(U\to T)\in\operatorname{AffPt}(T)^{\op}}
  \mathfrak{Form}^{q,\mathrm{closed},\mathrm{stsec}}_{\mathrm{dR},S}(U).
\end{aligned}
\]
\end{proposition}

\begin{proof}
\leavevmode
\begin{enumerate}[label=\textup{(\arabic*)}]
    \item \textbf{Apply the cited lemma.}
    Lemma~\ref{lem:de-Rham-affine-density} gives
    \[
    T
    \simeq
    \operatorname*{colim}_{(U\to T)\in\operatorname{AffPt}(T)}U
    \]
    in the slice over \(S\).
    \item \textbf{Use the stated construction.}
    Mapping this colimit into either comprehension object
    gives the corresponding limit of affine mapping spaces.
    \item \textbf{Identify the stated objects.}
    The affine characterizing equivalences of
    Construction~\ref{con:affine-descent-form-comprehension} identify these mapping
    spaces with the stable-section affine form spaces.
    \item \textbf{Establish the next claim.}
    The resulting limit is the
    pointwise formula for right Kan extension along \(j_S^{\op}\).
\end{enumerate}
\end{proof}

\begin{construction}[Stable-section-to-comprehended comparison]
\label{con:intrinsic-comprehended-form-comparison}
For every \(T\to S\), restriction of a stable-section form along every affine point
\(U\to T\) gives a compatible cone
\[
  \mathfrak{Form}^{q,\mathrm{stsec}}_{\mathrm{dR},S}(T)
  \longrightarrow
  \mathfrak{Form}^{q,\mathrm{stsec}}_{\mathrm{dR},S}(U).
\]
Using Proposition~\ref{prop:form-comprehension-right-kan}, this cone determines a
canonical natural transformation
\[
  \mathrm{cmp}^q:
  \mathfrak{Form}^{q,\mathrm{stsec}}_{\mathrm{dR},S}
  \longrightarrow
  \mathfrak{Form}^{q,\mathrm{comp}}_{\mathrm{dR},S}.
\]
The same construction gives
\[
  \mathrm{cmp}^{q,\mathrm{closed}}:
  \mathfrak{Form}^{q,\mathrm{closed},\mathrm{stsec}}_{\mathrm{dR},S}
  \longrightarrow
  \mathfrak{Form}^{q,\mathrm{closed},\mathrm{comp}}_{\mathrm{dR},S}.
\]
If \(T\to S\) is affine, the identity affine point \(T\to T\) is terminal in
\(\operatorname{AffPt}(T)\), so it is initial after passage to the opposite
category. The preceding limit therefore evaluates at \(T\), and
\[
  \mathrm{cmp}_T^q,
  \qquad
  \mathrm{cmp}_T^{q,\mathrm{closed}}
\]
are equivalences.
\end{construction}

\begin{theorem}[Comprehension of affine-descended de Rham form judgments]
\label{thm:internal-moduli-forms-detailed}
For every \(q\geq0\), the objects
\[
  \mathbf{Form}^q_{\mathrm{dR},S},
  \qquad
  \mathbf{Form}^{q,\mathrm{closed}}_{\mathrm{dR},S}
\]
comprehend the affine-descended de Rham form judgments:
\[
  \Map_S(T,\mathbf{Form}^q_{\mathrm{dR},S})
  \simeq
  \mathfrak{Form}^{q,\mathrm{comp}}_{\mathrm{dR},S}(T),
\]
\[
  \Map_S(T,\mathbf{Form}^{q,\mathrm{closed}}_{\mathrm{dR},S})
  \simeq
  \mathfrak{Form}^{q,\mathrm{closed},\mathrm{comp}}_{\mathrm{dR},S}(T)
\]
for every \(T\to S\), while on affine contexts these spaces are canonically the
stable-section form spaces. Their identity sections are the universal
terms for the comprehended judgments.

The underlying-form transformation induces
\[
  u_q:
  \mathbf{Form}^{q,\mathrm{closed}}_{\mathrm{dR},S}
  \longrightarrow
  \mathbf{Form}^q_{\mathrm{dR},S}.
\]
For \(q\geq1\), the exact-form lift induces
\[
  d^{\mathrm{closed}}_{\mathrm{dR}}:
  \mathbf{Form}^{q-1}_{\mathrm{dR},S}
  \longrightarrow
  \mathbf{Form}^{q,\mathrm{closed}}_{\mathrm{dR},S}.
\]
Their composite is the morphism of comprehension objects whose restriction to every
affine test represents the intrinsic de Rham differential.
\end{theorem}

\begin{proof}
\leavevmode
\begin{enumerate}[label=\textup{(\arabic*)}]
    \item \textbf{Use the stated construction.}
    Construction~\ref{con:affine-descent-form-comprehension} gives the comprehension
    objects and their universal terms.
    \item \textbf{Apply the cited proposition.}
    Proposition~\ref{prop:form-comprehension-right-kan}
    and Construction~\ref{con:intrinsic-comprehended-form-comparison} supply the
    right-Kan formula and the affine comparison.
    \item \textbf{Construct the indicated map.}
    The underlying-form maps and exact-form lifts are natural in affine substitution,
    hence define morphisms of affine-point sheaves.
    \item \textbf{Construct the indicated map.}
    Under
    Proposition~\ref{prop:affine-point-slice-site}, these natural transformations
    correspond to the displayed morphisms in the slice.
    \item \textbf{Apply the cited proposition.}
    Proposition~\ref{prop:exact-classifier-boundary} identifies the composite on
    every affine test with the homotopy-coherent de Rham differential.
    \item \textbf{Use the stated construction.}
    Affine
    descent gives the global statement.
\end{enumerate}
\end{proof}


We now have both semantic levels. On affine contexts the comprehended
semantics agrees with the stable-section semantics; on arbitrary contexts it
is the canonical affine-descent extension. The underlying-form and exact-form
maps descend with the same construction. The remaining sections can therefore
turn to comparison questions without reopening descent or representability.
\section{First-order cotangent comparison}
\label{sec:de-Rham-comparison-reflections}
\label{sec:first-order-linearization}

The de Rham calculus constructed above retains every finite infinitesimal stage
visible to the generated congruence. To compare it with the differentiated
spectral doctrine, we now restrict that nonlinear relation to the first
structured infinitesimal neighbourhood of the diagonal. This comparison is
deliberately one-way: first-order restriction extracts the cotangent
coefficient, but the full degree-one de Rham normalization can contain
higher-order infinitesimal classes that disappear under that restriction.

For a morphism of spectral schemes \(p:X\to S\), the relevant geometric
comparison begins with the first structured infinitesimal neighbourhood:
\[
\begin{tikzcd}[column sep=huge,row sep=large]
  X \arrow[r,hook,"i_1"] \arrow[dr,"\Delta_{X/S}"']
  & D^{(1)}_{X/S} \arrow[d]
  \\
  & X\times_SX,
\end{tikzcd}
\qquad
\text{first-order functions}
\longrightarrow
\operatorname{Real}^{\mathrm{st}}_X(L_{X/S}).
\]
The coefficient-level target of this restriction is the already constructed
stable realization of the cotangent complex:
\[
  \Omega^1_{\mathrm{dR}}(X/S)
  \longrightarrow
  \Pi_p\operatorname{Real}^{\mathrm{st}}_X(L_{X/S}),
\]

This target becomes meaningful through the actual restriction map. Restriction
of scalar functions along
\(D^{(1)}_{X/S}\to\mathcal R^{\mathrm{dR}}_{X/S,1}\) fits into a fibrewise
comparison:
\[
\begin{tikzcd}[column sep=large,row sep=large]
  C^{1,\mathrm{ctr}}_{\mathrm{dR}}(X/S)
      \arrow[r]
      \arrow[d,"\sigma"']
  &
  \mathscr F^{(1)}_{X/S}
      \arrow[d,"r_0"]
  \\
  \mathcal O_X^{\mathrm{add}}
      \arrow[r,equal]
  &
  \mathcal O_X^{\mathrm{add}}
  \\
  \operatorname{fib}(\sigma)
      \arrow[u]
      \arrow[r]
      \arrow[rr,bend right=18,"\mathscr\lambda^1_{X/S}"']
  &
  K^{(1)}_{X/S}
      \arrow[u]
      \arrow[r,"\simeq"']
  &
  \operatorname{Real}^{\mathrm{st}}_X(L_{X/S}) .
\end{tikzcd}
\]
After dependent product to \(S\), the bottom row becomes the first-order
linearization of intrinsic one-forms, and the degree-zero differential maps to
the stabilized universal derivation.

This section is restricted to spectral schemes so that
\(L_{X/S}\in\QCoh(X)\) is available. The first structured infinitesimal
neighbourhood, its base-change behaviour, and the identification of the
difference of its two retractions with the universal relative derivation were
constructed in Chapter~\Ref{chap:differentiation-spectral-doctrine}. We use
those results here as the first-order comparison interface.
Write
\[
  X\xrightarrow{i_1}D^{(1)}_{X/S}
  \xrightarrow{\nu_1}X\times_SX
\]
for that first structured infinitesimal neighbourhood, and
\[
  h_0:=\operatorname{pr}_1\nu_1,
  \qquad
  h_1:=\operatorname{pr}_2\nu_1
  :D^{(1)}_{X/S}\rightrightarrows X
\]
for its two retractions.

\begin{proposition}[First structured infinitesimal neighbourhood maps to the de Rham relation]
\label{prop:first-neighbourhood-relation}
There is a canonical map over \(X\times_SX\)
\[
  \phi_1:
  D^{(1)}_{X/S}
  \longrightarrow
  \mathcal R^{\mathrm{dR}}_{X/S,1}
\]
whose composites with the source and target projections are \(h_0\) and
\(h_1\). It is canonical and compatible with arbitrary Cartesian base change of
spectral schemes.
\end{proposition}

\begin{proof}
\leavevmode
\begin{enumerate}[label=\textup{(\arabic*)}]
    \item \textbf{Establish the indicated equivalence.}
    The structured immersion
    \(i_1:X\to D^{(1)}_{X/S}\) is a square-zero infinitesimal substitution, hence a
    de Rham equivalence over \(S\) by
    Proposition~\ref{prop:global-square-zero-de-Rham}.
    \item \textbf{Apply the indicated result.}
    Apply
    Proposition~\ref{prop:de-Rham-neighbourhood-factorization} to
    \[
    X\longrightarrow D^{(1)}_{X/S}
    \longrightarrow X\times_SX\longrightarrow S.
    \]
    \item \textbf{Establish the next claim.}
    The resulting target de Rham neighbourhood is
    \(\operatorname{Nbh}^{\mathrm{dR}}_{X/S}(X\times_SX)\), which is canonically
    \(\mathcal R^{\mathrm{dR}}_{X/S,1}\) by
    Theorem~\ref{thm:simplices-as-neighbourhoods}.
    \item \textbf{Construct the indicated map.}
    The two projection identities follow
    from the defining fibre-product map.
    \item \textbf{Apply naturality.}
    Naturality is inherited from
    the first-neighbourhood base-change theorem and the de Rham neighbourhood construction.
\end{enumerate}
\end{proof}

\subsection{The first-order function coefficient and realized cotangent complex}

\begin{construction}[First-order function coefficient]
\label{con:stable-cotangent-coefficient}
Put
\[
  D:=D^{(1)}_{X/S}.
\]
By base change of the additive scalar,
\[
  h_0^*\mathcal O_X^{\mathrm{add}}
  \simeq
  \mathcal O_D^{\mathrm{add}}.
\]
Form the first-order function object over the centre
\[
  \mathscr F^{(1)}_{X/S}
  :=
  \Pi_{h_0}\mathcal O_D^{\mathrm{add}}
  \simeq
  \Pi_{h_0}h_0^*\mathcal O_X^{\mathrm{add}}
  \in\mathcal D_{\Sp/X}.
\]
Evaluation along the section \(i_1:X\to D\) gives
\[
  r_0:
  \mathscr F^{(1)}_{X/S}
  \longrightarrow
  \mathcal O_X^{\mathrm{add}},
\]
while the unit of \(h_0^*\dashv\Pi_{h_0}\) gives the constant-function section
\[
  s_0:
  \mathcal O_X^{\mathrm{add}}
  \longrightarrow
  \mathscr F^{(1)}_{X/S}.
\]
Since \(h_0i_1\simeq\id_X\), the adjunction identities give
\[
  r_0s_0\simeq\id.
\]
Define the \emph{first-order function coefficient}
\[
  K^{(1)}_{X/S}
  :=
  \operatorname{fib}(r_0).
\]
Thus the splitting determined by \(s_0\) gives a canonical stable equivalence
\[
  \mathscr F^{(1)}_{X/S}
  \simeq
  \mathcal O_X^{\mathrm{add}}
  \oplus
  K^{(1)}_{X/S}.
\]
\end{construction}

\begin{proposition}[First-order coefficient and the realized cotangent complex]
\label{prop:stable-cotangent-realization}
There is a canonical equivalence in \(\mathcal D_{\Sp/X}\)
\[
  \xi_{X/S}:
  K^{(1)}_{X/S}
  \xrightarrow{\simeq}
  \operatorname{Real}^{\mathrm{st}}_X(L_{X/S}),
\]
compatible with Cartesian base change of spectral schemes. Consequently, for every affine spectral point \(a:T\to X\),
\[
  \Gamma^{\mathrm{st}}
  \bigl(T,a^*K^{(1)}_{X/S}\bigr)
  \xrightarrow{\simeq}
  \Gamma(T,a^*L_{X/S}),
\]
and
\[
  \Gamma^{\mathrm{st}}
  \bigl(S,\Pi_pK^{(1)}_{X/S}\bigr)
  \xrightarrow{\simeq}
  \Gamma(X,L_{X/S}).
\]
\end{proposition}

\begin{proof}
\leavevmode
\begin{enumerate}[label=\textup{(\arabic*)}]
    \item \textbf{Fix the indicated data.}
    Let \(a:T\to X\) be an affine spectral point and form
    \[
    \begin{tikzcd}
    D_T:=T\times_{X,h_0}D
    \ar[r,"\widetilde a"]
    \ar[d,"h_{0,T}"']
    & D \ar[d,"h_0"]\\
    T \ar[r,"a"'] & X.
    \end{tikzcd}
    \]
    \item \textbf{Establish the next claim.}
    Right Beck--Chevalley and base change of the additive scalar give
    \[
    a^*\mathscr F^{(1)}_{X/S}
    \simeq
    \Pi_{h_{0,T}}\mathcal O_{D_T}^{\mathrm{add}}.
    \]
    \item \textbf{Establish the next claim.}
    The retraction identity \(h_0i_1\simeq\id_X\) determines a canonical centre
    \[
    i_T:T\longrightarrow D_T,
    \qquad
    i_T=(\id_T,i_1a).
    \]
    \item \textbf{Establish the next claim.}
    The square
    \[
    \begin{tikzcd}
    T \ar[r,"i_T"] \ar[d,"a"']
    & D_T \ar[d,"\widetilde a"]\\
    X \ar[r,"i_1"']
    & D
    \end{tikzcd}
    \]
    is Cartesian.
    \item \textbf{Establish the next claim.}
    Indeed,
    \[
    \begin{aligned}
    D_T\times_DX
    &\simeq
    (T\times_XD)\times_DX\\
    &\simeq
    T\times_XX\\
    &\simeq T,
    \end{aligned}
    \]
    where the middle fibre product uses \(h_0i_1\simeq\id_X\).
    \item \textbf{Conclude the stated claim.}
    Thus this is precisely geometric Cartesian base change of the structured
    square-zero centre \(i_1:X\to D\).
    \item \textbf{Invoke the cited result.}
    By
    \Ref{thm:structured-square-zero-geometric-base-change}, \(i_T\) is therefore
    canonically a structured square-zero extension with coefficient
    \[
    a^*L_{X/S},
    \]
    as supplied by the structured base-change construction.
    \item \textbf{Use the stated hypothesis.}
    Since \(T\) is affine and \(i_T\) is a square-zero infinitesimal substitution,
    \Ref{lem:finite-infinitesimal-affineness-invariance} makes \(D_T\) affine.
    Write
    \[
    T=\Spec A,
    \qquad
    D_T=\Spec A'.
    \]
    \item \textbf{Construct the indicated map.}
    The opposite coordinate map
    \[
    A'\longrightarrow A
    \]
    is a structured spectral square-zero extension with coefficient
    \(\Gamma(T,a^*L_{X/S})\).
    \item \textbf{Invoke the cited result.}
    By
    \Ref{prop:structured-square-zero-fibre}, its homotopy fibre is canonically
    \[
    \Gamma(T,a^*L_{X/S}).
    \]
    \item \textbf{Establish the next claim.}
    Using the adjunction \(h_{0,T}^*\dashv\Pi_{h_{0,T}}\) and
    Proposition~\ref{prop:stable-scalar-sections-represented},
    \[
    \begin{aligned}
    \Gamma^{\mathrm{st}}
    \bigl(T,a^*\mathscr F^{(1)}_{X/S}\bigr)
    &\simeq
    \Gamma^{\mathrm{st}}
    \bigl(D_T,\mathcal O_{D_T}^{\mathrm{add}}\bigr)\\
    &\simeq A',
    \end{aligned}
    \]
    while
    \[
    \Gamma^{\mathrm{st}}(T,\mathcal O_T^{\mathrm{add}})
    \simeq A.
    \]
    \item \textbf{Identify the stated objects.}
    Under these identifications, the pulled-back evaluation map \(a^*r_0\) is exactly
    the square-zero augmentation \(A'\to A\).
    \item \textbf{Identify the relevant fibre.}
    Stable sections preserve fibres, so
    \[
    \begin{aligned}
    \Gamma^{\mathrm{st}}
    \bigl(T,a^*K^{(1)}_{X/S}\bigr)
    &\simeq
    \operatorname{fib}(A'\to A)\\
    &\simeq
    \Gamma(T,a^*L_{X/S}).
    \end{aligned}
    \]
    \item \textbf{Establish the next claim.}
    On the other hand, by the defining affine-point formula for stable realization,
    \[
    \Gamma^{\mathrm{st}}
    \bigl(T,a^*\operatorname{Real}^{\mathrm{st}}_X(L_{X/S})\bigr)
    \simeq
    \Gamma(T,a^*L_{X/S}).
    \]
    \item \textbf{Construct the indicated map.}
    Every comparison above is functorial under morphisms of affine points.
    \item \textbf{Conclude the stated claim.}
    Hence the
    stable affine-point presentation of
    Proposition~\ref{prop:stable-affine-point-presentation} identifies the two sheaves
    of spectra and produces the canonical equivalence
    \[
    \xi_{X/S}:
    K^{(1)}_{X/S}
    \xrightarrow{\simeq}
    \operatorname{Real}^{\mathrm{st}}_X(L_{X/S}).
    \]
    \item \textbf{Establish the next claim.}
    Geometric base change of the first structured infinitesimal neighbourhood, cotangent base change, right
    Beck--Chevalley, and base change of \(\operatorname{Real}^{\mathrm{st}}\) show that \(\xi_{X/S}\) is
    compatible with Cartesian base change of spectral schemes.
    \item \textbf{Establish the next claim.}
    The affine-point formula is the first displayed consequence.
    \item \textbf{Establish the next claim.}
    For the second,
    apply \(\Pi_p\) and stable global sections:
    \[
    \begin{aligned}
    \Gamma^{\mathrm{st}}
    \bigl(S,\Pi_pK^{(1)}_{X/S}\bigr)
    &\simeq
    \Gamma^{\mathrm{st}}
    \bigl(X,K^{(1)}_{X/S}\bigr)\\
    &\simeq
    \Gamma^{\mathrm{st}}
    \bigl(X,\operatorname{Real}^{\mathrm{st}}_X(L_{X/S})\bigr)\\
    &\simeq
    \Gamma(X,L_{X/S}),
    \end{aligned}
    \]
    where the first equivalence is \(p^*\dashv\Pi_p\) and the last is
    Proposition~\ref{prop:stable-realization-qcoh}.
\end{enumerate}
\end{proof}

\begin{construction}[Stabilized universal derivation over the base]
\label{con:stabilized-universal-derivation}
The stabilized universal derivation over the base is obtained from the additive
difference of the two first-neighbourhood retractions after dependent product to
\(S\).

Put
\[
  k:=ph_0\simeq ph_1:D\longrightarrow S
\]
and
\[
  \mathscr E_D
  :=
  \Pi_k\mathcal O_D^{\mathrm{add}}
  \in\mathcal D_{\Sp/S}.
\]
Composition of dependent products gives canonical equivalences
\[
  \mathscr E_D
  \simeq
  \Pi_p\Pi_{h_i}\mathcal O_D^{\mathrm{add}},
  \qquad i=0,1.
\]
For each \(i\), the unit of \(h_i^*\dashv\Pi_{h_i}\), followed by \(\Pi_p\),
gives
\[
  \overline s_i:
  \Pi_p\mathcal O_X^{\mathrm{add}}
  \longrightarrow
  \mathscr E_D.
\]
Restriction along \(i_1:X\to D\) gives
\[
  \overline r:
  \mathscr E_D
  \longrightarrow
  \Pi_p\mathcal O_X^{\mathrm{add}}.
\]
Since \(h_ii_1\simeq\id_X\), functoriality of the adjunction units gives canonical
homotopies
\[
  \overline r\,\overline s_i
  \simeq
  \id_{\Pi_p\mathcal O_X^{\mathrm{add}}},
  \qquad i=0,1.
\]
Using the \(h_0\)-description of \(\mathscr E_D\), exactness of \(\Pi_p\), and
Proposition~\ref{prop:stable-cotangent-realization}, we obtain canonical
equivalences
\[
  \operatorname{fib}(\overline r)
  \simeq
  \Pi_pK^{(1)}_{X/S}
  \xrightarrow{\simeq}
  \Pi_p\operatorname{Real}^{\mathrm{st}}_X(L_{X/S}).
\]
Therefore the additive difference
\[
  \overline s_1-\overline s_0:
  \Pi_p\mathcal O_X^{\mathrm{add}}
  \longrightarrow
  \mathscr E_D
\]
has canonically null composite with \(\overline r\) and hence determines, through
the stable fibre universal property, a canonical morphism
\[
  d^{\mathrm{univ,st}}_{X/S}:
  \Pi_p\mathcal O_X^{\mathrm{add}}
  \longrightarrow
  \Pi_p\operatorname{Real}^{\mathrm{st}}_X(L_{X/S}).
\]

\end{construction}

\begin{proposition}[Affine identification of the stabilized derivation]
\label{prop:stabilized-derivation-affine}
Suppose \(p:X\to S\) is affine, with
\[
  S=\Spec R,
  \qquad
  X=\Spec A.
\]
Then the map on stable global sections induced by
\(d^{\mathrm{univ,st}}_{X/S}\) is the universal spectral derivation
\[
  d_{A/R}:A\longrightarrow L_{A/R}.
\]
Thus, after restricting a morphism of spectral schemes to a compatible affine pair, the
corresponding affine construction recovers the universal relative derivation.
\end{proposition}

\begin{proof}
\leavevmode
\begin{enumerate}[label=\textup{(\arabic*)}]
    \item \textbf{Construct the indicated map.}
    For the affine morphism
    \[
    X=\Spec A\longrightarrow S=\Spec R,
    \]
    \item \textbf{Identify the stated objects.}
    Chapter~\Ref{chap:differentiation-spectral-doctrine},
    \Ref{thm:first-neighbourhood-universal-derivation}, identifies the first structured
    infinitesimal neighbourhood after splitting by \(h_0\) with a square-zero algebra
    \[
    A'\simeq A\oplus L_{A/R}
    \]
    for which the two algebra sections corresponding contravariantly to \(h_0\) and
    \(h_1\) are
    \[
    \alpha_0(a)=(a,0),
    \qquad
    \alpha_1(a)=(a,d_{A/R}a).
    \]
    \item \textbf{Identify the stated objects.}
    Under Proposition~\ref{prop:stable-scalar-sections-represented},
    \(\overline s_0\) and \(\overline s_1\) identify with these two sections.
    \item \textbf{Construct the indicated map.}
    Their
    additive difference is
    \[
    a\longmapsto(0,d_{A/R}a).
    \]
    \item \textbf{Use the stated construction.}
    Projection to the evaluation fibre gives \(d_{A/R}\), and
    Proposition~\ref{prop:stable-cotangent-realization} identifies that fibre with the
    stable realization of the cotangent coefficient.
    \item \textbf{Conclude the stated claim.}
    Thus the resulting map on stable
    global sections is exactly the universal spectral derivation.
\end{enumerate}
\end{proof}

\subsection{First-order restriction of intrinsic one-forms}

Use the degree-one centred coefficient
\(C^{1,\mathrm{ctr}}_{\mathrm{dR}}(X/S)\) of
Proposition~\ref{prop:centered-degreewise-etale-base-change}; here
\(v_1=\mathsf s:\mathcal R^{\mathrm{dR}}_{X/S,1}\to X\).
The diagonal \(X\to \mathcal R^{\mathrm{dR}}_{X/S,1}\) induces
\[
  \sigma:
  C^{1,\mathrm{ctr}}_{\mathrm{dR}}(X/S)
  \longrightarrow
  \mathcal O_X^{\mathrm{add}}.
\]
Since the structure map of the de Rham relation is \(p\mathsf s\),
\[
  C^1_{\mathrm{dR}}(X/S)
  \simeq
  \Pi_pC^{1,\mathrm{ctr}}_{\mathrm{dR}}(X/S),
  \qquad
  C^0_{\mathrm{dR}}(X/S)
  \simeq
  \Pi_p\mathcal O_X^{\mathrm{add}}.
\]
Because \(\Pi_p\) preserves fibres,
\[
  \Omega^1_{\mathrm{dR}}(X/S)
  \simeq
  \Pi_p\operatorname{fib}(\sigma).
\]

Restriction along
\(\phi_1:D^{(1)}_{X/S}\to \mathcal R^{\mathrm{dR}}_{X/S,1}\) gives a morphism
\[
  C^{1,\mathrm{ctr}}_{\mathrm{dR}}(X/S)
  \longrightarrow
  \mathscr F^{(1)}_{X/S}.
\]
It intertwines \(\sigma\) with the centre evaluation
\(r_0:\mathscr F^{(1)}\to\mathcal O_X^{\mathrm{add}}\), and hence induces
\[
  \operatorname{fib}(\sigma)
  \longrightarrow
  K^{(1)}_{X/S}
  \xrightarrow{\xi_{X/S}}
  \operatorname{Real}^{\mathrm{st}}_X(L_{X/S}).
\]
Denote this composite by
\[
  \mathscr\lambda^1_{X/S}:
  \operatorname{fib}(\sigma)
  \longrightarrow
  \operatorname{Real}^{\mathrm{st}}_X(L_{X/S}).
\]

\begin{theorem}[First-order linearization]
\label{thm:first-order-linearization}
Applying \(\Pi_p\) to \(\mathscr\lambda^1_{X/S}\) gives a canonical morphism
\[
  \lambda^1_{X/S}:
  \Omega^1_{\mathrm{dR}}(X/S)
  \longrightarrow
  \Pi_p\operatorname{Real}^{\mathrm{st}}_X(L_{X/S}).
\]
It satisfies
\[
  \lambda^1_{X/S}\circ d^0_{\mathrm{dR}}
  \simeq
  d^{\mathrm{univ,st}}_{X/S},
\]
equivalently the triangle
\[
\begin{tikzcd}[column sep=huge,row sep=large]
  \Omega^0_{\mathrm{dR}}(X/S)
    \arrow[r,"d^0_{\mathrm{dR}}"]
    \arrow[dr,"d^{\mathrm{univ,st}}_{X/S}"']
  &
  \Omega^1_{\mathrm{dR}}(X/S)
    \arrow[d,"\lambda^1_{X/S}"]
  \\
  &
  \Pi_p\operatorname{Real}^{\mathrm{st}}_X(L_{X/S})
\end{tikzcd}
\]
commutes up to the canonical homotopy,

where the right-hand side is the base-valued stabilized universal derivation of
Construction~\ref{con:stabilized-universal-derivation}.
\end{theorem}

\begin{proof}
\leavevmode
\begin{enumerate}[label=\textup{(\arabic*)}]
    \item \textbf{Establish the next claim.}
    Under the ordered face convention, the degree-zero boundary is the typed
    scalar-restriction difference
    \[
    \operatorname{res}_{\mathsf t}-\operatorname{res}_{\mathsf s}.
    \]
    \item \textbf{Establish the next claim.}
    Restriction along \(\phi_1\) replaces the target and source projections by \(h_1\) and
    \(h_0\), respectively.
    \item \textbf{Construct the indicated map.}
    After dependent product to \(S\), these two restrictions are precisely
    the maps \(\overline s_1\) and \(\overline s_0\) into the common stable function
    object \(\mathscr E_D\).
    \item \textbf{Identify the relevant fibre.}
    Their difference factors through the evaluation fibre
    \[
    \operatorname{fib}(\overline r)
    \simeq
    \Pi_pK^{(1)}_{X/S}
    \xrightarrow{\simeq}
    \Pi_p\operatorname{Real}^{\mathrm{st}}_X(L_{X/S}),
    \]
    and by Construction~\ref{con:stabilized-universal-derivation} this factor is
    \(d^{\mathrm{univ,st}}_{X/S}\).
    \item \textbf{Construct the indicated map.}
    The map on the same fibre induced by restriction of
    intrinsic one-forms is \(\lambda^1_{X/S}\), giving the displayed identity.
    \item \textbf{Apply the cited proposition.}
    Proposition~\ref{prop:stabilized-derivation-affine} identifies the resulting map on
    compatible affine charts with the universal relative derivation.
\end{enumerate}
\end{proof}

\begin{remark}[Higher-order kernel of first-order linearization]
The morphism \(\lambda^1_{X/S}\) records restriction of normalized
degree-one de Rham functions to the first structured infinitesimal neighbourhood. Higher-order
finite-infinitesimal functions can lie in its kernel, as the quadratic example below
shows.
\end{remark}

\begin{proposition}[A quadratic class beyond first order]
\label{prop:higher-infinitesimal-terms}
Let
\[
  S=\Spec H\mathbb Q,
  \qquad
  X=\Spec H(\mathbb Q[t]).
\]
Then the map
\[
  \pi_0\Gamma^{\mathrm{st}}
  \bigl(S,\Omega^1_{\mathrm{dR}}(X/S)\bigr)
  \longrightarrow
  \pi_0\Gamma(X,L_{X/S})
\]
induced by \(\lambda^1_{X/S}\) has nonzero kernel.
\end{proposition}

\begin{proof}
\leavevmode
\begin{enumerate}[label=\textup{(\arabic*)}]
    \item \textbf{Fix the indicated data.}
    Let
    \[
    t\in
    \pi_0\Gamma^{\mathrm{st}}
    \bigl(S,C^0_{\mathrm{dR}}(X/S)\bigr)
    \]
    be the coordinate function.
    \item \textbf{Introduce the notation.}
    Define its degree-one infinitesimal difference
    \[
    \delta t
    :=
    \operatorname{res}_{\mathsf t}(t)
    -
    \operatorname{res}_{\mathsf s}(t)
    \in
    \pi_0\Gamma^{\mathrm{st}}
    \bigl(S,C^1_{\mathrm{dR}}(X/S)\bigr).
    \]
    \item \textbf{Fix the indicated data.}
    Let
    \[
    \sigma:
    C^1_{\mathrm{dR}}(X/S)\longrightarrow
    C^0_{\mathrm{dR}}(X/S)
    \]
    be the degree-one codegeneracy induced by the diagonal.
    \item \textbf{Establish the next claim.}
    The simplicial
    identities give canonical homotopies
    \[
    \sigma\operatorname{res}_{\mathsf t}\simeq\id
    \simeq
    \sigma\operatorname{res}_{\mathsf s}.
    \]
    \item \textbf{Conclude the stated claim.}
    Hence \(\sigma(\delta t)\) is canonically null.
    \item \textbf{Use the stated construction.}
    Multiplication in the
    commutative algebra \(C^1_{\mathrm{dR}}(X/S)\) therefore gives
    \[
    c_2:=(\delta t)^2
    \]
    together with the induced canonical nullhomotopy
    \[
    \sigma(c_2)\simeq0.
    \]
    \item \textbf{Invoke the cited result.}
    By the stable fibre description
    \[
    \Omega^1_{\mathrm{dR}}(X/S)
    \simeq
    \operatorname{fib}
    \bigl(
    C^1_{\mathrm{dR}}(X/S)\to C^0_{\mathrm{dR}}(X/S)
    \bigr),
    \]
    this pair determines a canonical class
    \[
    \omega_2
    \in
    \pi_0\Gamma^{\mathrm{st}}
    \bigl(S,\Omega^1_{\mathrm{dR}}(X/S)\bigr).
    \]
    \item \textbf{Conclude the stated claim.}
    Thus the normalization nullhomotopy is determined by the displayed fibre
    construction.
    \item \textbf{Introduce the notation.}
    Put
    \[
    Y=X\times_SX
    =
    \Spec H(\mathbb Q[t_0,t_1]),
    \qquad
    u=t_1-t_0,
    \]
    and let
    \[
    T_2
    :=
    \Spec H(\mathbb Q[t,u]/(u^3))
    \]
    map to \(Y\) by
    \[
    t_0\longmapsto t,
    \qquad
    t_1\longmapsto t+u.
    \]
    \item \textbf{Establish the next claim.}
    The centre \(X\to T_2\) factors through
    \[
    \mathbb Q[t,u]/(u^3)
    \longrightarrow
    \mathbb Q[t,u]/(u^2)
    \longrightarrow
    \mathbb Q[t],
    \]
    whose successive kernels \((u^2)/(u^3)\) and \((u)/(u^2)\) are square-zero.
    \item \textbf{Conclude the stated claim.}
    Hence \(X\to T_2\) is a finite spectral infinitesimal substitution, and
    Proposition~\ref{prop:de-Rham-neighbourhood-factorization} supplies
    \[
    T_2\longrightarrow
    \mathcal R^{\mathrm{dR}}_{X/S,1}.
    \]
    \item \textbf{Establish the next claim.}
    Under the represented function-spectrum formula,
    \(\delta t\) restricts to \(u\), so \(c_2\) restricts to
    \[
    u^2\neq0
    \qquad\text{in}\qquad
    \mathbb Q[t,u]/(u^3).
    \]
    \item \textbf{Conclude the stated claim.}
    Therefore the image of \(\omega_2\) in degree-one cochains is nonzero, and
    hence \(\omega_2\neq0\).
    \item \textbf{Establish the next claim.}
    Restrict \(\delta t\) along the canonical first structured infinitesimal neighbourhood
    \[
    D^{(1)}_{X/S}
    \longrightarrow
    \mathcal R^{\mathrm{dR}}_{X/S,1}.
    \]
    \item \textbf{Establish the next claim.}
    Under the splitting by \(h_0\), the difference of the two retractions takes
    values in the square-zero coefficient
    \(L_{X/S}\).
    \item \textbf{Invoke the cited result.}
    By
    \Ref{prop:structured-square-zero-fibre}, the induced nonunital multiplication
    on that coefficient fibre is coherently null.
    \item \textbf{Establish the next claim.}
    Consequently the square
    \[
    (\delta t)^2
    \]
    restricts canonically to zero on \(D^{(1)}_{X/S}\).
    \item \textbf{Use the stated hypothesis.}
    Since
    \(\lambda^1_{X/S}\) is induced precisely by restriction of normalized
    degree-one functions to this first structured infinitesimal neighbourhood, it follows that
    \[
    \lambda^1_{X/S}(\omega_2)=0.
    \]
    \item \textbf{Conclude the stated claim.}
    Thus \(\omega_2\) is a nonzero element of the displayed kernel.
\end{enumerate}
\end{proof}


The cotangent comparison therefore has a precise scope. First-order
restriction identifies the evaluation fibre with the realized cotangent
complex and sends the intrinsic de Rham differential to the universal
derivation. It does not collapse the finite-infinitesimal calculus to first
order: the quadratic class above is detected by a second-order thickening and
lies in the kernel of \(\lambda^1_{X/S}\). The remaining comparison is of a
different kind altogether. We now ask when finite-infinitesimal invariance,
augmented by convergence, agrees with reduced-classical invariance.
\section{Reduced-classical convergence of de Rham reflection}
\label{sec:reduced-classical-shadow}

The first comparison restricted the de Rham relation to first order. The
present comparison instead changes the localization. Finite-infinitesimal
invariance already identifies every finite Postnikov or nilpotent stage.
Reduced-classical reflection adds the convergence requirements for passage
through an infinite Postnikov tower and an arbitrary nilradical.

This separates three locality conditions that are easy to conflate:
\[
\begin{tikzcd}[column sep=large,row sep=large]
  & \text{de Rham-local}
      \arrow[dl,phantom,"\substack{\text{finite square-zero}\\\text{invariance}}" description]
      \arrow[dr,"\substack{\text{Postnikov continuity}\\+\ \text{nilradical continuity}}"]
  &
  \\
  \text{finite infinitesimal geometry}
  &&
  \text{reduced-classical local}
  \\
  & \text{reconstruction-local}
      \arrow[ul,phantom,"\substack{\text{Postnikov continuity}\\+\ \text{nilpotent excision}}" description]
  &
\end{tikzcd}
\]
The lower and upper-right conditions encode different requirements:
reconstruction uses nilpotent excision, whereas reduced-classical locality uses
nilradical continuity after finite-infinitesimal invariance has already been
imposed.
The goal of this section is to construct the comparison
\(\operatorname{dR}\to\operatorname{Red}_{\mathrm{cl}}\) and characterize
exactly when it is invertible.
For a connective commutative ring
spectrum \(A\), let \(r_{\mathrm{red}}\) send
\[
  \Spec A
  \longmapsto
  \Spec H\bigl((\pi_0A)_{\mathrm{red}}\bigr),
\]
and define on affine tests
\[
  (\operatorname{Red}_{\mathrm{cl}}F)(\Spec A)
  :=
  F\bigl(\Spec H((\pi_0A)_{\mathrm{red}})\bigr).
\]

\begin{lemma}[Reduced Zariski base change]
\label{lem:reduced-zariski-base-change}
Let \(A\to B\) represent an affine spectral Zariski open immersion. Then there
is a canonical equivalence
\[
  B\otimes_AH((\pi_0A)_{\mathrm{red}})
  \xrightarrow{\simeq}
  H((\pi_0B)_{\mathrm{red}}).
\]
In particular, for a principal spectral localization \(A\to A[f^{-1}]\),
\[
  A[f^{-1}]
  \otimes_A
  H((\pi_0A)_{\mathrm{red}})
  \xrightarrow{\simeq}
  H\bigl((\pi_0A[f^{-1}])_{\mathrm{red}}\bigr).
\]
Reduced-classical reflection therefore carries every finite affine spectral
Zariski \v{C}ech nerve to the \v{C}ech nerve of the corresponding reduced
classical cover.
\end{lemma}

\begin{proof}
\leavevmode
\begin{enumerate}[label=\textup{(\arabic*)}]
    \item \textbf{Introduce the notation.}
    Put
    \[
    R:=\pi_0A,
    \qquad
    S:=\pi_0B.
    \]
    \item \textbf{Construct the indicated map.}
    The maps
    \[
    A\longrightarrow B,
    \qquad
    H(R_{\mathrm{red}})\longrightarrow H(S_{\mathrm{red}})
    \]
    induce a canonical \(B\)-algebra morphism
    \[
    \theta:
    B\otimes_AH(R_{\mathrm{red}})
    \longrightarrow
    H(S_{\mathrm{red}}).
    \]
    \item \textbf{Establish the indicated equivalence.}
    We prove that \(\theta\) is an equivalence after a finite affine Zariski
    covering family of \(\Spec B\).
    \item \textbf{Compute the principal-open case.}
    For \(f\in R\), spectral localization is flat. The homotopy base-change
    formula for flat modules from
    \Ref{thm:spectral-lazard} gives
    \[
    \begin{aligned}
    A[f^{-1}]\otimes_AH(R_{\mathrm{red}})
    &\simeq
    H(R_{\mathrm{red}}[\bar f^{-1}])\\
    &\simeq
    H\bigl((R[\bar f^{-1}])_{\mathrm{red}}\bigr)\\
    &\simeq
    H\bigl((\pi_0A[f^{-1}])_{\mathrm{red}}\bigr).
    \end{aligned}
    \]
    Here the middle equivalence is the elementary identity that ordinary
    reduction commutes with localization.
    \item \textbf{Invoke the cited result.}
    By \Ref{prop:principal-opens-basis} and
    \Ref{prop:finite-principal-refinement}, there are finitely many
    \(f_i\in R\) such that
    \[
    D_A^{\mathrm{sp}}(f_i)\longrightarrow\Spec A
    \]
    factors through \(\Spec B\), and the resulting family covers \(\Spec B\).
    \item \textbf{Use the stated hypothesis.}
    Since \(\Spec B\to\Spec A\) is a monomorphism, the factorization implies
    \[
    B\otimes_AA[f_i^{-1}]
    \xrightarrow{\simeq}
    A[f_i^{-1}].
    \]
    \item \textbf{Conclude the stated claim.}
    Thus the morphisms \(B\to A[f_i^{-1}]\) form a finite affine spectral
    Zariski cover of \(\Spec B\).
    \item \textbf{Establish the next claim.}
    Base-change \(\theta\) along \(B\to A[f_i^{-1}]\).
    \item \textbf{Establish the next claim.}
    On the source,
    associativity and Step \textup{(2)} give
    \[
    \begin{aligned}
    \bigl(B\otimes_AH(R_{\mathrm{red}})\bigr)
    \otimes_BA[f_i^{-1}]
    &\simeq
    A[f_i^{-1}]\otimes_AH(R_{\mathrm{red}})\\
    &\simeq
    H\bigl((R[\bar f_i^{-1}])_{\mathrm{red}}\bigr).
    \end{aligned}
    \]
    \item \textbf{Establish the next claim.}
    On degree zero the equality
    \[
    S\otimes_RR[\bar f_i^{-1}]
    \simeq
    R[\bar f_i^{-1}]
    \]
    follows from
    \(B\otimes_AA[f_i^{-1}]\simeq A[f_i^{-1}]\).
    \item \textbf{Use the stated hypothesis.}
    Since
    \(B\to A[f_i^{-1}]\) is again a flat spectral localization, the flat homotopy
    base-change formula gives
    \[
    \begin{aligned}
    H(S_{\mathrm{red}})
    \otimes_BA[f_i^{-1}]
    &\simeq
    H\bigl(
    S_{\mathrm{red}}\otimes_SR[\bar f_i^{-1}]
    \bigr)\\
    &\simeq
    H\bigl((S[\bar f_i^{-1}])_{\mathrm{red}}\bigr)\\
    &\simeq
    H\bigl((R[\bar f_i^{-1}])_{\mathrm{red}}\bigr).
    \end{aligned}
    \]
    \item \textbf{Identify the stated objects.}
    Under these identifications the base change of \(\theta\) is the identity.
    \item \textbf{Establish the indicated equivalence.}
    The covering family is jointly conservative on \(B\)-modules, so
    \(\theta\) is an equivalence.
    \item \textbf{Fix the indicated data.}
    Let \(\{\Spec(B_i)\to\Spec(A)\}\) be a finite affine spectral Zariski
    covering family.
    \item \textbf{Establish the next claim.}
    Every iterated intersection
    \[
    B_{i_0}\otimes_A\cdots\otimes_AB_{i_n}
    \]
    is again an affine spectral Zariski open over \(A\).
    \item \textbf{Identify the stated objects.}
    Applying the first
    assertion degreewise identifies its reduction with base change to
    \(H(R_{\mathrm{red}})\).
    \item \textbf{Identify the stated objects.}
    Associativity of relative tensor products then
    identifies these degreewise reductions with the iterated fibre products of
    the reduced classical opens.
    \item \textbf{Establish the indicated equivalence.}
    The comparisons are natural in all face and
    degeneracy maps, so they assemble into the asserted equivalence of
    \v{C}ech nerves.
\end{enumerate}
\end{proof}

\begin{proposition}[Reduced-classical localization]
\label{prop:reduced-classical-localization}
The endofunctor \(\operatorname{Red}_{\mathrm{cl}}\) preserves spectral Zariski sheaves and defines
an accessible left-exact localization of \(\mathcal H_{\Sp}\). Its local
objects are precisely the \(F\) for which
\[
  F(\Spec A)
  \xrightarrow{\simeq}
  F\bigl(\Spec H((\pi_0A)_{\mathrm{red}})\bigr)
\]
for every connective \(A\).
\end{proposition}

\begin{proof}
\leavevmode
\begin{enumerate}[label=\textup{(\arabic*)}]
    \item \textbf{Fix the indicated data.}
    Let
    \[
    \iota_{\mathrm{Zar}}:
    \mathcal H_{\Sp}
    \hookrightarrow
    \PSh(\Aff_{\Sp})
    \]
    be the fully faithful inclusion and let \(r_{\mathrm{red}}^*\) denote precomposition with
    reduced-classical reflection on affine spectral schemes.
    \item \textbf{Establish the next claim.}
    At the presheaf level, \(r_{\mathrm{red}}^*\) is
    accessible and preserves all limits.
    \item \textbf{Apply the cited lemma.}
    Lemma~\ref{lem:reduced-zariski-base-change}
    shows directly that it carries every finite affine spectral Zariski \v{C}ech
    nerve to the corresponding reduced classical \v{C}ech diagram.
    \item \textbf{Use the stated hypothesis.}
    Since these
    finite affine covers generate the spectral Zariski topology of
    Chapter~\Ref{chap:spectral-algebraic-geometry} and \cite{LurieSAG},
    \(r_{\mathrm{red}}^*\iota_{\mathrm{Zar}}\) already lands in spectral Zariski
    sheaves.
    \item \textbf{Conclude the stated claim.}
    Thus the sheaf endofunctor is canonically the
    composite
    \[
    \operatorname{Red}_{\mathrm{cl}}
    \simeq
    a^{\mathrm{sp}}_{\mathrm{Zar}}
    r_{\mathrm{red}}^*\iota_{\mathrm{Zar}},
    \]
    and hence is accessible.
    \item \textbf{Use the stated identification.}
    Because the displayed presheaf precomposition preserves
    limits objectwise and the sheafification is left exact, \(\operatorname{Red}_{\mathrm{cl}}\)
    preserves finite limits.
    \item \textbf{Construct the indicated map.}
    The canonical map
    \[
    A\longrightarrow H((\pi_0A)_{\mathrm{red}})
    \]
    induces a coaugmentation \(\eta^{\mathrm{red}}:\id\to \operatorname{Red}_{\mathrm{cl}}\).
    \item \textbf{Establish the next claim.}
    The reduced-classical
    affine endofunctor is canonically idempotent, so both transformations
    \[
    \operatorname{Red}_{\mathrm{cl}}\eta^{\mathrm{red}},
    \qquad
    \eta^{\mathrm{red}}\operatorname{Red}_{\mathrm{cl}}:
    \operatorname{Red}_{\mathrm{cl}}\longrightarrow \operatorname{Red}_{\mathrm{cl}}^2
    \]
    are equivalences.
    \item \textbf{Conclude the stated claim.}
    Therefore the coaugmented endofunctor is an accessible
    idempotent localization.
    \item \textbf{Establish the indicated equivalence.}
    Its finite-limit preservation makes it left exact, and
    its fixed objects are exactly those satisfying the displayed affine-test
    equivalences.
\end{enumerate}
\end{proof}

\begin{lemma}[Finite nilpotent quotients are de Rham-equivalent]
\label{lem:finite-nilpotent-towers-de-Rham}
Let \(R\) be an ordinary commutative ring and \(I\subseteq R\) nilpotent. Then
\[
  HR\longrightarrow H(R/I)
\]
is a finite composite of spectral square-zero extensions. Consequently
every de Rham-local object evaluates \(\Spec HR\) and \(\Spec H(R/I)\)
equivalently.
\end{lemma}

\begin{proof}
\leavevmode
\begin{enumerate}[label=\textup{(\arabic*)}]
    \item \textbf{Establish the next claim.}
    This is the coordinate form of
    \Ref{prop:discrete-nilpotent-finite-square-zero} from
    Chapter~\Ref{chap:differentiation-spectral-doctrine}.
    \item \textbf{Establish the next claim.}
    The proof of that proposition
    constructs the finite tower by powers of \(I\), using
    \Ref{lem:discrete-square-zero-recognition} there to recognize each quotient stage
    as a spectral square-zero extension.
    \item \textbf{Apply the cited theorem.}
    Theorem~\ref{thm:de-Rham-locality} then gives the displayed invariance.
\end{enumerate}
\end{proof}

\subsection{From de Rham reflection to reduced-classical reflection}

Reduced-classical reflection forgets both higher homotopy and nilpotent
classical structure. The finite part of that forgetting is already invisible
to de Rham-local objects. The following comparison therefore exists by the
universal property of de Rham reflection.

\begin{theorem}[Comparison with reduced-classical reflection]
\label{thm:de-Rham-to-reduced}
Every de Rham generator is inverted by \(\operatorname{Red}_{\mathrm{cl}}\). Hence the universal
property of \(L^{\mathrm{dR}}\) supplies a canonical transformation of left-exact
ambient endofunctors
\[
  \kappa:\operatorname{dR}\longrightarrow\operatorname{Red}_{\mathrm{cl}}.
\]
Thus the two localizations fit into
\[
\begin{tikzcd}[column sep=huge]
  \operatorname{dR}
    \arrow[r,"\kappa"]
  &
  \operatorname{Red}_{\mathrm{cl}}.
\end{tikzcd}
\]

For \(X\to S\), it induces
\[
  \kappa_{X/S}:
  X_{\mathrm{dR}/S}
  \longrightarrow
  S\times_{\operatorname{Red}_{\mathrm{cl}}(S)}\operatorname{Red}_{\mathrm{cl}}(X),
\]
compatible with arbitrary base change.

The absolute comparison is compatible with the two localization units:
\[
\begin{tikzcd}[column sep=huge,row sep=large]
  X
    \arrow[r,"\eta_X^{\mathrm{dR}}"]
    \arrow[dr]
  &
  X_{\mathrm{dR}}
    \arrow[d,"\kappa_X"]
  \\
  &
  \operatorname{Red}_{\mathrm{cl}}(X).
\end{tikzcd}
\]

\end{theorem}

\begin{proof}
\leavevmode
\begin{enumerate}[label=\textup{(\arabic*)}]
    \item \textbf{Fix the indicated data.}
    Let \(A'\to A\) be a spectral square-zero extension, so the associated
    generator is
    \[
    h_A\longrightarrow h_{A'}.
    \]
    \item \textbf{Establish the next claim.}
    For a connective test algebra \(B\), since the target
    \(H((\pi_0B)_{\mathrm{red}})\) is discrete, the \(0\)-truncated adjunction of
    Chapter~\Ref{chap:spectral-algebraic-geometry} gives
    \[
    \begin{aligned}
    (\operatorname{Red}_{\mathrm{cl}}h_A)(\Spec B)
    &=
    \Map_{\CAlg(\Sp)}\bigl(A,H((\pi_0B)_{\mathrm{red}})\bigr)\\
    &\simeq
    \Hom_{\mathrm{CRing}}
    \bigl(\pi_0A,(\pi_0B)_{\mathrm{red}}\bigr)\\
    &\simeq
    \Hom_{\mathrm{CRing}}
    \bigl((\pi_0A)_{\mathrm{red}},(\pi_0B)_{\mathrm{red}}\bigr).
    \end{aligned}
    \]
    \item \textbf{Establish the next claim.}
    The same formula holds with \(A'\) in place of \(A\).
    \item \textbf{Establish the next claim.}
    A square-zero extension
    induces a square-zero quotient on \(\pi_0\), hence
    \[
    (\pi_0A')_{\mathrm{red}}
    \simeq
    (\pi_0A)_{\mathrm{red}}.
    \]
    \item \textbf{Conclude the stated claim.}
    Therefore \(\operatorname{Red}_{\mathrm{cl}}\) sends the generator to a pointwise equivalence.
    \item \textbf{Establish the next claim.}
    Initiality of Theorem~\ref{thm:intrinsic-de-Rham-localization} now constructs
    \(\kappa\).
    \item \textbf{Use pullback stability.}
    The relative comparison is obtained by applying the two left-exact
    constructions to the defining pullback of the relative de Rham object; its
    base-change coherence is inherited from finite-limit functoriality.
\end{enumerate}
\end{proof}

\subsection{Finite collapse and the convergence gap}

The comparison above is automatically an equivalence on tests for which the
passage to the reduced classical ring is finite. This isolates the only
remaining issue: convergence when the Postnikov tower or nilradical cannot be
removed in finitely many square-zero stages.

\begin{proposition}[Finite Postnikov and nilpotent collapse]
\label{prop:finite-collapse}
Let \(A\) be connective and eventually coconnective, and suppose the nilradical of
\(\pi_0A\) is nilpotent. Then the coordinate morphism
\[
  A\longrightarrow H((\pi_0A)_{\mathrm{red}})
\]
is a finite composite of spectral square-zero extensions. Equivalently,
\[
  \Spec H((\pi_0A)_{\mathrm{red}})
  \longrightarrow
  \Spec A
\]
is a finite spectral infinitesimal substitution. Consequently every de Rham-local
\(F\) satisfies
\[
  F(\Spec A)
  \simeq
  F\bigl(\Spec H((\pi_0A)_{\mathrm{red}})\bigr).
\]
\end{proposition}

\begin{proof}
\leavevmode
\begin{enumerate}[label=\textup{(\arabic*)}]
    \item \textbf{Invoke the cited result.}
    By \Ref{thm:postnikov-square-zero-stages} of
    Chapter~\Ref{chap:formal-reconstruction-spectral-doctrine}, the multiplicative
    Postnikov tower has canonical structured square-zero transition maps.
    \item \textbf{Establish the next claim.}
    Eventual
    coconnectivity makes the tower from \(A\) to \(H\pi_0A\) finite.
    \item \textbf{Establish the next claim.}
    The nilpotent
    quotient from \(\pi_0A\) to its reduction is a finite composite of underlying
    spectral square-zero extensions by
    Lemma~\ref{lem:finite-nilpotent-towers-de-Rham}.
    \item \textbf{Establish the next claim.}
    Concatenating these two finite
    chains proves the first assertion.
    \item \textbf{Establish the next claim.}
    The canonical Postnikov stages and the finite nilpotent quotient tower together
    give the required finite composite.
    \item \textbf{Establish the locality claim.}
    The de Rham-locality conclusion follows from
    Theorem~\ref{thm:de-Rham-locality}.
\end{enumerate}
\end{proof}

\subsection{Postnikov and nilradical convergence}

We now formulate the two continuity conditions that close the finite-to-
infinite gap. Their combination is exactly what distinguishes
reduced-classical locality inside the de Rham-local sector.

\begin{definition}[Nilradical continuity on affine tests]
A de Rham-local object \(F\) is \emph{nilradical-continuous on affine tests} if,
for every ordinary ring \(R\) with nilradical \(N\),
\[
  \operatorname*{colim}_{\substack{J\subseteq N\\J\text{ finitely generated}}}
  F(\Spec H(R/J))
  \xrightarrow{\simeq}
  F(\Spec H(R/N)).
\]
The transition maps are those induced by inclusion of finitely generated ideals.
\end{definition}

\begin{theorem}[Convergence criterion]
\label{thm:reduced-convergence}
Let \(F\) be de Rham-local. The following are equivalent:
\begin{enumerate}
\item \(F\) is \(\operatorname{Red}_{\mathrm{cl}}\)-local;
\item \(F\) is local for the Postnikov-continuity generators of
\Ref{def:global-spectral-reconstruction-generators} and is
nilradical-continuous on affine tests.
\end{enumerate}
Thus reduced-classical reflection is the further left-exact localization of the de
Rham topos imposing Postnikov continuity together with nilradical continuity.
\end{theorem}

\begin{proof}
\leavevmode
\begin{enumerate}[label=\textup{(\arabic*)}]
    \item \textbf{Fix the indicated data.}
    Let \(R\) be an ordinary ring with nilradical \(N\), and let
    \(J=(x_1,\ldots,x_r)\subseteq N\) be finitely generated.
    \item \textbf{Establish the next claim.}
    Choose
    \(n_i\geq1\) with \(x_i^{n_i}=0\).
    \item \textbf{Establish the next claim.}
    Then
    \[
    J^{n_1+\cdots+n_r}=0,
    \]
    so \(J\) is nilpotent.
    Hence
    \[
    \sqrt{J}=N
    \qquad\text{and}\qquad
    (R/J)_{\mathrm{red}}\simeq R/N.
    \]
    \item \textbf{Treat the stated case.}
    If \(F\) is reduced-classical-local, then \(A\), every
    \(\tau_{\leq n}^{\CAlg}A\), and
    \(H((\pi_0A)_{\mathrm{red}})\) have the same reduced classical test.
    \item \textbf{Conclude the stated claim.}
    Hence
    \(F\) is local for the Postnikov-continuity generators.
    \item \textbf{Invoke the cited result.}
    By the preceding
    calculation, every finitely generated \(J\subseteq N\) has
    \[
    (R/J)_{\mathrm{red}}\simeq R/N,
    \]
    so the nilradical-continuity diagram is canonically constant up to equivalence.
    \item \textbf{Prove the converse implication.}
    Conversely, by \Ref{cor:infinitesimal-postnikov-reconstruction}, every finite map
    \[
    \tau_{\leq n}^{\CAlg}A\longrightarrow H\pi_0A
    \]
    is a finite infinitesimal substitution.
    \item \textbf{Use the stated construction.}
    De Rham locality therefore gives
    \[
    F(\Spec\tau_{\leq n}^{\CAlg}A)
    \simeq
    F(\Spec H\pi_0A)
    \]
    for every \(n\).
    \item \textbf{Establish the locality claim.}
    Locality for the Postnikov-continuity generators yields
    \[
    F(\Spec A)
    \simeq
    \varprojlim\nolimits_nF(\Spec\tau_{\leq n}^{\CAlg}A)
    \simeq
    F(\Spec H\pi_0A).
    \]
    \item \textbf{Establish the next claim.}
    For every finitely generated \(J\subseteq N\), the first step shows that \(J\)
    is nilpotent.
    \item \textbf{Apply the cited lemma.}
    Lemma~\ref{lem:finite-nilpotent-towers-de-Rham} therefore gives
    \[
    F(\Spec HR)
    \simeq
    F(\Spec H(R/J)).
    \]
    \item \textbf{Establish the indicated equivalence.}
    The quotient maps \(R\to R/J\) assemble these equivalences into a natural
    objectwise equivalence from the constant diagram at \(F(\Spec HR)\) to
    \[
    J\longmapsto F(\Spec H(R/J)).
    \]
    \item \textbf{Establish the next claim.}
    The indexing poset of finitely generated ideals \(J\subseteq N\) is filtered and
    has the initial object \(0\), so its nerve is contractible.
    \item \textbf{Establish the next claim.}
    Its colimit is
    therefore canonically equivalent to \(F(\Spec HR)\).
    \item \textbf{Use the stated construction.}
    Nilradical continuity gives
    \[
    F(\Spec HR)
    \simeq
    F(\Spec H(R/N)).
    \]
    \item \textbf{Take the indicated specialization.}
    Taking \(R=\pi_0A\) and combining with the Postnikov comparison proves
    \[
    F(\Spec A)
    \simeq
    F\bigl(\Spec H((\pi_0A)_{\mathrm{red}})\bigr),
    \]
    which is reduced-classical locality.
    \item \textbf{Invoke the cited result.}
    By Theorem~\ref{thm:de-Rham-to-reduced}, every
    \(\operatorname{Red}_{\mathrm{cl}}\)-local object is de Rham-local and
    \[
    \operatorname{Red}_{\mathrm{cl}}\circ\operatorname{dR}
    \simeq
    \operatorname{Red}_{\mathrm{cl}}.
    \]
    \item \textbf{Conclude the stated claim.}
    Therefore restriction of \(\operatorname{Red}_{\mathrm{cl}}\) to the de
    Rham-local full subcategory is itself an accessible left-exact reflector, and
    the equivalence just proved identifies its fixed objects precisely by the two
    stated convergence properties.
\end{enumerate}
\end{proof}

\begin{remark}[Comparison of the three left-exact reflections]
\label{rem:de-Rham-reduced-reconstruction-distinction}
The de Rham reflector, the reduced-classical reflector, and the reconstruction
reflector of Chapter~\Ref{chap:formal-reconstruction-spectral-doctrine} encode
three locality conditions. The functor \(\operatorname{dR}\) inverts finite
square-zero infinitesimal variation. On the de Rham-local subcategory,
\(\operatorname{Red}_{\mathrm{cl}}\) imposes the Postnikov and nilradical
convergence conditions of Theorem~\ref{thm:reduced-convergence}. The reflector
\(\operatorname{Recon}_S\) is generated by Postnikov continuity together with
nilpotent excision. These descriptions provide the comparison points
used in the preceding results.
\end{remark}

\begin{corollary}[Invertibility criterion for the comparison]
\label{cor:invertibility-reduced-comparison}
For every \(G\in\mathcal H_{\Sp}\),
\[
  \kappa_G:\operatorname{dR}(G)\longrightarrow \operatorname{Red}_{\mathrm{cl}}(G)
\]
is an equivalence if and only if \(\operatorname{dR}(G)\) is local for the
Postnikov-continuity generators of
\Ref{def:global-spectral-reconstruction-generators} and is
nilradical-continuous on affine tests.
\end{corollary}

\begin{proof}
\leavevmode
\begin{enumerate}[label=\textup{(\arabic*)}]
    \item \textbf{Establish the locality claim.}
    The functor \(\operatorname{Red}_{\mathrm{cl}}\) is a further localization of the de Rham-local
    subcategory.
    \item \textbf{Apply the indicated result.}
    Apply Theorem~\ref{thm:reduced-convergence} to
    \(\operatorname{dR}(G)\).
\end{enumerate}
\end{proof}


We can now state the comparison without collapsing distinct notions.
De Rham reflection enforces invariance under finite infinitesimal variation.
Reduced-classical reflection is the further localization obtained by adding
Postnikov and nilradical continuity. Reconstruction, by contrast, combines
Postnikov continuity with nilpotent excision. These three descriptions give
the precise interfaces among the localizations and complete the comparison
part of the chapter.
\section{The spectral de Rham interface}
\label{sec:de-Rham-construction}

All constructions named in the chapter are now active, so the global
dependency diagrams can be read without introducing future notation. The
first records how the effective relation feeds the stable descent calculus:
\[
\begin{tikzcd}[column sep=large,row sep=large]
  X
    \arrow[r,"\eta_{X/S}^{\mathrm{dR}}"]
    \arrow[d,two heads,"e_{X/S}"']
  &
  X_{\mathrm{dR}/S}
  &
  \mathcal R^{\mathrm{dR}}_{X/S,\bullet}
    \arrow[l,phantom,"\leadsto" description]
    \arrow[r,"\text{stable scalar functions}"]
  &
  C^\bullet_{\mathrm{dR}}(X/S)
    \arrow[d,"\operatorname*{Tot}"]
  \\
  Q_{X/S}
    \arrow[ur,hook,"m_{X/S}"']
  &&
  \bigl(\Omega^\bullet_{\mathrm{dR}}(X/S),d_{\mathrm{dR}}\bigr)
    \arrow[r,phantom,"\longleftarrow" description]
  &
  \mathrm{DR}_{X/S}.
\end{tikzcd}
\]

The second records the two comparison exits from that calculus: the
reduced-classical comparison from the full reflection and the first-order
cotangent comparison from intrinsic one-forms.
\[
\begin{tikzcd}[column sep=large,row sep=large]
  X \arrow[r,"\eta_{X/S}^{\mathrm{dR}}"] \arrow[d,two heads,"e_{X/S}"']
  & X_{\mathrm{dR}/S}
    \arrow[r,"\kappa_{X/S}"]
  & S\times_{\operatorname{Red}_{\mathrm{cl}}(S)}\operatorname{Red}_{\mathrm{cl}}(X)
  \\
  Q_{X/S}
    \arrow[r,phantom,"\leadsto" description]
  & \mathcal R^{\mathrm{dR}}_{X/S,\bullet}
    \arrow[r]
  & C^\bullet_{\mathrm{dR}}(X/S)
    \arrow[d,"\operatorname*{Tot}"]
  \\
  && \mathrm{DR}_{X/S}
    \arrow[d,phantom,"\leadsto" description]
  \\
  && \Omega^1_{\mathrm{dR}}(X/S)
    \arrow[d,"\lambda^1_{X/S}"]
  \\
  && \Pi_p\operatorname{Real}^{\mathrm{st}}_X(L_{X/S}).
\end{tikzcd}
\]

\begin{theorem}[Structural summary of the spectral de Rham construction]
\label{thm:doctrinal-recognition-spectral-de-Rham}
The constructions of this chapter assemble into the following sequence.
\begin{enumerate}[label=(\roman*)]
  \item The underlying square-zero substitutions between affine spectral schemes generate a pullback-stable
  infinitesimal congruence \(W^{\mathrm{dR}}\) on spectral contexts.

  \item The universal accessible left-exact localization at that congruence is the
  indexed de Rham reflection \(\operatorname{dR}_S\). Its local objects are
  characterized by finite-infinitesimal invariance, and its right
  orthogonal class consists of the formally \(\acute{e}\)tale morphisms.

  \item Effective image factorization of the relative de Rham unit produces the
  effective quotient \(Q_{X/S}\), whose \v{C}ech nerve
  \[
    \mathcal R^{\mathrm{dR}}_{X/S,\bullet}
  \]
  is the de Rham relation groupoid generated on \(X\).

  \item The stabilized slice categories of
  Chapter~\Ref{chap:formal-reconstruction-spectral-doctrine}, with the dependent-adjoint
  structure established here, give the parameterized stable slice doctrine
  \[
    S\longmapsto\mathcal D_{\Sp/S},
  \]
  and stable realization interprets quasi-coherent coefficients inside that
  doctrine. The realized structure sheaf
  \[
    \mathcal O_S^{\mathrm{add}}
  \]
  is its distinguished stable additive scalar.

  \item The cosimplicial object
  \[
    C^\bullet_{\mathrm{dR}}(X/S)
  \]
  consists of stable scalar functions on de Rham relation simplices.
  Comonadic effective descent identifies its totalization with invariant scalar
  functions on the effective quotient.

  \item The complete totalization filtration is the totalization filtration of that
  descent problem. Its normalized layers are the intrinsic de Rham forms, its
  connecting maps form the homotopy-coherent de Rham differential, and its higher
  stages classify higher closedness, exact-form lifts, and de Rham primitive
  torsors.

  \item Affine descent of the stable-section form-space functors constructs the internal
  comprehension objects
  \[
    \mathbf{Form}^q_{\mathrm{dR},S},
    \qquad
    \mathbf{Form}^{q,\mathrm{closed}}_{\mathrm{dR},S},
  \]
  together with the stable-section-to-comprehended transformations
  \[
    \mathrm{cmp}^q,
    \qquad
    \mathrm{cmp}^{q,\mathrm{closed}}.
  \]
  These transformations are equivalences on affine contexts, and the displayed
  objects comprehend the corresponding affine-descended form judgments on
  arbitrary contexts.

  \item For a morphism \(p:X\to S\) of spectral schemes, restriction to the
  canonical first structured infinitesimal neighbourhood gives the first-order
  coefficient
  \[
    K^{(1)}_{X/S}
    =
    \operatorname{fib}
    \bigl(
      \mathscr F^{(1)}_{X/S}
      \longrightarrow
      \mathcal O_X^{\mathrm{add}}
    \bigr)
  \]
  and a canonical equivalence
  \[
    K^{(1)}_{X/S}
    \xrightarrow{\simeq}
    \operatorname{Real}^{\mathrm{st}}_X(L_{X/S}).
  \]
  The resulting first-order linearization
  \[
    \lambda^1_{X/S}:
    \Omega^1_{\mathrm{dR}}(X/S)
    \longrightarrow
    \Pi_p\operatorname{Real}^{\mathrm{st}}_X(L_{X/S})
  \]
  satisfies
  \[
    \lambda^1_{X/S}d^0_{\mathrm{dR}}
    \simeq
    d^{\mathrm{univ,st}}_{X/S}.
  \]
\item Reduced-classical reflection is a further left-exact localization of the
  de Rham-local sector characterized by Postnikov continuity and nilradical
  continuity. The comparison
  \(\kappa:\operatorname{dR}\to\operatorname{Red}_{\mathrm{cl}}\) is invertible
  exactly on the corresponding convergence domain. This localization is
  distinct from the first-order cotangent comparison of the preceding item.
\end{enumerate}
\end{theorem}

\begin{proof}
\leavevmode
\begin{enumerate}[label=\textup{(\arabic*)}]
    \item \textbf{Collect localization, factorization, and the effective relation.}
    Items \textup{(i)}--\textup{(iii)} follow from
    Theorems~\ref{thm:intrinsic-de-Rham-localization},
    \ref{thm:indexed-spectral-de-Rham-reflection},
    \ref{thm:de-Rham-etale-factorization-system}, and
    \ref{thm:relation-effective-presentation}, together with
    Theorem~\ref{thm:crystalline-recognition}.
    \item \textbf{Establish the next claim.}
    Item \textup{(iv)} is
    Corollary~\ref{cor:stable-linear-hyperdoctrine-recognition} together with
    Construction~\ref{con:stable-realization-qcoh} and
    Proposition~\ref{prop:stable-realization-qcoh}.
    \item \textbf{Establish the next claim.}
    Item \textup{(v)} is
    Theorems~\ref{thm:internal-function-cochains-expanded},
    \ref{thm:comonadic-effective-descent-de-Rham}, and
    \ref{thm:effective-DR-functions}.
    \item \textbf{Establish the next claim.}
    Item \textup{(vi)} follows from
    Theorems~\ref{thm:complete-filtration} and
    \ref{thm:complete-filtrations-coherent-cochains}, together with the subsequent
    closed-form and de Rham primitive constructions.
    \item \textbf{Establish the next claim.}
    Item \textup{(vii)} follows
    from Construction~\ref{con:affine-descent-form-comprehension},
    Proposition~\ref{prop:form-comprehension-right-kan},
    Construction~\ref{con:intrinsic-comprehended-form-comparison}, and
    Theorem~\ref{thm:internal-moduli-forms-detailed}.
    \item \textbf{Establish the next claim.}
    Item \textup{(viii)} follows from
    Proposition~\ref{prop:stable-cotangent-realization},
    Construction~\ref{con:stabilized-universal-derivation}, and
    Theorem~\ref{thm:first-order-linearization}.
    \item \textbf{Establish the next claim.}
    Item \textup{(ix)} follows from
    Theorems~\ref{thm:de-Rham-to-reduced} and
    \ref{thm:reduced-convergence}.
\end{enumerate}
\end{proof}



The theorem above is the citation interface of the chapter. Later arguments
can use the de Rham reflector, effective relation, stable scalar cochains,
complete form filtration, comprehension objects, and the two comparison maps
without reopening their constructions.

\section{Chapter summary}
\label{sec:de-Rham-summary}

The chapter now supplies a de Rham descent calculus for finite infinitesimal
variation. The indexed left-exact reflector identifies precisely the contexts
that are indistinguishable under finite infinitesimal substitution, and its
right orthogonal class is the formally \(\acute{e}\)tale class. This is the
localization interface used by every later construction in the chapter.

The relative unit contains two different pieces of information. Its target
\(X_{\mathrm{dR}/S}\) is the universal de Rham-local reflection, whereas its
effective image \(Q_{X/S}\) is the object actually presented by the
generalized points of \(X\). The de Rham relation is the effective
equivalence-relation groupoid of \(X\to Q_{X/S}\). On the formally smooth
effectivity domain the distinction disappears, but the construction does not
require that special case.

Stable scalar functions on this relation form a cosimplicial commutative
algebra. Effective descent identifies its totalization with scalar functions
on \(Q_{X/S}\). The complete totalization filtration then exposes the
successive descent layers: normalization gives the intrinsic de Rham forms,
the connecting morphisms give the homotopy-coherent differential, and the
remaining filtration data record multiplication, higher closedness,
exact-form lifts, and primitive torsors. Formal \(\acute{e}\)tale descent
and affine density turn the resulting affine stable-section semantics into
internal comprehension objects on arbitrary contexts.

Two distinct comparison interfaces delimit what the resulting calculus
remembers. Restriction to the first structured infinitesimal neighbourhood
recovers the realized cotangent complex and the universal derivation, while the
quadratic example shows that higher finite-infinitesimal information can survive
outside that first-order image. Reduced-classical reflection tests a different
boundary: it adds Postnikov and nilradical convergence to the de Rham-local
sector. Later arguments may therefore use the cotangent comparison to extract
first-order information and the reduced-classical comparison to test convergence,
without conflating the two operations.
